%% ---- begin expanded/00_preamble.tex ----
%% =====================================================================
%%  MTP Report -- Khem Singh (25cs06004), IIT Bhubaneswar
%%  Fully Dynamic Algorithms for Chordal Graphs and Split Graphs:
%%  implementation of L. Ibarra, ACM TALG 4(4), Article 40, 2008.
%%  Source code: https://github.com/khemsinghbhati/MTP-3rd-sem
%%
%%  Compile with:  pdflatex MTP_Report.tex   (run three times)
%%  Only standard TeX Live packages are used; no external image files.
%% =====================================================================
\documentclass[11pt,a4paper,oneside]{report}

\usepackage[T1]{fontenc}
\usepackage[utf8]{inputenc}
\usepackage[sc]{mathpazo}
\linespread{1.07}
\renewcommand{\ttdefault}{lmtt}
\renewcommand{\sfdefault}{lmss}
\usepackage[a4paper,top=25mm,bottom=25mm,left=27mm,right=25mm,headheight=14pt]{geometry}
\usepackage{microtype}
\usepackage{amsmath,amssymb,amsthm,mathtools}
\usepackage{graphicx}
\usepackage[dvipsnames,table]{xcolor}
\usepackage{tikz}
\usetikzlibrary{positioning,arrows.meta,shapes.geometric,shapes.misc,fit,calc,
  decorations.pathreplacing,backgrounds,matrix}
\usepackage{pgfplots}
\pgfplotsset{compat=1.18}
\usepackage{booktabs,tabularx,longtable,multirow,makecell,array}
\usepackage{enumitem}
\usepackage{listings}
\usepackage{algorithm}
\usepackage{algpseudocode}
\usepackage[most]{tcolorbox}
\usepackage{caption}
\usepackage{subcaption}
\usepackage{float}
\usepackage{fancyhdr}
\usepackage{titlesec}
\usepackage[titles]{tocloft}
\usepackage{setspace}
\usepackage{url}
\usepackage[colorlinks=true,linkcolor=accent,citecolor=accent!70!black,
  urlcolor=accent!80!black,bookmarksnumbered=true,
  pdftitle={MTP Report: Fully Dynamic Algorithms for Chordal Graphs and Split Graphs},
  pdfauthor={Khem Singh (25cs06004)},
  pdfsubject={Implementation of Ibarra (2008), ACM Transactions on Algorithms}]{hyperref}
\usepackage[nameinlink,capitalise,noabbrev]{cleveref}

%% ---------------------------------------------------------------------
%% Colours
%% ---------------------------------------------------------------------
\definecolor{accent}{RGB}{0,72,128}
\definecolor{accentlight}{RGB}{226,236,246}
\definecolor{codeback}{RGB}{248,248,246}
\definecolor{okgreen}{RGB}{0,115,60}
\definecolor{warnorange}{RGB}{196,96,0}
\definecolor{badred}{RGB}{170,20,30}

%% ---------------------------------------------------------------------
%% Page style and headings
%% ---------------------------------------------------------------------
\pagestyle{fancy}
\fancyhf{}
\fancyhead[L]{\small\color{black!60}\nouppercase{\leftmark}}
\fancyhead[R]{\small\color{black!60}MTP Report \textbar\ Khem Singh}
\fancyfoot[C]{\small\thepage}
\renewcommand{\headrulewidth}{0.3pt}
\fancypagestyle{plain}{\fancyhf{}\fancyfoot[C]{\small\thepage}\renewcommand{\headrulewidth}{0pt}}

\titleformat{\chapter}[display]
  {\normalfont\bfseries\color{accent}}
  {\normalfont\large\scshape\color{accent}\chaptertitlename\ \thechapter}
  {4pt}
  {\titlerule[0.8pt]\vspace{6pt}\Huge}
\titlespacing*{\chapter}{0pt}{-12pt}{22pt}
\titleformat{\section}{\normalfont\Large\bfseries\color{accent}}{\thesection}{0.8em}{}
\titleformat{\subsection}{\normalfont\large\bfseries\color{accent!85!black}}{\thesubsection}{0.8em}{}
\titleformat{\subsubsection}{\normalfont\normalsize\bfseries}{\thesubsubsection}{0.8em}{}
\setcounter{secnumdepth}{3}
\renewcommand{\bibname}{References}
\setcounter{tocdepth}{2}

\setlength{\parskip}{3pt plus 1pt}
\setlist{itemsep=2pt,topsep=4pt,parsep=1pt}
\emergencystretch=2.5em
\setlength{\cftfignumwidth}{2.8em}
\setlength{\cfttabnumwidth}{2.8em}

\captionsetup{font=small,labelfont={bf,color=accent},format=hang,margin=6pt}

%% ---------------------------------------------------------------------
%% Theorem-like environments
%% ---------------------------------------------------------------------
\theoremstyle{plain}
\newtheorem{theorem}{Theorem}[chapter]
\newtheorem{lemma}[theorem]{Lemma}
\newtheorem{proposition}[theorem]{Proposition}
\newtheorem{corollary}[theorem]{Corollary}
\newtheorem{claim}[theorem]{Claim}
\theoremstyle{definition}
\newtheorem{definition}[theorem]{Definition}
\newtheorem{example}[theorem]{Example}
\theoremstyle{remark}
\newtheorem{remark}[theorem]{Remark}
\crefname{claim}{Claim}{Claims}

%% ---------------------------------------------------------------------
%% Boxes
%% ---------------------------------------------------------------------
\tcbset{boxsep=2pt,left=6pt,right=6pt,top=4pt,bottom=4pt,fonttitle=\bfseries\small,
  breakable,enhanced}
\newtcolorbox{keyidea}[1][Key idea]{colback=accentlight!55,colframe=accent!80,title={#1},
  coltitle=white}
\newtcolorbox{incode}[1][Where this lives in the code]{colback=teal!4,colframe=teal!60!black,
  title={#1},coltitle=white}
\newtcolorbox{paperbox}[1][What the paper says]{colback=gray!5,colframe=gray!55!black,
  title={#1},coltitle=white}
\newtcolorbox{caution}[1][Paper versus implementation]{colback=orange!6,colframe=warnorange,
  title={#1},coltitle=white}
\newtcolorbox{summarybox}[1][Summary]{colback=okgreen!5,colframe=okgreen!80!black,
  title={#1},coltitle=white}

%% ---------------------------------------------------------------------
%% Code listings
%% ---------------------------------------------------------------------
\lstdefinestyle{cpp}{language=C++,
  basicstyle=\ttfamily\footnotesize,
  keywordstyle=\color{accent}\bfseries,
  commentstyle=\color{green!40!black}\itshape,
  stringstyle=\color{red!55!black},
  numbers=left,numberstyle=\tiny\color{gray},numbersep=7pt,
  frame=tb,framerule=0.5pt,rulecolor=\color{gray!60},
  backgroundcolor=\color{codeback},
  breaklines=true,breakatwhitespace=false,columns=fullflexible,keepspaces=true,
  showstringspaces=false,xleftmargin=2.2em,framexleftmargin=2.2em,
  aboveskip=8pt,belowskip=6pt,tabsize=4,
  morekeywords={nullptr,constexpr,noexcept,override,size_t},
  escapeinside={(*@}{@*)}}
\lstdefinestyle{shell}{language={},numbers=none,keywordstyle={},commentstyle={},
  stringstyle={},basicstyle=\ttfamily\footnotesize,
  frame=single,framerule=0.4pt,rulecolor=\color{gray!50},
  backgroundcolor=\color{codeback},breaklines=true,columns=fullflexible,
  keepspaces=true,showstringspaces=false,xleftmargin=0.6em,framexleftmargin=0.6em,
  aboveskip=6pt,belowskip=6pt}
\lstset{style=cpp}

%% ---------------------------------------------------------------------
%% Algorithms
%% ---------------------------------------------------------------------
\algrenewcommand\algorithmicrequire{\textbf{Input:}}
\algrenewcommand\algorithmicensure{\textbf{Output:}}
\algnewcommand{\LineComment}[1]{\State\textcolor{black!55}{\(\triangleright\) \emph{#1}}}
\renewcommand{\algorithmiccomment}[1]{\hfill\textcolor{black!55}{\(\triangleright\) #1}}
\crefname{algorithm}{Algorithm}{Algorithms}
\makeatletter
\providecommand{\theHALG@line}{}\renewcommand{\theHALG@line}{\thealgorithm.\arabic{ALG@line}}
\makeatother

%% ---------------------------------------------------------------------
%% TikZ styles
%% ---------------------------------------------------------------------
\tikzset{
  vtx/.style={circle,draw=accent!85!black,fill=accentlight,thick,minimum size=6.6mm,
    inner sep=0pt,font=\small},
  vtxh/.style={vtx,fill=orange!25,draw=orange!80!black},
  vtxg/.style={vtx,fill=okgreen!15,draw=okgreen!80!black},
  cl/.style={rounded corners=3pt,draw=teal!70!black,fill=teal!7,thick,inner sep=3pt,
    font=\small,align=center,minimum height=7mm},
  clh/.style={cl,fill=orange!14,draw=orange!80!black},
  clg/.style={cl,fill=okgreen!12,draw=okgreen!80!black},
  clr/.style={cl,fill=red!8,draw=badred,dashed},
  tedge/.style={thick,draw=teal!60!black},
  dedge/.style={thick,draw=teal!60!black,dashed},
  wlab/.style={font=\scriptsize,fill=white,inner sep=1.2pt},
  gedge/.style={thick,draw=black!70},
  newedge/.style={very thick,draw=badred,dashed},
  boxnode/.style={draw=accent!70,fill=accentlight!60,rounded corners=2pt,thick,
    align=center,font=\small,inner sep=4pt},
  arr/.style={-{Stealth[length=2.4mm]},thick},
}

%% ---------------------------------------------------------------------
%% Macros
%% ---------------------------------------------------------------------
\newcommand{\yes}{\textsf{yes}}
\newcommand{\no}{\textsf{no}}
\newcommand{\op}[1]{\textsf{#1}}
% Typewriter identifiers that may break after _ / . : , and before capitals
\ExplSyntaxOn
\str_new:N \l__mtp_code_str
\bool_new:N \l__mtp_code_first_bool
\cs_new_protected:Npn \mtp_breakable_tt:n #1
  {
    \str_set:Nn \l__mtp_code_str {#1}
    \mode_leave_vertical:
    \group_begin:
    \ttfamily
    \bool_set_true:N \l__mtp_code_first_bool
    \str_map_inline:Nn \l__mtp_code_str
      {
        \bool_if:NF \l__mtp_code_first_bool
          { \str_if_in:nnT { ABCDEFGHIJKLMNOPQRSTUVWXYZ } {##1} { \penalty 300 \scan_stop: } }
        \bool_set_false:N \l__mtp_code_first_bool
        ##1
        \str_if_in:nnT { _/.:, } {##1} { \penalty 50 \scan_stop: }
      }
    \group_end:
  }
\DeclareRobustCommand{\code}[1]{\mtp_breakable_tt:n {#1}}
\DeclareRobustCommand{\file}[1]{\mtp_breakable_tt:n {#1}}
\ExplSyntaxOff
\newcommand{\Kset}{\mathcal{K}}
\newcommand{\KG}{\mathcal{K}_G}
\newcommand{\Tu}{T_u}
\newcommand{\Tv}{T_v}
\newcommand{\polylog}{\operatorname{polylog}}
\newcommand{\size}{\operatorname{size}}
\newcommand{\wt}{\operatorname{wt}}
\newcommand{\repo}{\href{https://github.com/khemsinghbhati/MTP-3rd-sem}{\nolinkurl{github.com/khemsinghbhati/MTP-3rd-sem}}}
\newcommand{\status}[1]{\textbf{#1}}
\newcommand{\ok}{\textcolor{okgreen}{\textbf{Achieved}}}
\newcommand{\notok}{\textcolor{badred}{\textbf{Not achieved}}}
\newcommand{\partly}{\textcolor{warnorange}{\textbf{Partly}}}

\begin{document}
%% ---- end expanded/00_preamble.tex ----
%% ---- begin expanded/01_front.tex ----
%% =====================================================================
%% FRONT MATTER
%% =====================================================================
\pagenumbering{roman}
\hypersetup{pageanchor=false}
\begin{titlepage}
\newgeometry{top=20mm,bottom=20mm,left=22mm,right=22mm}
\centering
{\large\scshape Indian Institute of Technology Bhubaneswar\par}
\vspace{1mm}
{\normalsize School of Electrical and Computer Sciences\\
Department of Computer Science and Engineering\par}
\vspace{8mm}
{\color{accent}\rule{\textwidth}{1.2pt}}\par
\vspace{6mm}
{\large M.Tech Project (MTP) Report --- Third Semester\par}
\vspace{7mm}
{\color{accent}\fontsize{25}{31}\selectfont\bfseries Fully Dynamic Algorithms for\\[1mm]
Chordal Graphs and Split Graphs\par}
\vspace{5mm}
{\Large\itshape A detailed, tested C++ implementation of the core algorithms in\\ Louis Ibarra,
\emph{ACM Transactions on Algorithms} 4(4), Article 40 (2008)\par}
\vspace{6mm}
{\color{accent}\rule{\textwidth}{1.2pt}}\par
\vspace{9mm}
{\normalsize A report submitted in partial fulfilment of the requirements for the degree of\par}
\vspace{1mm}
{\large\bfseries Master of Technology\\ in\\ Computer Science and Engineering\par}
\vspace{9mm}
{\normalsize Submitted by\par}
\vspace{1mm}
{\LARGE\bfseries Khem Singh\par}
\vspace{1mm}
{\large Roll No.\ \texttt{25cs06004}\par}
\vspace{8mm}
{\normalsize Under the guidance of\par}
\vspace{1mm}
{\LARGE\bfseries Dr.\ Manas Jyoti Kashyop\par}
\vspace{1mm}
{\normalsize Assistant Professor, School of Electrical and Computer Sciences\\
Indian Institute of Technology Bhubaneswar\par}
\vspace{10mm}
\begin{tcolorbox}[enhanced,colback=accentlight!60,colframe=accent,boxrule=1pt,
  width=0.9\textwidth,halign=center,title={\centering Complete source code, tests, benchmarks and data},
  fonttitle=\bfseries\normalsize,coltitle=white]
{\large \href{https://github.com/khemsinghbhati/MTP-3rd-sem}{\nolinkurl{https://github.com/khemsinghbhati/MTP-3rd-sem}}}\\[1mm]
{\small Frozen algorithm core: \href{https://github.com/khemsinghbhati/MTP-3rd-sem/tree/2bf723c2b3b57a6b532058f21742b8125292852a}{\texttt{2bf723c2b3b5\ldots}} (branch \texttt{freeze/core-2026-11-21});\\
benchmarks run against \texttt{main} at commit \texttt{651b236}.}
\end{tcolorbox}
\vfill
{\large November 2026\par}
\restoregeometry
\end{titlepage}
\hypersetup{pageanchor=true}

%% ---------------------------------------------------------------------
\chapter*{Certificate}
\addcontentsline{toc}{chapter}{Certificate}
\thispagestyle{plain}
\noindent\textit{Draft for supervisor review; not signed or certified.}\par
\vspace{4mm}
\noindent This is to certify that the project report entitled
\textbf{``Fully Dynamic Algorithms for Chordal Graphs and Split Graphs: a detailed C++
implementation and verification of the core algorithms in Ibarra (2008)''}, submitted by \textbf{Khem Singh} (Roll No.\ \texttt{25cs06004})
to the Indian Institute of Technology Bhubaneswar in partial fulfilment of the requirements for
the degree of Master of Technology in Computer Science and Engineering, is a record of the
M.Tech project work carried out under my supervision and guidance during the third semester.
The implementation described in this report is available in the public repository
\repo.

\vspace{30mm}
\noindent\begin{tabular}{@{}l@{}}
\rule{65mm}{0.4pt}\\
\textbf{Dr.\ Manas Jyoti Kashyop}\\
Assistant Professor\\
School of Electrical and Computer Sciences\\
Indian Institute of Technology Bhubaneswar
\end{tabular}
\hfill
\begin{tabular}{@{}l@{}}
Place: Bhubaneswar\\
Date: \rule{34mm}{0.4pt}
\end{tabular}

%% ---------------------------------------------------------------------
\chapter*{Declaration}
\addcontentsline{toc}{chapter}{Declaration}
\thispagestyle{plain}
\noindent\textit{Student declaration --- signature and date pending.}\par
\vspace{4mm}
\noindent I declare that the work presented in this report was carried out by me under the
guidance of Dr.\ Manas Jyoti Kashyop, and that every published algorithm, data structure,
theorem and result used in it is attributed to its original source and cited. Statements that
are \emph{not} taken from a published source --- in particular the implementation-level design
choices, the measurements, and the analysis of the preprocessing claim in
\cref{ch:preprocessing} --- are explicitly marked as such. The algorithmic core of the
implementation is preserved unchanged at commit
\texttt{2bf723c2b3b57a6b532058f21742b8125292852a} of the repository \repo.

\vspace{15mm}
\noindent\begin{tabular}{@{}l@{}}
\textbf{Khem Singh}\\
Roll No.\ \texttt{25cs06004}\\
M.Tech, Computer Science and Engineering\\
Indian Institute of Technology Bhubaneswar
\end{tabular}
\hfill
\begin{tabular}{@{}l@{}}
Place: Bhubaneswar\\
Date: \rule{34mm}{0.4pt}
\end{tabular}

%% ---------------------------------------------------------------------
\chapter*{Acknowledgements}
\addcontentsline{toc}{chapter}{Acknowledgements}
\thispagestyle{plain}
I am grateful to my supervisor, Dr.\ Manas Jyoti Kashyop, for the guidance, encouragement
and feedback given throughout this project. I also thank the faculty and staff of the School of
Electrical and Computer Sciences, IIT Bhubaneswar, for the facilities and the academic
environment that made this work possible.

\vspace{10mm}
\hfill Khem Singh

%% ---------------------------------------------------------------------
\chapter*{Abstract}
\addcontentsline{toc}{chapter}{Abstract}
\thispagestyle{plain}
A graph is \emph{chordal} if every cycle of length four or more has a chord. Chordal graphs
appear in sparse Gaussian elimination, in relational databases, in phylogenetics and in
probabilistic graphical models, and many problems that are NP-hard in general are solvable in
linear time on them. A \emph{fully dynamic} algorithm for chordal graphs maintains a chordal
graph under edge insertions and deletions and answers, before each update, whether the update
keeps the graph chordal. Louis Ibarra (ACM Transactions on Algorithms, 2008) gave the first
such algorithm. It keeps a \emph{clique tree} of the graph --- a tree whose nodes are the
maximal cliques --- and decides each update by a purely local test on that tree.

This project implements the four principal dynamic chordal-graph variants and the
dynamic split-graph procedures of Ibarra (2008) in C++17. This report explains their
implementation from definitions and algorithmic invariants to the corresponding code.
The optional ET-tree connectivity extension discussed in the paper is not implemented,
and the published A2 preprocessing target is not achieved, as detailed in
\cref{ch:preprocessing}. The implementation contains:
(i) the linear-time static preprocessing (Maximum Cardinality Search, perfect-elimination
check, maximal cliques and the Blair--Peyton clique tree);
(ii) Ibarra's first implementation \textbf{A1}, with a sparse clique representation that
achieves the paper's worst-case bounds of \(O(m+n)\) preprocessing and \(O(n)\) per query and
update;
(iii) the second implementation \textbf{A2}, built on a from-scratch implementation of the
\emph{worst-case} Sleator--Tarjan dynamic tree (solid paths and path sets stored as
Bent--Sleator--Tarjan globally biased binary trees), giving \(O(\log^2 n)\) insertion queries,
\(O(n)\) insertions and \(O(n\log n)\) deletions;
(iv) the faster deletion query of Section~6 (\textbf{A1\(^+\)}, \textbf{A2\(^+\)}), based on a
modified version of Yellin's dynamic set-intersection structure, with \(O(\sqrt m)\) queries; and
(v) Ibarra's constant-time dynamic recognition of \emph{split graphs} via the Hammer--Simeone
degree-sequence theorem, together with the clique, chromatic, independence and clique-cover
numbers.

Correctness is supported by a suite of 23 automated tests, including exhaustive checking on all
19{,}048 labelled chordal graphs with at most six vertices, randomised legal update sequences
checked against independent oracles, differential testing of A2 against A1, structural
invariant checks of the dynamic trees, sanitizer builds and independent Python/NetworkX audits.

The report is also explicit about what is \emph{not} achieved. The paper's summary table claims
\(O(m+n\log n)\) preprocessing for A2, while the paper's own Section~5.1 builds the
\(n\times n\) intersection matrix \(W\) in \(O(n^3)\) or \(O(mn)\) time. The implementation
follows Section~5.1 and therefore preprocesses A2 in \(O(n^2+mn)\) time. We show that any
algorithm that writes an explicit \(W\) needs \(\Omega(n^2)\) time on some graphs with \(m=0\),
and we give a conditional argument, based on Boolean matrix multiplication, that achieving both
the claimed preprocessing time and the claimed \(O(\log^2 n)\) insertion query for every chordal
graph would require a major algorithmic breakthrough. These findings are reported as limitations
rather than hidden.

\vspace{3mm}
\noindent\textbf{Keywords:} chordal graphs, clique trees, fully dynamic graph algorithms,
dynamic trees, biased search trees, split graphs, perfect elimination ordering.

%% ---------------------------------------------------------------------
\cleardoublepage
\phantomsection
\addcontentsline{toc}{chapter}{Contents}
\tableofcontents
\cleardoublepage
\phantomsection
\addcontentsline{toc}{chapter}{List of Figures}
\listoffigures
\cleardoublepage
\phantomsection
\addcontentsline{toc}{chapter}{List of Tables}
\listoftables
\cleardoublepage
\phantomsection
\addcontentsline{toc}{chapter}{List of Algorithms}
\listofalgorithms

%% ---------------------------------------------------------------------
\chapter*{How to Read This Report}
\addcontentsline{toc}{chapter}{How to Read This Report}
\markboth{How to Read This Report}{How to Read This Report}
\thispagestyle{plain}

The report is written for a reader who knows basic graph theory and data structures (graphs,
trees, breadth-first search, balanced binary search trees, big-\(O\) notation) but has
\emph{not} read Ibarra's paper. Every concept that the implementation depends on is defined
before it is used, and every algorithm is shown three times: as mathematics, as pseudocode, and
as the actual C++ code from the repository (with the real file name and line numbers, so that
the reader can open the file and follow along).

\paragraph{Suggested paths through the report.}
\begin{itemize}
\item \textbf{To understand the algorithm:} \cref{ch:background} (background), then
  \cref{ch:ibarra} (Ibarra's method, with small worked examples), then
  \cref{ch:a1} (the simplest complete implementation, A1, with traces of real program runs).
\item \textbf{To review the code:} \cref{ch:architecture} (repository and module map), then
  \cref{ch:static,ch:a1,ch:a2,ch:fastdelete,ch:split}; \cref{app:files} lists every file.
\item \textbf{To check the complexity claims:} the ``Complexity'' sections at the end of
  \cref{ch:static,ch:a1,ch:a2,ch:fastdelete,ch:split}, then \cref{ch:preprocessing} (the one
  claim of the paper that is not achieved) and the summary in \cref{tab:final-status}.
\item \textbf{To reproduce the results:} \cref{app:build} (build and run guide),
  \cref{ch:verification} (tests) and \cref{ch:experiments} (measurements).
\end{itemize}

\paragraph{Boxes used in the text.}
\begin{keyidea}[Key idea]
Blue boxes state the one idea that makes a section work.
\end{keyidea}
\begin{incode}
Green boxes point to the exact files, classes and functions that implement the section.
\end{incode}
\begin{paperbox}
Grey boxes summarise what Ibarra's paper (or another cited source) states, with its section
and page number, so that the reader can separate the paper's claims from ours.
\end{paperbox}
\begin{caution}
Orange boxes mark places where the implementation deliberately differs from the literal text of
the paper, or where a claim of the paper is not achieved.
\end{caution}

\paragraph{Page references to the paper.} Ibarra's paper is Article~40 of volume~4 of
\emph{ACM Transactions on Algorithms}; its pages are numbered 40:1--40:20. A reference such as
``Ibarra \S5.1, p.\,40:13'' therefore means Section~5.1 on page~13 of that article.

\clearpage
\section*{Notation}
\addcontentsline{toc}{section}{Notation}
\begin{longtable}{@{}p{0.2\textwidth}p{0.74\textwidth}@{}}
\toprule
\textbf{Symbol} & \textbf{Meaning}\\
\midrule
\endhead
\(G=(V,E)\) & the current (chordal) graph; simple and undirected\\
\(n,\ m\) & \(n=|V|\), \(m=|E|\) (current number of edges)\\
\(N(v),\ d(v)\) & neighbourhood and degree of vertex \(v\)\\
\(uv\) or \(\{u,v\}\) & the (possible) edge between \(u\) and \(v\)\\
\(G+uv,\ G-uv\) & \(G\) with edge \(uv\) inserted / deleted\\
\(\KG\) & the set of maximal cliques of \(G\); \(c=|\KG|\le n\)\\
\(\KG(v)\) & the maximal cliques that contain vertex \(v\)\\
\(T\) & a clique tree of \(G\) (nodes are maximal cliques)\\
\(x,y,z,w\) & nodes of \(T\); \(K_x\) is the maximal clique stored at node \(x\)\\
\(w(x,y)\) & weight of tree edge \(\{x,y\}\): \(w(x,y)=|K_x\cap K_y|\)\\
\(W(x,y)\) & the matrix of A2: \(W(x,y)=|K_x\cap K_y|\) for \emph{all} pairs of nodes\\
\(T_v\) & the subtree of \(T\) induced by the cliques containing \(v\)\\
\(N_T(x)\) & the neighbours of node \(x\) in the clique tree\\
PEO & perfect elimination ordering\\
MCS & maximum cardinality search\\
ST83 & Sleator and Tarjan, ``A data structure for dynamic trees'' (1983)\\
BST85 & Bent, Sleator and Tarjan, ``Biased search trees'' (1985)\\
A1, A2 & Ibarra's first and second implementations (\S4 and \S5 of the paper)\\
A1\(^+\), A2\(^+\) & A1 and A2 with the faster Delete-Query structure of \S6\\
\(\log\) & logarithm to base 2\\
\S & ``Section'' of a paper; \S5.1 of Ibarra's paper is its Section~5.1\\
\bottomrule
\end{longtable}

\cleardoublepage
\pagenumbering{arabic}
%% ---- end expanded/01_front.tex ----
%% ---- begin expanded/02_intro.tex ----
%% =====================================================================
%% Reusable drawing of the 7-vertex example graph (data/example_7_vertices.txt)
%% =====================================================================
\newcommand{\ExGraphNodes}{%
  \node[vtx] (g0) at (0,1.6) {0};
  \node[vtx] (g3) at (1.6,1.6) {3};
  \node[vtx] (g1) at (0,0) {1};
  \node[vtx] (g2) at (1.6,0) {2};
  \node[vtx] (g6) at (3.1,1.6) {6};
  \node[vtx] (g4) at (3.1,0.15) {4};
  \node[vtx] (g5) at (2.45,-1.15) {5};
}
\newcommand{\ExGraphEdges}{%
  \draw[gedge] (g0)--(g1) (g1)--(g3) (g0)--(g2) (g2)--(g3) (g2)--(g4) (g2)--(g5)
               (g4)--(g5) (g3)--(g6) (g0)--(g3) (g1)--(g2);
}

\chapter{Introduction}\label{ch:intro}

\noindent\textbf{AI Assistance Declaration.}
AI assistance was used to identify relevant resources and research papers, better understand their methods, and support implementation, testing, and documentation. The project work was verified by a human reviewer against the research paper's requirements. Any deviations from the paper's published guarantees, including the A2 preprocessing bound, are explicitly documented.

\section{Dynamic graph algorithms and graph classes}
Many graphs that arise in practice change over time: links are added to and removed from a
network, constraints are added to and removed from a model. A \emph{dynamic graph algorithm}
maintains information about a graph while the graph is modified by a sequence of updates, and
it tries to process each update much faster than recomputing everything from scratch. An
algorithm is \emph{fully dynamic} if it supports both insertions and deletions of edges.

A natural dynamic question is whether a graph \emph{stays inside a graph class}. If an
application relies on a structural property --- for example, that its graph is chordal, so that
a sparse matrix can be factorised without fill-in, or so that a probabilistic model is
decomposable --- then before applying a change it must know whether the change destroys that
property. Ibarra's paper~\cite{ibarra2008} solves this question for two important classes,
\emph{chordal graphs} and \emph{split graphs}. For chordal graphs it supports the four
operations in \cref{tab:ops}.

\begin{table}[ht]
\centering
\caption[The four operations]{The four operations of a fully dynamic chordal graph algorithm
(Ibarra \S1.3, p.\,40:3). The graph \(G\) is always chordal.}\label{tab:ops}
\begin{tabularx}{\textwidth}{@{}lX@{}}
\toprule
\textbf{Operation} & \textbf{Meaning}\\
\midrule
\op{Delete-Query}\((u,v)\) & For an edge \(uv\in E\): answer \yes\ if \(G-uv\) is chordal, otherwise \no.\\
\op{Delete}\((u,v)\) & If \(G-uv\) is chordal, delete \(uv\) and update the data structure; otherwise reject.\\
\op{Insert-Query}\((u,v)\) & For a non-edge \(uv\notin E\): answer \yes\ if \(G+uv\) is chordal, otherwise \no.\\
\op{Insert}\((u,v)\) & If \(G+uv\) is chordal, insert \(uv\) and update the data structure; otherwise reject.\\
\bottomrule
\end{tabularx}
\end{table}

Recomputing chordality from scratch after every change costs \(\Theta(n+m)\) time (with the
classic linear-time recognition algorithms of Rose, Tarjan and Lueker~\cite{rose1976} and
Tarjan and Yannakakis~\cite{tarjan1984}). Since a chordal graph can have \(\Theta(n^2)\) edges,
this is \(\Theta(n^2)\) in the worst case. Ibarra's algorithm answers every query and performs
every update in \(O(n)\) worst-case time, and its second implementation answers insertion
queries in \(O(\log^2 n)\) time.

\section{Why chordal graphs?}
A chordal graph is one in which every cycle of length at least four has a \emph{chord} (an edge
joining two non-consecutive vertices of the cycle). Equivalently (\cref{ch:background}), its
vertices can be eliminated one by one so that each eliminated vertex's remaining neighbours
form a clique. This one property explains why the class is so useful:
\begin{itemize}
\item In \textbf{sparse Gaussian elimination}, a symmetric matrix whose nonzero pattern is a
  chordal graph can be factorised with no fill-in, and the elimination order is exactly a
  perfect elimination ordering~\cite{rose1976,blair1993}.
\item In \textbf{relational databases} and \textbf{graphical models}, chordal graphs
  correspond to acyclic hypergraphs and decomposable models, and the clique tree is the
  ``join tree'' or ``junction tree''~\cite{tarjan1984}.
\item Chordal graphs are \textbf{perfect}: the clique number equals the chromatic number, and
  maximum clique, colouring, maximum independent set and minimum clique cover are all solvable
  in linear time~\cite{golumbic1980}, although they are NP-hard on general graphs.
\end{itemize}

\section{The idea of Ibarra's algorithm in one picture}
Every chordal graph has a \emph{clique tree}: a tree whose nodes are the maximal cliques of the
graph and in which, for every vertex \(v\), the cliques that contain \(v\) form a connected
subtree (\cref{sec:clique-trees}). Ibarra's insight is that the four questions of
\cref{tab:ops} can be answered by looking at this tree, and that after an accepted update the
tree changes only \emph{locally} --- one clique splits, or one new clique appears and at most
two old ones disappear. \Cref{fig:overview} shows the example graph used throughout this report
and its clique tree.

\begin{figure}[ht]
\centering
\begin{tikzpicture}
\begin{scope}
\ExGraphNodes
\ExGraphEdges
\node[font=\small\bfseries,anchor=south] at (1.5,2.1) {graph \(G\) (7 vertices, 10 edges)};
\end{scope}
\begin{scope}[xshift=7.2cm,yshift=0.2cm]
\node[cl] (K0) at (0,1.8) {\(K_0=\{3,6\}\)};
\node[cl] (K1) at (0,0.5) {\(K_1=\{0,1,2,3\}\)};
\node[cl] (K2) at (0,-0.8) {\(K_2=\{2,4,5\}\)};
\draw[tedge] (K0)--node[wlab,right=2pt]{\(|K_0\cap K_1|=1\)}(K1);
\draw[tedge] (K1)--node[wlab,right=2pt]{\(|K_1\cap K_2|=1\)}(K2);
\node[font=\small\bfseries,anchor=south] at (0,2.3) {clique tree \(T\)};
\end{scope}
\draw[arr,accent] (3.7,0.6) -- node[above,font=\small,align=center]{preprocess\\\(O(m+n)\)} (5.6,0.6);
\end{tikzpicture}
\caption[The running example and its clique tree]{The running example (file \file{data/example_7_vertices.txt}) and the clique tree that
the program computes for it. Each tree edge is labelled with the size of the intersection of the
two cliques it joins. All four dynamic operations are answered from \(T\).}
\label{fig:overview}
\end{figure}

For example, inserting the non-edge \(\{0,4\}\) keeps the graph chordal, and the program
answers \yes; inserting \(\{4,6\}\) would create the chordless cycle \(4\!-\!2\!-\!3\!-\!6\!-\!4\)
and the program answers \no. \Cref{ch:ibarra} explains how
both answers are read off the tree, and \cref{sec:a1-traces} shows the actual program output.

\section{Ibarra's results}
\Cref{tab:paper-bounds} reproduces the running-time summary of the paper. All bounds are
\emph{worst-case} (Ibarra \S1.3, p.\,40:3): they hold for every single operation, not on average over a
sequence and not with high probability.

\begin{table}[ht]
\centering
\caption[Running times claimed by Ibarra]{Running times claimed by Ibarra (\S7, p.\,40:17). \(n\) and \(m\) are the current
numbers of vertices and edges. All bounds are worst-case.}\label{tab:paper-bounds}
\begin{tabular}{@{}lccccc@{}}
\toprule
& \op{Delete-Query} & \op{Delete} & \op{Insert-Query} & \op{Insert} & preprocessing\\
\midrule
A1 & \(O(n)\) & \(O(n)\) & \(O(n)\) & \(O(n)\) & \(O(m+n)\)\\
A2 & \(O(n)\) & \(O(n\log n)\) & \(O(\log^2 n)\) & \(O(n)\) & \(O(m+n\log n)\)\\
A1\(^+\) & \(O(\sqrt m)\) & \(O(m+n)\) & \(O(n)\) & \(O(m+n)\) & \(O(mn+n^2)\)\\
A2\(^+\) & \(O(\sqrt m)\) & \(O(m+n\log n)\) & \(O(\log^2 n)\) & \(O(m+n)\) & \(O(mn+n^2)\)\\
\bottomrule
\end{tabular}
\end{table}

For split graphs (\S8 of the paper), Ibarra shows how to maintain, on an arbitrary graph, in
\(O(1)\) time per edge insertion or deletion, whether the graph is split, and --- when it is ---
its clique number, chromatic number, independence number and clique cover number, after
\(O(m+n)\) preprocessing.

\section{Objectives of the project}
The objectives, fixed at the start of the project and recorded in the repository
\file{README.md}, were:
\begin{enumerate}
\item Implement \emph{every} algorithm of the paper --- A1, A2, A1\(^+\), A2\(^+\) and the split
  graph algorithm --- using the methods the paper describes, and the published data structures
  it cites (in particular the worst-case Sleator--Tarjan dynamic tree~\cite{sleator1983}, not an
  amortised or randomised substitute).
\item Achieve the paper's \emph{worst-case} time bounds in the implementation, and argue each
  bound from the code, not from timing measurements.
\item Never present an invented algorithm as a published one: where the implementation
  composes published results in a way the paper does not literally describe, say so.
\item Verify correctness thoroughly against independent oracles.
\item Report honestly any claim of the paper that the implementation does not achieve.
\end{enumerate}

\section{Summary of outcomes}
\Cref{tab:status-intro} summarises the outcome; \cref{tab:final-status} in the conclusion gives
the detailed version with references to the relevant sections.

\begin{table}[ht]
\centering
\caption[Outcome of the project, by component]{Outcome of the project, by component. The
numbers in brackets are the sections of Ibarra's paper that describe each component
(\S\ = ``Section'').}\label{tab:status-intro}
\begin{tabularx}{\textwidth}{@{}lXl@{}}
\toprule
\textbf{Component} & \textbf{Implemented bound (worst case)} & \textbf{Status}\\
\midrule
Static preprocessing & MCS, PEO test, cliques, clique tree: \(O(n+m)\) & \ok\\
A1 (\S3--4) & \(O(n+m)\) preprocessing; all four operations \(O(n)\) & \ok\\
A2 operations (\S5) & \op{Insert-Query} \(O(\log^2 n)\), \op{Insert} \(O(n)\), \op{Delete} \(O(n\log n)\), \op{Delete-Query} \(O(n)\) & \ok\\
A2 preprocessing (\S5.1, \S7) & \(O(n^2+mn)\) instead of the claimed \(O(m+n\log n)\) & \notok\\
A1\(^+\)/A2\(^+\) (\S6) & \op{Delete-Query} \(O(\sqrt m)\); updates \(O(m+n)\) extra; \(O(mn+n^2)\) preprocessing & \ok\\
Split graphs (\S8) & \(O(1)\) per update and query; \(O(n+m)\) preprocessing & \ok\\
\bottomrule
\end{tabularx}
\end{table}

\section{Contributions of this project}
The algorithms are Ibarra's; the contribution of the project is a complete and carefully
verified implementation, together with the following implementation-level results, each of
which is stated and proved in this report.
\begin{enumerate}
\item A \textbf{sparse representation for A1} (\cref{ch:a1}). The paper stores every maximal
  clique as a characteristic vector of length \(n\). Allocating these vectors costs
  \(\Theta(cn)\), which is \(\Theta(n^2)\) for a graph with many small cliques --- more than
  the claimed \(O(m+n)\) preprocessing. The implementation stores clique member lists plus
  inverse lists and meets the \(O(m+n)\) bound exactly.
\item A \textbf{from-scratch implementation of the worst-case Sleator--Tarjan dynamic
  tree}~\cite{sleator1983} with Bent--Sleator--Tarjan \emph{globally biased} binary
  trees~\cite{bent1985} (\cref{ch:a2}), including the linear-time bottom-up construction of
  BST85.
\item An \textbf{\(O(n)\) worst-case A2 insertion} even when an insertion merges two
  high-degree cliques (\cref{sec:a2-insert-on}). A literal implementation would perform
  \(\Theta(n)\) dynamic-tree operations, i.e.\ \(\Theta(n\log n)\) time; the implementation
  rebuilds the same Sleator--Tarjan structure in \(O(n)\) using a linear-time construction of
  globally biased trees (attributed in the code to BST85; see \cref{sec:bst85}).
\item An \textbf{analysis of the A2 preprocessing claim} (\cref{ch:preprocessing}): an
  \(\Omega(n^2)\) bound for every method that writes an explicit matrix \(W\), and a
  conditional barrier via Boolean matrix multiplication for every method whatsoever.
\item A \textbf{verification suite} (\cref{ch:verification}) and \textbf{reproducible
  experiments} (\cref{ch:experiments}).
\end{enumerate}

\section{Organisation of the report}
\Cref{ch:background} introduces chordal graphs, perfect elimination orderings, clique trees,
split graphs and dynamic trees. \Cref{ch:ibarra} presents Ibarra's algorithm mathematically.
\Cref{ch:architecture} describes the repository and its modules. \Cref{ch:static} covers the
static preprocessing, \cref{ch:a1} the A1 implementation, \cref{ch:a2} the A2 implementation
with the Sleator--Tarjan dynamic tree, \cref{ch:fastdelete} the faster deletion query, and
\cref{ch:split} the split graph algorithm. \Cref{ch:preprocessing} analyses the A2
preprocessing claim. \Cref{ch:verification,ch:experiments} present testing and experiments, and
\cref{ch:conclusion} concludes. The appendices contain a build guide, a file-by-file reference,
complete program traces and a glossary.
%% ---- end expanded/02_intro.tex ----
%% ---- begin expanded/03_background.tex ----
%% =====================================================================
\chapter{Background}\label{ch:background}
%% =====================================================================
This chapter collects the graph theory and data-structure background that the rest of the
report uses. Everything here is standard; the sources are cited for each result. Readers
familiar with chordal graphs and clique trees can skim to \cref{sec:dummy-edges}, which
describes a convention of Ibarra's paper that is essential for the implementation.

\section{Basic notation}
All graphs are finite, simple and undirected. For a graph \(G=(V,E)\) we write \(n=|V|\) and
\(m=|E|\); in the implementation vertices are the integers \(0,1,\dots,n-1\). The
\emph{neighbourhood} of \(v\) is \(N(v)=\{u : uv\in E\}\) and its degree is \(d(v)=|N(v)|\).
For \(S\subseteq V\), \(G[S]\) is the subgraph \emph{induced} by \(S\) (all edges of \(G\) with
both ends in \(S\)), and \(G-S=G[V\setminus S]\).

A \emph{clique} is a set of pairwise adjacent vertices. A clique is \emph{maximal} if it is not
contained in a larger clique, equivalently if no vertex outside it is adjacent to all of its
vertices. We write \(\KG\) for the set of maximal cliques of \(G\) and \(\KG(v)\) for those that
contain \(v\). A \emph{cycle} \(v_1v_2\cdots v_kv_1\) (\(k\ge 3\)) has a \emph{chord} if two
non-consecutive vertices of the cycle are adjacent.

\section{Chordal graphs}
\begin{definition}[Chordal graph]
A graph is \emph{chordal} (also called \emph{triangulated}) if every cycle of length at least
four has a chord. Equivalently, \(G\) has no \emph{induced} cycle \(C_k\) with \(k\ge 4\).
\end{definition}

\begin{figure}[ht]
\centering
\begin{subfigure}[b]{0.42\textwidth}\centering
\begin{tikzpicture}
\ExGraphNodes
\ExGraphEdges
\end{tikzpicture}
\caption{Chordal: the running example.}\label{fig:ex-chordal}
\end{subfigure}\hfill
\begin{subfigure}[b]{0.25\textwidth}\centering
\begin{tikzpicture}
\node[vtx] (a) at (0,1.4) {0};\node[vtx] (b) at (1.4,1.4) {1};
\node[vtx] (c) at (1.4,0) {2};\node[vtx] (d) at (0,0) {3};
\draw[gedge] (a)--(b)--(c)--(d)--(a);
\end{tikzpicture}
\caption{Not chordal: \(C_4\).}\label{fig:c4}
\end{subfigure}\hfill
\begin{subfigure}[b]{0.3\textwidth}\centering
\begin{tikzpicture}
\node[vtx] (a) at (0,1.4) {0};\node[vtx] (b) at (1.4,1.4) {1};
\node[vtx] (c) at (1.4,0) {2};\node[vtx] (d) at (0,0) {3};
\draw[gedge] (a)--(b)--(c)--(d)--(a);
\draw[newedge] (a)--(c);
\end{tikzpicture}
\caption{Adding a chord makes it chordal.}\label{fig:c4chord}
\end{subfigure}
\caption[Chordal and non-chordal graphs]{Chordal and non-chordal graphs. The 4-cycle in (b) is the smallest non-chordal graph;
the program \file{mtp_cli} rejects the file \file{data/nonchordal_cycle4.txt}, which contains
exactly this graph.}
\label{fig:chordal-examples}
\end{figure}

In the running example (\cref{fig:ex-chordal}) the vertices \(\{0,1,2,3\}\) form a clique of
size four, \(\{2,4,5\}\) a triangle and \(\{3,6\}\) an edge; every cycle of length four or more
stays inside the 4-clique or passes through it in a way that has a chord, so the graph is
chordal. Two facts are used repeatedly:
\begin{itemize}
\item Every \emph{induced} subgraph of a chordal graph is chordal (a chordless cycle in
  \(G[S]\) is a chordless cycle in \(G\)).
\item Deleting an \emph{edge} or inserting an edge can destroy chordality; that is why the
  queries of \cref{tab:ops} are needed. For instance, deleting the chord of
  \cref{fig:c4chord} produces \cref{fig:c4}.
\end{itemize}

\section{Perfect elimination orderings}\label{sec:peo}
\begin{definition}[Simplicial vertex, PEO]
A vertex is \emph{simplicial} if its neighbourhood is a clique. An ordering
\(\sigma=(v_1,v_2,\dots,v_n)\) of \(V\) is a \emph{perfect elimination ordering} (PEO) if, for
every \(i\), the set of \emph{later neighbours}
\[ L(v_i)=\{v_j\in N(v_i) : j>i\} \]
is a clique; that is, every \(v_i\) is simplicial in \(G[\{v_i,\dots,v_n\}]\).
\end{definition}

\begin{theorem}[Fulkerson and Gross~\cite{fulkerson1965}]\label{thm:fg}
A graph is chordal if and only if it has a perfect elimination ordering.
\end{theorem}

The name comes from Gaussian elimination: eliminating the variables of a symmetric sparse
system in a PEO order creates no new nonzero entries. \Cref{fig:peo} shows the PEO that the
implementation computes for the running example.

\begin{figure}[ht]
\centering
\begin{tikzpicture}[x=1.75cm]
\foreach \v/\i in {4/0,5/1,1/2,2/3,0/4,3/5,6/6}{
  \node[vtx] (p\v) at (\i,0) {\v};
  \node[font=\scriptsize,below=1mm of p\v] {\(v_{\the\numexpr\i+1\relax}\)};
}
% later-neighbour arcs (drawn above): only from each vertex to its later neighbours
\draw[gedge,accent!80] (p4) to[bend left=35] (p5);
\draw[gedge,accent!80] (p4) to[bend left=45] (p2);
\draw[gedge,accent!80] (p5) to[bend left=40] (p2);
\draw[gedge,accent!80] (p1) to[bend left=30] (p2);
\draw[gedge,accent!80] (p1) to[bend left=40] (p0);
\draw[gedge,accent!80] (p1) to[bend left=45] (p3);
\draw[gedge,accent!80] (p2) to[bend left=30] (p0);
\draw[gedge,accent!80] (p2) to[bend left=40] (p3);
\draw[gedge,accent!80] (p0) to[bend left=30] (p3);
\draw[gedge,accent!80] (p3) to[bend left=30] (p6);
\end{tikzpicture}
\caption[A perfect elimination ordering of the running example]{The ordering \(\sigma=(4,5,1,2,0,3,6)\) printed by \file{mtp_cli} for the running
example, with every edge drawn as an arc. The later neighbours of each vertex form a clique:
\(L(4)=\{5,2\}\), \(L(5)=\{2\}\), \(L(1)=\{2,0,3\}\), \(L(2)=\{0,3\}\), \(L(0)=\{3\}\),
\(L(3)=\{6\}\), \(L(6)=\emptyset\). Hence \(\sigma\) is a PEO.}
\label{fig:peo}
\end{figure}

\section{Maximum cardinality search}\label{sec:mcs-background}
A PEO of a chordal graph can be found in linear time by \emph{maximum cardinality search}
(MCS)~\cite{tarjan1984}. MCS numbers the vertices from \(n\) down to \(1\): at each step it
selects an unnumbered vertex with the largest number of already-numbered neighbours (ties are
broken arbitrarily).

\begin{theorem}[Tarjan and Yannakakis~\cite{tarjan1984}]\label{thm:mcs}
If \(G\) is chordal, the order \(v_1,\dots,v_n\) produced by MCS (where \(v_n\) is selected
first) is a perfect elimination ordering. MCS can be implemented in \(O(n+m)\) time.
\end{theorem}

MCS alone does not tell us whether \(G\) is chordal --- it produces \emph{some} ordering for
every graph. Chordality is decided by testing whether that ordering is a PEO, which
\cref{lem:peo-test} below does in linear time.

\begin{lemma}[Linear-time PEO test~\cite{tarjan1984,rose1976}]\label{lem:peo-test}
For an ordering \(\sigma\), let the \emph{parent} \(p(u)\) be the earliest vertex of \(L(u)\)
(undefined if \(L(u)=\emptyset\)). Then \(\sigma\) is a PEO if and only if, for every vertex
\(u\) with \(L(u)\ne\emptyset\),
\[ L(u)\setminus\{p(u)\}\ \subseteq\ N(p(u)). \]
\end{lemma}
\begin{proof}[Proof sketch]
If \(\sigma\) is a PEO, \(L(u)\) is a clique containing \(p(u)\), so the condition holds.
Conversely, suppose the condition holds and, for contradiction, let \(u\) be the
\emph{latest} vertex whose \(L(u)\) is not a clique, with non-adjacent \(a,b\in L(u)\). Neither
equals \(p(u)\) (all of \(L(u)\setminus\{p(u)\}\) is adjacent to \(p(u)\)), so both lie in
\(L(p(u))\), which is a clique because \(p(u)\) is later than \(u\). This contradicts
\(ab\notin E\).
\end{proof}
The test needs only one marking pass per parent, so it costs \(O(n+m)\); the code is shown in
\cref{sec:peo-code}. If \(G\) is not chordal, \cref{thm:fg} guarantees that no ordering is a
PEO, so the test fails on the MCS ordering and the program reports ``not chordal''.

\section{Maximal cliques of a chordal graph}\label{sec:cliques-background}
For a PEO \(\sigma\), call \(C(v)=\{v\}\cup L(v)\) the \emph{candidate clique} of \(v\). Every
candidate is a clique, and every maximal clique is a candidate:

\begin{lemma}[Cliques from a PEO~\cite{fulkerson1965,golumbic1980}]\label{lem:cliques}
Let \(\sigma\) be a PEO of \(G\). For every maximal clique \(K\), if \(v\) is the earliest
vertex of \(K\) in \(\sigma\), then \(K=C(v)\). Consequently:
\begin{enumerate}
\item \(G\) has at most \(n\) maximal cliques, \(c=|\KG|\le n\);
\item \(\displaystyle\sum_{K\in\KG}|K|\;\le\;\sum_{v\in V}\bigl(1+|L(v)|\bigr)\;=\;n+m.\)
\end{enumerate}
\end{lemma}
\begin{proof}
All other vertices of \(K\) are later neighbours of \(v\), so \(K\subseteq C(v)\); since \(C(v)\)
is a clique and \(K\) is maximal, \(K=C(v)\). Distinct maximal cliques therefore have distinct
earliest vertices, which gives (1), and (2) follows because every edge is counted in exactly one
\(L(\cdot)\).
\end{proof}

Bound (2) is the reason why a sparse representation of the cliques fits in \(O(n+m)\) space
(\cref{ch:a1}). Ibarra also uses the following closely related bound (Ibarra Lemma~7,
p.\,40:14): for a chordal graph,
\(\sum_{v:\,d(v)>0}|\KG(v)|\le 2m\). Note that \(\sum_{K\in\KG}|K|=\sum_{v\in V}|\KG(v)|\),
because both count the pairs (vertex, maximal clique containing it).

Which candidates are maximal? The following test is the one implemented in
\file{src/maximal_cliques.cpp}.

\begin{lemma}[Parent subsumption test]\label{lem:subsumption}
Let \(\sigma\) be a PEO. The candidate \(C(p)\) is \emph{not} maximal if and only if there is a
vertex \(u\) with \(p(u)=p\) and \(|L(u)|=|L(p)|+1\).
\end{lemma}
\begin{proof}
For every \(u\) with parent \(p\) we have \(L(u)\setminus\{p\}\subseteq L(p)\) (the members of
\(L(u)\) other than \(p\) are later than \(p\) and adjacent to it), so \(|L(u)|\le|L(p)|+1\),
with equality exactly when \(L(u)=C(p)\), in which case \(C(p)\subsetneq C(u)\).
Conversely, suppose \(C(p)\) is not maximal. Among the vertices \(u\) with
\(C(p)\subseteq L(u)\) (there is one: the earliest vertex of a maximal clique strictly
containing \(C(p)\)), choose the \emph{latest}. If \(p(u)\ne p\), then \(p(u)\) is earlier than
\(p\), is not in \(C(p)\), and \(C(p)\subseteq L(u)\setminus\{p(u)\}\subseteq L(p(u))\) --- so
\(p(u)\) is a later choice, a contradiction. Hence \(p(u)=p\) and \(L(u)=C(p)\).
\end{proof}

\section{Clique trees}\label{sec:clique-trees}
\begin{definition}[Clique tree; Ibarra \S2, p.\,40:4]\label{def:clique-tree}
A \emph{clique tree} of \(G\) is a tree \(T\) whose node set is \(\KG\) and which has the
\emph{induced subtree property}: for every vertex \(v\in V\), the nodes whose cliques contain
\(v\) induce a connected subtree \(T_v\) of \(T\). Each tree edge \(\{x,y\}\) gets the
\emph{weight} \(w(x,y)=|K_x\cap K_y|\).
\end{definition}

\begin{theorem}[Buneman; Gavril; Walter~\cite{buneman1974,gavril1974,walter1978}]\label{thm:clique-tree-exists}
A graph is chordal if and only if it has a clique tree.
\end{theorem}

The induced subtree property has several equivalent forms, and the paper switches between them.
Let the \emph{weighted clique intersection graph} \(W_G\) be the complete graph on \(\KG\) in
which the pair \(\{K,K'\}\) has weight \(|K\cap K'|\).

\begin{theorem}[Equivalent characterisations; Ibarra \S2, p.\,40:4]\label{thm:ct-equiv}
For a tree \(T\) on \(\KG\) of a connected chordal graph, the following are equivalent:
\begin{enumerate}
\item \(T\) has the induced subtree property;
\item \(T\) has the \emph{clique intersection property}: for any two nodes \(x,y\),
  \(K_x\cap K_y\subseteq K_z\) for every node \(z\) on the \(x\)--\(y\) path of \(T\);
\item \(T\) is a maximum-weight spanning tree of \(W_G\)
  (Bernstein and Goodman~\cite{bernstein1981}).
\end{enumerate}
\end{theorem}

The clique intersection property has an immediate consequence that is the heart of Ibarra's
insertion test:
\begin{equation}\label{eq:path-min}
\text{if } x=z_0,z_1,\dots,z_k=y \text{ is a path in } T, \text{ then }
|K_x\cap K_y|\ \le\ \min_{1\le i\le k} w(z_{i-1},z_i).
\end{equation}
Indeed \(K_x\cap K_y\) is contained in every clique on the path, hence in every intersection
\(K_{z_{i-1}}\cap K_{z_i}\).

\begin{figure}[ht]
\centering
\begin{subfigure}[b]{0.46\textwidth}\centering
\begin{tikzpicture}
\node[cl] (K0) at (0,2.6) {\(K_0=\{3,6\}\)};
\node[clh] (K1) at (0,1.3) {\(K_1=\{0,1,\mathbf{2},3\}\)};
\node[clh] (K2) at (0,0) {\(K_2=\{\mathbf{2},4,5\}\)};
\draw[tedge] (K0)--node[wlab,right=2pt]{1}(K1);
\draw[tedge,line width=1.6pt,draw=orange!80!black] (K1)--node[wlab,right=2pt]{1}(K2);
\end{tikzpicture}
\caption{Clique tree of the running example; the subtree \(T_2\) (cliques containing vertex 2)
is highlighted and is connected.}\label{fig:ct-example}
\end{subfigure}\hfill
\begin{subfigure}[b]{0.5\textwidth}\centering
\begin{tikzpicture}
\begin{scope}
\node[vtx] (s0) at (0,0) {0};
\node[vtx] (s1) at (-1.1,1.0) {1};
\node[vtx] (s2) at (1.1,1.0) {2};
\node[vtx] (s3) at (0,-1.3) {3};
\draw[gedge] (s0)--(s1) (s0)--(s2) (s0)--(s3);
\end{scope}
\begin{scope}[xshift=3.1cm]
\node[cl] (a) at (0,1.1) {\(\{0,1\}\)};
\node[cl] (b) at (-0.95,-0.6) {\(\{0,2\}\)};
\node[cl] (c) at (0.95,-0.6) {\(\{0,3\}\)};
\draw[tedge] (a)--node[wlab,left]{1}(b);
\draw[tedge] (a)--node[wlab,right]{1}(c);
\draw[dedge,draw=black!35] (b)--node[wlab,below]{1}(c);
\end{scope}
\end{tikzpicture}
\caption{A star \(K_{1,3}\) and its intersection graph \(W_G\): every pair of cliques meets in
\(\{0\}\), so \emph{any} two of the three edges form a clique tree.}\label{fig:ct-nonunique}
\end{subfigure}
\caption[Clique trees and their non-uniqueness]{Clique trees. (a) The tree of the running example. (b) Clique trees are not unique.}
\label{fig:clique-trees}
\end{figure}

\paragraph{Minimal separators.} A set \(S\subseteq V\) is a \emph{\(uv\)-separator} if \(u\) and
\(v\) are in the same component of \(G\) but in different components of \(G-S\); it is
\emph{minimal} if no proper subset is a \(uv\)-separator. The edges of a clique tree correspond
exactly to the minimal separators:

\begin{theorem}[Ho and Lee~\cite{holee1989}; Lundquist~\cite{lundquist1990}; Ibarra Theorem~1, p.\,40:5]\label{thm:separators}
Let \(G\) be a connected chordal graph and \(T\) any clique tree of \(G\).
\begin{enumerate}
\item \(S\) is a minimal vertex separator of \(G\) if and only if \(S=K_x\cap K_y\) for some
  edge \(\{x,y\}\) of \(T\).
\item If \(S=K_x\cap K_y\) for an edge \(\{x,y\}\) of \(T\), then \(S\) is a minimal
  \(uv\)-separator for every \(u\in K_x\setminus S\) and \(v\in K_y\setminus S\).
\end{enumerate}
\end{theorem}

In the running example the edge \(K_1K_2\) has \(K_1\cap K_2=\{2\}\): vertex 2 separates
\(\{0,1,3,6\}\) from \(\{4,5\}\).

\section{Disconnected graphs and dummy edges}\label{sec:dummy-edges}
Most texts define clique trees only for connected graphs and treat components separately.
Ibarra instead maintains \emph{one} tree for the whole graph, because an insertion can join two
components and a deletion can split one.

\begin{paperbox}[Ibarra \S2, p.\,40:6 --- dummy edges]
If \(G\) is disconnected with component clique trees \(T_1,\dots,T_k\), join them into a single
tree by arbitrary \emph{dummy edges} \(\{x,y\}\) of weight \(w(x,y)=0\) (the two cliques are in
different components, so they are disjoint). Adding all dummy edges to \(W_G\), the resulting
tree \(T\) is a maximum-weight spanning tree and satisfies the induced subtree and clique
intersection properties. \(T\) is called a clique tree; the \(T_i\) are \emph{strict} clique
trees.
\end{paperbox}

An isolated vertex \(v\) is a maximal clique \(\{v\}\) of its own. \Cref{fig:dummy} shows the
graph used later in \cref{sec:a1-traces}, which has an isolated vertex.

\begin{figure}[ht]
\centering
\begin{tikzpicture}
\begin{scope}
\node[vtx] (f0) at (0,0) {0};
\node[vtx] (f1) at (1.2,1.0) {1};
\node[vtx] (f2) at (1.2,-1.0) {2};
\node[vtx] (f3) at (2.6,1.0) {3};
\node[vtx] (f4) at (2.6,-1.0) {4};
\node[vtx] (f5) at (4.0,0) {5};
\draw[gedge] (f0)--(f1) (f0)--(f2) (f1)--(f2) (f1)--(f3) (f2)--(f3) (f2)--(f4) (f3)--(f4);
\end{scope}
\begin{scope}[xshift=6.6cm]
\node[cl] (c3) at (0,1.3) {\(\{0,1,2\}\)};
\node[cl] (c2) at (0,0) {\(\{1,2,3\}\)};
\node[cl] (c1) at (0,-1.3) {\(\{2,3,4\}\)};
\node[cl] (c0) at (2.7,-1.3) {\(\{5\}\)};
\draw[tedge] (c3)--node[wlab,right]{2}(c2);
\draw[tedge] (c2)--node[wlab,right]{2}(c1);
\draw[dedge] (c1)--node[wlab,below=2pt]{0 (dummy)}(c0);
\end{scope}
\end{tikzpicture}
\caption[A clique tree with a dummy edge]{The ``fan'' graph of \file{fan.txt} (\cref{app:traces}) with an isolated vertex 5,
and the clique tree computed by the program. The singleton clique \(\{5\}\) is attached by a
weight-0 dummy edge.}
\label{fig:dummy}
\end{figure}

\section{Split graphs}\label{sec:split-background}
\begin{definition}[Split graph]
\(G\) is a \emph{split graph} if \(V\) can be partitioned into a clique \(K\) and an
independent set \(I\) (edges between \(K\) and \(I\) are allowed).
\end{definition}
The \emph{complement} \(\bar G\) has the same vertices and exactly the edges that \(G\) lacks.
Földes and Hammer showed that \(G\) is split if and only if both \(G\) and \(\bar G\) are
chordal~\cite{foldes1977}. Hence a dynamic chordal algorithm run on \(G\) and \(\bar G\) gives a
dynamic split algorithm with \(O(n)\) operations; Ibarra (\S8) does much better using a degree
characterisation.

\begin{theorem}[Hammer and Simeone~\cite{hammer1981}; Ibarra Theorem~8, p.\,40:17]\label{thm:hs}
Let \(d_1\ge d_2\ge\dots\ge d_n\) be the degree sequence of \(G\) and
\(j=\max\{i : d_i\ge i-1\}\). Then \(G\) is split if and only if
\begin{equation}\label{eq:hs}
\sum_{i=1}^{j} d_i \;=\; j(j-1)+\sum_{i=j+1}^{n} d_i .
\end{equation}
Moreover, if \eqref{eq:hs} holds, the clique number of \(G\) is \(j\).
\end{theorem}
Since \eqref{eq:hs} depends only on the degree sequence, an edge update changes two degrees by
one, and \cref{ch:split} shows that \(j\) and both sums can then be updated in \(O(1)\) time.

\section{Dynamic trees}\label{sec:dynamic-trees-background}
A \emph{dynamic tree} (or link--cut tree) data structure maintains a forest of rooted trees
with weighted edges under the operations below, each in \(O(\log n)\) time for an \(n\)-node
forest~\cite{sleator1983}:
\begin{itemize}
\item \op{link}\((u,v,c)\): add an edge of cost \(c\) between nodes of different trees;
\item \op{cut}\((u,v)\): delete an edge;
\item \op{evert}\((u)\): make \(u\) the root of its tree;
\item queries on the path from \(u\) to the root, such as the minimum edge cost; with simple
  augmentations also the number of nodes on the path and its \(k\)-th node.
\end{itemize}
The basic idea is to partition the tree into vertex-disjoint paths and store each path in a
balanced binary search tree ordered by depth, so that any root path is the concatenation of
\(O(\log n)\) pieces of stored paths. Sleator and Tarjan give two versions: one with
\emph{amortised} \(O(\log n)\) bounds (later simplified into splay-tree-based link--cut trees)
and one with \emph{worst-case} \(O(\log n)\) bounds, which stores each path in a
\emph{globally biased} binary tree~\cite{bent1985}. Ibarra's A2 cites the dynamic tree of
Sleator and Tarjan and requires worst-case bounds, so the project implements the worst-case
version (ST83 \S5, ``partitioning by size''); it is described in detail in \cref{sec:st83}.

\section{Model of computation}\label{sec:model}
All bounds are in the standard \emph{word RAM} model: memory is an array of words of
\(\Theta(\log n)\) bits, and arithmetic, comparison, memory access and pointer operations on
words take constant time. We follow two conventions throughout:
\begin{itemize}
\item \textbf{Worst case.} A bound \(O(f)\) for an operation means that \emph{every} execution
  of that operation takes \(O(f)\) time. Amortised, expected and randomised bounds are never
  substituted for worst-case ones.
\item \textbf{Initialisation is counted.} Allocating and initialising an array of \(s\) words
  costs \(\Theta(s)\) time. In particular, an \(n\times n\) matrix costs \(\Theta(n^2)\) to
  create, and a characteristic vector of length \(n\) costs \(\Theta(n)\). (Tricks that
  avoid initialisation, such as the ``initialisable array'' of Aho, Hopcroft and
  Ullman~\cite{aho1974} or the sparse sets of Briggs and Torczon~\cite{briggs1993}, still
  \emph{allocate} the memory; we do not use them to hide costs.)
\end{itemize}
The C++ containers used by the implementation have the following worst-case costs, which are
used in all analyses: \code{std::vector} indexing \(O(1)\); \code{std::list} insertion and
erasure at a known iterator \(O(1)\); \code{std::set} search, insertion and deletion
\(O(\log n)\) (red--black tree), and construction from a \emph{sorted} range of length \(k\)
in \(O(k)\) (guaranteed by the C++ standard).
%% ---- end expanded/03_background.tex ----
%% ---- begin expanded/04_ibarra.tex ----
%% =====================================================================
\chapter{Ibarra's Fully Dynamic Algorithm}\label{ch:ibarra}
%% =====================================================================
This chapter presents the algorithm of Ibarra~\cite{ibarra2008} at the level of mathematics and
pseudocode, independent of any data structure. The next chapters implement it. Throughout,
\(G\) is the current chordal graph and \(T\) the current clique tree of \(G\), including dummy
edges (\cref{sec:dummy-edges}).

\begin{keyidea}[The invariant maintained by every implementation]
Before and after every operation, \(T\) is a clique tree of the current graph \(G\):
its nodes are in bijection with the maximal cliques of \(G\), it has the induced subtree
property, and every edge \(\{x,y\}\) stores the correct weight \(w(x,y)=|K_x\cap K_y|\).
Queries only \emph{read} \(T\); updates change \(T\) locally and restore the invariant.
\end{keyidea}

\section{Deletion query}\label{sec:delete-query-math}
\begin{theorem}[Ibarra Theorem~4, p.\,40:6]\label{thm:delete}
Let \(G\) be chordal and \(uv\in E\). Then \(G-uv\) is chordal if and only if exactly one maximal
clique of \(G\) contains both \(u\) and \(v\).
\end{theorem}
\begin{proof}
The proof uses two classical lemmas. (Rose, Tarjan, Lueker~\cite{rose1976}) If \(G\) is chordal
then \(G-uv\) is either chordal or has a chordless 4-cycle. (Grone et al.~\cite{grone1984}) A
graph has no chordless 4-cycle iff for every non-edge \(st\), adding \(st\) creates exactly one
maximal clique containing \(\{s,t\}\).

(\(\Rightarrow\)) Apply the second lemma to the chordal graph \(G-uv\) and the non-edge \(uv\).

(\(\Leftarrow\)) If \(G-uv\) is not chordal it has a chordless 4-cycle, which must be
\((u,s,v,t)\) for some \(s,t\) with \(st\notin E\). Then \(\{u,v,s\}\) and \(\{u,v,t\}\) are cliques
of \(G\) that cannot lie in a common clique (because \(st\notin E\)), so \(uv\) lies in at least
two maximal cliques.
\end{proof}

\begin{algorithm}[ht]
\caption{\op{Delete-Query}\((u,v)\) (Ibarra \S3.1, p.\,40:7)}\label{alg:delete-query}
\begin{algorithmic}[1]
\Require chordal \(G\) with clique tree \(T\); an edge \(uv\in E\)
\State count \(\gets\) number of nodes \(x\) of \(T\) with \(\{u,v\}\subseteq K_x\)
\State \Return \yes\ if count \(=1\), else \no
\end{algorithmic}
\end{algorithm}

\begin{example}
In the running example after inserting the edge \(\{0,4\}\) (\cref{fig:trace-ex} on
\cpageref{fig:trace-ex}) the maximal cliques are \(\{3,6\},\{0,1,2,3\},\{2,4,5\},\{0,2,4\}\). The
edge \(\{0,2\}\) lies in two of them, so \op{Delete-Query}\((0,2)\) answers \no: deleting it
would leave the chordless cycle \(0\!-\!1\!-\!2\!-\!4\!-\!0\). The edge \(\{3,6\}\) lies only in
\(\{3,6\}\), so \op{Delete-Query}\((3,6)\) answers \yes. (Both answers were checked with the
program; see \cref{app:traces}.)
\end{example}

\section{Deletion}\label{sec:delete-math}
Suppose \(uv\) lies in the unique maximal clique \(K_x\), and let \(k=|K_x|\). In \(G-uv\) every
other maximal clique is still a maximal clique, while \(K_x\) is no longer a clique; it is
replaced by the two cliques
\[ K_{xu}=K_x\setminus\{v\},\qquad K_{xv}=K_x\setminus\{u\}, \]
which may or may not be maximal. Partition the tree neighbours \(N_T(x)\) of \(x\) into
\begin{align*}
N_u&=\{y\in N_T(x): u\in K_y\}, & N_v&=\{z\in N_T(x): v\in K_z\}, &
N_0&=\{w\in N_T(x): u,v\notin K_w\}.
\end{align*}
(No neighbour contains both \(u\) and \(v\), because \(K_x\) is the only clique containing
\(uv\); the paper calls \(N_0\) ``\(N_{uv}\)''.) Every neighbour \(y\) has \(w(x,y)\le k-1\), and
\begin{align}\label{eq:delete-maximal}
K_{xu}\text{ is not maximal in }G-uv
\iff \exists y\ne x: K_{xu}\subsetneq K_y
\iff \exists y\in N_u: w(x,y)=k-1 .
\end{align}
The middle step uses the clique intersection property: if \(K_{xu}\subseteq K_y\) for some node
\(y\), then \(K_{xu}=K_x\cap K_y\) is contained in the neighbour of \(x\) on the \(x\)--\(y\) path,
and that neighbour contains \(u\). The same holds for \(K_{xv}\) with \(N_v\).

\begin{algorithm}[ht]
\caption{\op{Delete}\((u,v)\) (Ibarra \S3.2, p.\,40:8)}\label{alg:delete}
\begin{algorithmic}[1]
\Require \(uv\) lies in exactly one maximal clique \(K_x\) (otherwise reject); \(k\gets|K_x|\)
\State Replace node \(x\) by two nodes \(x_1,x_2\) with \(K_{x_1}=K_x\setminus\{v\}\) and
  \(K_{x_2}=K_x\setminus\{u\}\); add the edge \(\{x_1,x_2\}\) with weight \(k-2\)
\For{each neighbour \(y\) of \(x\) (keep the old weight)}
  \If{\(y\in N_u\)} reconnect \(y\) to \(x_1\)
  \ElsIf{\(y\in N_v\)} reconnect \(y\) to \(x_2\)
  \Else\ reconnect \(y\) to \(x_1\) or \(x_2\) (arbitrary)
  \EndIf
\EndFor
\If{some \(y\in N_u\) has \(w(x,y)=k-1\)} contract the edge \(\{x_1,y\}\), keeping \(y\)\EndIf
\If{some \(z\in N_v\) has \(w(x,z)=k-1\)} contract the edge \(\{x_2,z\}\), keeping \(z\)\EndIf
\State delete \(uv\) from \(G\)
\end{algorithmic}
\end{algorithm}

\begin{figure}[ht]
\centering
\begin{tikzpicture}
\begin{scope}
\node[clh,minimum width=2.4cm] (x) at (0,0) {\(K_x\ni u,v\)};
\node[cl] (y) at (-2.2,1.4) {\(K_y\ni u\)};
\node[cl] (z) at (2.2,1.4) {\(K_z\ni v\)};
\node[cl] (w) at (0,-1.6) {\(K_w\not\ni u,v\)};
\draw[tedge] (x)--(y); \draw[tedge] (x)--(z); \draw[tedge] (x)--(w);
\node[font=\small\bfseries] at (0,2.3) {before: \(uv\) only in \(K_x\)};
\end{scope}
\draw[arr,accent] (3.2,0)--node[above,font=\small]{\op{Delete}\((u,v)\)}(4.8,0);
\begin{scope}[xshift=8.4cm]
\node[clg] (x1) at (-1.7,0) {\(K_{x_1}=K_x\!\setminus\!\{v\}\)};
\node[clg] (x2) at (2.1,0) {\(K_{x_2}=K_x\!\setminus\!\{u\}\)};
\node[cl] (y) at (-2.6,1.4) {\(K_y\ni u\)};
\node[cl] (z) at (2.8,1.4) {\(K_z\ni v\)};
\node[cl] (w) at (-1.7,-1.6) {\(K_w\)};
\draw[tedge] (x1)--node[wlab,below=1pt]{\(k-2\)}(x2);
\draw[tedge] (x1)--(y); \draw[tedge] (x2)--(z); \draw[tedge] (x1)--(w);
\node[font=\small\bfseries] at (0.2,2.3) {after (before contractions)};
\end{scope}
\end{tikzpicture}
\caption[The deletion update]{\op{Delete}: the unique clique containing \(uv\) is split into two cliques joined by an
edge of weight \(k-2\). Neighbours containing \(u\) go to \(x_1\), neighbours containing \(v\)
go to \(x_2\). If, e.g., \(w(x,y)=k-1\), then \(K_{x_1}\subseteq K_y\) and \(x_1\) is contracted
into \(y\). (After Ibarra Figs.~3--5.)}
\label{fig:delete}
\end{figure}

After \op{Delete}, the nodes of \(T\) are exactly the maximal cliques of \(G-uv\), the induced
subtree property holds, and all weights are correct; hence \(T\) is a clique tree of \(G-uv\)
(Ibarra p.\,40:8). (Intuitively: \(T_u\) now reaches \(x_1\) through the neighbours in \(N_u\),
\(T_v\) reaches \(x_2\) through \(N_v\), and every other vertex of \(K_x\) lies in both \(x_1\)
and \(x_2\), which are adjacent.) A useful observation: the edge \(\{x_1,x_2\}\) is never
contracted, and its endpoints (which may have become \(y\) and \(z\)) contain \(u\) and \(v\).
In particular, if \(k=2\) the clique \(\{u,v\}\) splits into \(\{u\}\) and \(\{v\}\) joined by a
weight-0 dummy edge: \(uv\) was a cut edge and the component splits.

\section{Insertion query}\label{sec:insert-query-math}
Let \(uv\notin E\). Since no clique contains both \(u\) and \(v\), the subtrees \(T_u\) and
\(T_v\) are disjoint. Let \(x\in T_u\) and \(y\in T_v\) be the \emph{closest} pair of nodes,
i.e.\ the endpoints of the shortest path in \(T\) between the two subtrees. (In a tree this
pair is unique.)

\begin{theorem}[Ibarra Theorem~5, p.\,40:8]\label{thm:insert5}
\(G+uv\) is chordal if and only if there is a clique tree \(T'\) of \(G\) with an edge
\(\{x',y'\}\) such that \(u\in K_{x'}\) and \(v\in K_{y'}\).
\end{theorem}
\begin{proof}[Proof of (\(\Leftarrow\))]
If \(\{x',y'\}\) is a dummy edge, \(u\) and \(v\) lie in different components and \(G+uv\) is
clearly chordal. Otherwise \(I=K_{x'}\cap K_{y'}\) is a \(uv\)-separator by
\cref{thm:separators}. Any cycle of length \(\ge 4\) through the new edge \(uv\) consists of
\(uv\) and a \(u\)--\(v\) path of length \(\ge 3\), which must pass through some \(s\in I\); as
\(s\) is adjacent to both \(u\) and \(v\), the cycle has a chord. (The direction
(\(\Rightarrow\)) follows from the observation at the end of \cref{sec:delete-math}.)
\end{proof}

\begin{theorem}[Ibarra Theorem~6, p.\,40:9]\label{thm:insert6}
Let \(x,y\) be the closest nodes with \(u\in K_x\), \(v\in K_y\), and suppose \(\{x,y\}\notin T\).
A clique tree \(T'\) as in \cref{thm:insert5} exists if and only if the minimum-weight edge
\(e\) on the \(x\)--\(y\) path of \(T\) satisfies \(w(e)=|K_x\cap K_y|\). In that case
\(T'=T-e+\{x,y\}\) is such a tree.
\end{theorem}

The ``if'' direction is easy: by \eqref{eq:path-min} every edge on the path has weight at least
\(|K_x\cap K_y|\), so replacing a minimum edge \(e\) by \(\{x,y\}\) with
\(w(e)=|K_x\cap K_y|\) gives another maximum-weight spanning tree of \(W_G\), i.e.\ another
clique tree. The ``only if'' direction is an exchange argument on maximum spanning trees
(Ibarra p.\,40:9).

\begin{algorithm}[ht]
\caption{\op{Insert-Query}\((u,v)\) (Ibarra \S3.3, p.\,40:10)}\label{alg:insert-query}
\begin{algorithmic}[1]
\Require chordal \(G\) with clique tree \(T\); a non-edge \(uv\)
\State find the closest nodes \(x\in T_u\), \(y\in T_v\)
\If{\(\{x,y\}\) is an edge of \(T\)} \Return \yes \EndIf
\State \(\mu\gets\) minimum weight of an edge on the \(x\)--\(y\) path of \(T\)
\State \Return \yes\ if \(\mu=|K_x\cap K_y|\), and \no\ if \(\mu>|K_x\cap K_y|\)
\end{algorithmic}
\end{algorithm}

\begin{figure}[ht]
\centering
\begin{tikzpicture}
\begin{scope}
\node[clh] (a) at (0,0) {\(\{0,1,2\}\)};
\node[cl] (b) at (2.4,0) {\(\{1,2,3\}\)};
\node[clh] (c) at (4.8,0) {\(\{2,3,4\}\)};
\draw[tedge] (a)--node[wlab,above]{2}(b);
\draw[tedge] (b)--node[wlab,above]{2}(c);
\node[font=\small,align=center] at (2.4,-1.1) {\op{Insert-Query}\((0,4)\): \(x=\{0,1,2\}\), \(y=\{2,3,4\}\)\\
path minimum \(=2\), but \(|K_x\cap K_y|=|\{2\}|=1\) \(\Rightarrow\) \textbf{\no}};
\end{scope}
\begin{scope}[xshift=8.4cm]
\node[vtx] (f0) at (0,0) {0};
\node[vtx] (f1) at (1.0,0.9) {1};
\node[vtx] (f2) at (1.0,-0.9) {2};
\node[vtx] (f3) at (2.2,0.9) {3};
\node[vtx] (f4) at (2.2,-0.9) {4};
\draw[gedge] (f0)--(f1) (f0)--(f2) (f1)--(f2) (f1)--(f3) (f2)--(f3) (f2)--(f4) (f3)--(f4);
\draw[newedge] (f0) to[bend right=50] (f4);
\node[font=\small,align=center] at (1.1,-2.0) {\(0\!-\!1\!-\!3\!-\!4\!-\!0\) would be chordless};
\end{scope}
\end{tikzpicture}
\caption[A rejected insertion query]{Why the test compares the path minimum with the intersection of the \emph{endpoints}.
In the fan graph of \cref{fig:dummy}, every tree edge on the path has weight 2 but the end
cliques share only vertex 2. Inserting \(\{0,4\}\) would create the chordless 4-cycle on the
right, and the program answers \no. This example also shows that \(|K_x\cap K_y|\) cannot be
replaced by a path minimum (\cref{sec:path-min-counterexample}).}
\label{fig:insert-query-no}
\end{figure}

\section{Insertion}\label{sec:insert-math}
Suppose the query answered \yes; after the exchange of \cref{thm:insert6} (if needed) we have
\(u\in K_x\), \(v\in K_y\) and \(\{x,y\}\in T\). Let \(I=K_x\cap K_y\). Then
\[ K_z=I\cup\{u,v\} \]
is a clique of \(G+uv\). It is \emph{maximal}: a vertex adjacent to all of \(K_z\) would connect
\(u\) and \(v\) in \(G-I\), contradicting \cref{thm:separators}; and by \cref{thm:delete} it is
the \emph{only} maximal clique of \(G+uv\) that contains \(uv\). Which old cliques stop being
maximal? A clique \(K_w\) with \(w\ne x,y\) cannot be contained in \(K_z\) (it would be
contained in \(I\cup\{u\}\subseteq K_x\) or \(I\cup\{v\}\subseteq K_y\)), so only \(K_x\) and
\(K_y\) can disappear, and
\begin{equation}\label{eq:insert-maximal}
K_x \text{ is not maximal in } G+uv \iff K_x=I\cup\{u\} \iff |K_x|=|I|+1,
\end{equation}
and symmetrically for \(K_y\). So maximality is decided by comparing three numbers.

\begin{algorithm}[ht]
\caption{\op{Insert}\((u,v)\) (Ibarra \S3.4, p.\,40:11)}\label{alg:insert}
\begin{algorithmic}[1]
\Require a non-edge \(uv\)
\State find the closest nodes \(x\in T_u\), \(y\in T_v\); \(I\gets K_x\cap K_y\)
\If{\(\{x,y\}\notin T\)}
  \State \(e\gets\) a minimum-weight edge on the \(x\)--\(y\) path
  \If{\(w(e)>|I|\)} \textbf{reject} \Comment{\(G+uv\) is not chordal}\EndIf
  \State replace \(e\) by the edge \(\{x,y\}\) with weight \(|I|\) \Comment{\cref{thm:insert6}}
\EndIf
\State replace the edge \(\{x,y\}\) by a new node \(z\), \(K_z=I\cup\{u,v\}\), with edges
  \(\{x,z\},\{z,y\}\) of weight \(|I|+1\)
\If{\(|K_x|=|I|+1\)} contract \(\{x,z\}\), keeping \(z\) \EndIf
\If{\(|K_y|=|I|+1\)} contract \(\{y,z\}\), keeping \(z\) \EndIf
\State insert \(uv\) into \(G\)
\end{algorithmic}
\end{algorithm}

\begin{figure}[ht]
\centering
\begin{tikzpicture}
\begin{scope}
\node[clh] (x) at (0,0) {\(K_x\ni u\)};
\node[cl] (p1) at (1.9,0) {\(\cdot\)};
\node[cl] (p2) at (3.6,0) {\(\cdot\)};
\node[clh] (y) at (5.5,0) {\(K_y\ni v\)};
\draw[tedge] (x)--node[wlab,above]{\(\ge|I|\)}(p1);
\draw[tedge,line width=1.6pt,draw=badred] (p1)--node[wlab,above]{\(e:|I|\)}(p2);
\draw[tedge] (p2)--node[wlab,above]{\(\ge|I|\)}(y);
\draw[tedge,dashed,draw=okgreen!70!black] (x) to[bend right=30] node[wlab,below]{new edge, weight \(|I|\)} (y);
\node[font=\small\bfseries] at (2.75,1.0) {1. exchange (only if \(x,y\) not adjacent)};
\end{scope}
\begin{scope}[yshift=-3.1cm]
\node[clh] (x) at (0,0) {\(K_x\)};
\node[clg] (z) at (2.75,0) {\(K_z=I\cup\{u,v\}\)};
\node[clh] (y) at (5.5,0) {\(K_y\)};
\draw[tedge] (x)--node[wlab,above]{\(|I|+1\)}(z);
\draw[tedge] (z)--node[wlab,above]{\(|I|+1\)}(y);
\node[font=\small\bfseries] at (2.75,0.95) {2. subdivide \(\{x,y\}\) by the new clique \(z\)};
\node[font=\small,align=center] at (2.75,-1.0) {3. if \(|K_x|=|I|+1\), contract \(x\) into \(z\); same for \(y\)};
\end{scope}
\end{tikzpicture}
\caption[The insertion update]{\op{Insert}: the three steps (after Ibarra Figs.~6--8). The tree path between \(T_u\)
and \(T_v\) is shortened to a single edge, which is then subdivided by the unique new maximal
clique \(K_z\).}
\label{fig:insert}
\end{figure}

\begin{example}[Three shapes of insertion]
\begin{enumerate}
\item \emph{No clique disappears.} Running example, \op{Insert}\((0,4)\): \(x=\{0,1,2,3\}\),
  \(y=\{2,4,5\}\) are adjacent, \(I=\{2\}\), and \(K_z=\{0,2,4\}\) is inserted between them with
  edges of weight 2 (\cref{fig:trace-ex}).
\item \emph{Both cliques disappear} (``double subsumption''). In the star with centre 0 and
  leaves 1,2,3, \op{Insert}\((1,3)\): \(x=\{0,1\}\), \(y=\{0,3\}\), \(I=\{0\}\) and
  \(|K_x|=|K_y|=2=|I|+1\), so both are absorbed into \(K_z=\{0,1,3\}\) (\cref{fig:trace-star}).
\item \emph{With an exchange.} In the same star with tree \(\{0,1\}\!-\!\{0,3\}\!-\!\{0,2\}\),
  \op{Insert}\((1,2)\): \(x=\{0,1\}\) and \(y=\{0,2\}\) are not adjacent, the path minimum is
  \(1=|I|\), so one path edge is exchanged for \(\{x,y\}\) before the insertion.
\end{enumerate}
\end{example}

Note that \(\{x,y\}\) is a dummy edge (\(|I|=0\)) exactly when \(uv\) joins two components; then
\(K_z=\{u,v\}\).

\section{Correctness and cost at the level of the paper}
Ibarra proves (pp.\,40:8 and 40:11) that \op{Delete} and \op{Insert} leave a clique tree of the
updated graph. Every step of \cref{alg:delete,alg:insert} touches \(O(1)\) cliques and \(O(n)\)
tree edges: the only non-local work is (i) finding the clique(s) containing \(u\) and \(v\),
(ii) finding the closest pair and the path minimum, and (iii) reconnecting the neighbours of one
or two nodes. The cost of these three steps depends entirely on how \(T\) and the cliques are
represented --- which is what distinguishes Ibarra's implementations.

\section{The implementations described in the paper}\label{sec:paper-implementations}
\begin{paperbox}[Ibarra \S4 (p.\,40:11) --- first implementation, A1]
Represent \(T\) by adjacency lists and each clique by a \emph{characteristic vector} of length
\(n\), so membership is \(O(1)\). Preprocessing computes the maximal cliques and a clique tree
in \(O(m+n)\). All four operations search \(T\) (at most \(n\) nodes) and compare vectors of
length \(n\): \(O(n)\) each. When \(K_x\) or \(K_y\) is absorbed during \op{Insert}, the paper
updates the absorbed clique's vector in place instead of building a new one.
\end{paperbox}

\begin{paperbox}[Ibarra \S5--5.1 (pp.\,40:12--40:13) --- second implementation, A2]
Represent \(T\) as a Sleator--Tarjan dynamic tree (``building the initial dynamic tree requires
\(O(n\log n)\) time''). Keep, for every vertex, a pointer to some clique containing it.
\op{Insert-Query} reroots the tree, binary-searches the path between the two pointed-to
cliques with \(O(\log n)\) ``\(i\)-th node on a path'' operations to find \(x\) and \(y\), and
finds the path minimum: \(O(\log^2 n)\). The intersection \(|K_x\cap K_y|\) is read in
\(O(1)\) from an \(n\times n\) matrix \(W\) with \(W(x,y)=|K_x\cap K_y|\), which is updated in
\(O(n)\) after each \op{Delete} or \op{Insert}. \op{Delete} makes \(O(n)\) link/cut changes:
\(O(n\log n)\). \(W\) is initialised in \(O(n^3)\) by comparing vectors, or in \(O(mn)\) by
inserting the edges of \(G\) one by one into an empty graph in reverse PEO order.
\end{paperbox}

\begin{paperbox}[Ibarra \S6 (pp.\,40:13--40:16) --- faster \op{Delete-Query}, A1\(^+\)/A2\(^+\)]
Maintain, for every vertex \(v\), the set \(\KG(v)\) in a modified version of Yellin's dynamic
set-intersection structure~\cite{yellin1994}; then \op{Delete-Query}\((u,v)\) asks whether
\(|\KG(u)\cap\KG(v)|=1\), in \(O(\sqrt m)\) time. Updates cost \(O(m)\) extra; preprocessing
\(O(mn+n^2)\).
\end{paperbox}

\Cref{tab:paper-bounds} on \cpageref{tab:paper-bounds} lists the resulting bounds. The paper
also remarks (\S4.1) that connected components could be maintained alongside with ET-trees of
Henzinger and King~\cite{henzinger1999}; this remark is not part of the four operations or of the summary table, and
it is not implemented in this project.

\paragraph{From the paper to the code.} The rest of the report follows the paper's structure:
\cref{ch:static} builds the initial clique tree, \cref{ch:a1} implements A1, \cref{ch:a2}
implements A2, \cref{ch:fastdelete} implements \S6 and \cref{ch:split} implements \S8.
%% ---- end expanded/04_ibarra.tex ----
%% ---- begin expanded/05_arch.tex ----
%% =====================================================================
\chapter{Software Architecture}\label{ch:architecture}
%% =====================================================================
This chapter gives a map of the repository \repo{} before the individual algorithms are
discussed. The implementation is written in standard C++17 with no third-party dependencies
and is built with CMake. The algorithmic core (\file{include/} and \file{src/}) consists of
about 4{,}250 lines; the test programs add about 2{,}450 lines and the benchmark programs about
360 lines.

\section{Repository layout}
\begin{lstlisting}[style=shell,caption={Top-level layout of the repository (the build directory
is not part of the repository).},label={lst:tree}]
MTP-3rd-sem/
|-- CMakeLists.txt          build: library mtp_graphs, 4 CLIs, 23 tests, 6 benchmarks
|-- README.md               overview, build instructions, expected output
|-- include/mtp/            public headers (one per module, see Table 4.1)
|-- src/                    implementations + the four command-line programs
|-- tests/                  19 test programs (23 CTest entries incl. 4 CLI runs)
|-- benchmarks/             6 measurement programs + archived CSV results
|-- data/                   example graph, example operation script, non-chordal C4
|-- audit/                  independent Python/NetworkX verification scripts
|-- scripts/                check_markdown_math.py (documentation lint)
|-- docs/                   earlier reports, audit notes, build notes
`-- .github/workflows/      5 continuous-integration workflows
\end{lstlisting}

\section{Modules}
\Cref{tab:modules} lists the modules of the core library \code{mtp_graphs}, the paper section
each one implements, and the published source of its method.

\begin{table}[ht]
\centering\small
\caption[Modules of the core library]{Modules of the core library. ``Header'' is under \file{include/mtp/}, ``Source''
under \file{src/}.}\label{tab:modules}
\begin{tabularx}{\textwidth}{@{}>{\raggedright\arraybackslash}p{3.4cm}>{\raggedright\arraybackslash}p{4.6cm}X@{}}
\toprule
\textbf{Header / source} & \textbf{Main type or function} & \textbf{Role and published source}\\
\midrule
\file{graph} & \code{Graph} & adjacency vectors plus \code{std::set} membership; linear bulk loader\\
\file{mcs} & \code{maximum_cardinality_search}, \code{is_perfect_elimination_ordering} & MCS and the linear PEO test (Tarjan--Yannakakis~\cite{tarjan1984})\\
\file{maximal_cliques} & \code{extract_maximal_cliques} & cliques from a PEO (\cref{lem:cliques,lem:subsumption}); used as a test oracle\\
\file{clique_tree_blair_peyton} & \code{build_clique_tree_blair_peyton} & linear clique tree from MCS (Blair--Peyton~\cite{blair1993}) + dummy edges (Ibarra \S2)\\
\file{clique_graph}, \file{clique_tree} & \code{build_clique_tree} & explicit \(W_G\) + Kruskal; \emph{test baseline only} (library \code{mtp_clique_baseline})\\
\file{sparse_ibarra_a1} & \code{SparseIbarraChordalGraphA1} & \textbf{A1} (Ibarra \S3--4) with sparse clique incidence\\
\file{dynamic_ibarra} & \code{DynamicChordalGraph} & A1 with dense characteristic vectors (Ibarra \S4 literally); used inside A2\\
\file{globally_biased_binary_tree} & \code{GloballyBiasedBinaryTree} & biased search trees, global join/split, linear build (BST85~\cite{bent1985})\\
\file{worst_case_dynamic_tree} & \code{SleatorTarjanWorstCaseForest} & worst-case dynamic trees (ST83~\cite{sleator1983} \S5)\\
\file{dynamic_ibarra_a2} & \code{SleatorTarjanChordalGraphA2} & \textbf{A2} (Ibarra \S5--5.1): dynamic tree + matrix \(W\)\\
\file{indexed_dynamic_ibarra} & \code{IndexedDynamicChordalGraph} & comparison variant: heavy-path index rebuilt after each update (\emph{not} ST83)\\
\file{fast_delete_ibarra} & \code{FastCliqueIntersectionIndex} & modified Yellin structure (Ibarra \S6) for A1\(^+\)/A2\(^+\)\\
\file{dynamic_split_ibarra} & \code{DynamicSplitGraph}, \code{DynamicSplitStatistics} & split graphs (Ibarra \S8, Hammer--Simeone~\cite{hammer1981})\\
\bottomrule
\end{tabularx}
\end{table}

\Cref{fig:modules} shows how the modules depend on each other. The arrows point from a module
to the modules it uses.

\begin{figure}[ht]
\centering
\begin{tikzpicture}[node distance=6mm and 6mm,
  m/.style={boxnode,minimum width=2.9cm,minimum height=8mm},
  mm/.style={m,fill=okgreen!12,draw=okgreen!70!black},
  mt/.style={m,fill=gray!10,draw=gray!60,dashed},
  ex/.style={m,fill=orange!12,draw=orange!70!black},
  dep/.style={-{Stealth[length=2mm]},draw=black!55,thick}]
\node[m] (graph) at (0,0) {\code{Graph}};
\node[m] (mcs) at (0,1.4) {MCS + PEO test};
\node[m] (bp) at (-3.4,2.8) {Blair--Peyton\\clique tree};
\node[mt] (mc) at (3.4,2.8) {maximal cliques\\(oracle)};
\node[m] (fast) at (3.6,4.2) {modified Yellin\\index (\S6)};
\node[mm] (a1) at (-3.8,4.6) {\textbf{A1} sparse\\\code{SparseIbarra...A1}};
\node[m] (dense) at (0,4.6) {dense A1\\\code{DynamicChordalGraph}};
\node[m] (bst) at (-3.8,6.2) {globally biased\\binary tree (BST85)};
\node[m] (st) at (-3.8,7.6) {Sleator--Tarjan\\forest (ST83)};
\node[mm] (a2) at (0,7.6) {\textbf{A2}\\\code{SleatorTarjan...A2}};
\node[mt] (idx) at (3.6,6.4) {indexed variant\\(comparison)};
\node[mm] (split) at (3.6,0) {\textbf{split graphs}\\(\S8)};
\draw[dep] (mcs)--(graph);
\draw[dep] (bp)--(mcs);
\draw[dep] (mc)--(mcs);
\draw[dep] (a1)--(bp);
\draw[dep] (a1) to[bend right=8] (fast);
\draw[dep] (dense)--(bp);
\draw[dep] (dense)--(fast);
\draw[dep] (st)--(bst);
\draw[dep] (a2)--(st);
\draw[dep] (a2)--(dense);
\draw[dep] (idx)--(dense);
\draw[dep] (split)--(graph);
\end{tikzpicture}
\caption[Module dependencies]{Module dependencies. Green: the three user-facing algorithms of the paper.
Dashed grey: modules used only for testing or comparison.}
\label{fig:modules}
\end{figure}

\paragraph{Why are there two A1 classes?} The class \code{DynamicChordalGraph} is a literal
implementation of Ibarra \S4: every clique is a characteristic vector of length \(n\). It was
written first, and A2 is built on top of it because A2's matrix \(W\) (whose rows are indexed by
clique-node identifiers) and its weight-update formulas (\cref{sec:w-updates}) use \(O(1)\)
membership tests \(v\in K_w\) for arbitrary cliques \(w\). Its constructor, however, costs
\(\Theta(cn)\), which can be \(\Theta(n^2)\) and exceeds the \(O(m+n)\) preprocessing bound of
A1. The class \code{SparseIbarraChordalGraphA1} was therefore written to meet that bound
(\cref{ch:a1}); it is the class used by the A1 command-line program and is the A1 of this
report. Both implement exactly the same queries and updates and are tested against each other.

\section{Programs and input formats}\label{sec:programs}
Four command-line programs are built (\cref{tab:clis}).

\begin{table}[ht]
\centering\small
\caption{Command-line programs.}\label{tab:clis}
\begin{tabularx}{\textwidth}{@{}llX@{}}
\toprule
\textbf{Program} & \textbf{Source} & \textbf{Purpose}\\
\midrule
\code{mtp_cli} & \file{src/main.cpp} & static: print the MCS ordering, the maximal cliques and the clique tree, or report ``not chordal'' (exit code 2)\\
\code{mtp_dynamic_cli} & \file{src/dynamic_main.cpp} & run an operation script with \textbf{A1}; flag \code{--fast-delete} enables A1\(^+\)\\
\code{mtp_a2_cli} & \file{src/a2_main.cpp} & run an operation script with \textbf{A2}; \code{--fast-delete} enables A2\(^+\)\\
\code{mtp_indexed_cli} & \file{src/indexed_main.cpp} & run an operation script with the indexed comparison variant\\
\bottomrule
\end{tabularx}
\end{table}

\paragraph{Graph files.} The first line contains \(n\) and \(m\); then \(m\) lines each contain
one edge \(u\ v\) with zero-based vertex numbers. Self-loops, duplicate edges, out-of-range
vertices and trailing text are rejected with an error message. The running example
\file{data/example_7_vertices.txt} is:
\begin{lstlisting}[style=shell]
7 10
0 1
1 3
0 2
2 3
2 4
2 5
4 5
3 6
0 3
1 2
\end{lstlisting}

\paragraph{Operation scripts.} One operation per line: \code{Insert-Query u v},
\code{Insert u v}, \code{Delete-Query u v}, \code{Delete u v} or \code{Show} (print the current
cliques and clique tree). Blank lines and lines starting with \texttt{\#} are ignored. Each
operation prints its answer, \code{yes} or \code{no}; an update that is not allowed is
rejected (prints \code{no}) and leaves the graph unchanged. Example:
\begin{lstlisting}[style=shell]
$ ./build/mtp_dynamic_cli data/example_7_vertices.txt ops.txt
Insert-Query 0 4 => yes
Insert 0 4 => yes
Delete-Query 0 2 => no
\end{lstlisting}

\section{Build system and quality gates}
\code{CMakeLists.txt} builds the library \code{mtp_graphs} from the twelve core source files,
the four programs, the test programs (registered with CTest) and, with the option
\code{-DMTP_BUILD_BENCHMARKS=ON}, six benchmark programs. Every target is compiled with
\code{-Wall -Wextra -Wpedantic}; the Release configuration uses \code{-O3 -DNDEBUG}. Five
GitHub Actions workflows (\cref{sec:ci}) build and test the code automatically, including
builds with AddressSanitizer and UndefinedBehaviorSanitizer.

\section{Design principles of the implementation}
\begin{enumerate}
\item \textbf{Source traceability.} Each function that implements a published step carries a
  comment naming the paper, section and figure (for example
  \code{// Blair & Peyton Fig. 4.2, i = n,...,1}), so that the code can be audited against the
  paper line by line.
\item \textbf{Fail loudly.} Preconditions (chordal input, valid vertex numbers, a valid PEO)
  are checked and violations throw \code{std::invalid_argument}; internal inconsistencies
  (which would indicate a bug) throw \code{std::logic_error}. No operation silently produces a
  wrong clique tree.
\item \textbf{Determinism.} No hashing or randomisation is used in the algorithms: adjacency
  membership uses \code{std::set} (worst-case \(O(\log n)\)) rather than a hash table (whose
  \(O(1)\) bound is only expected). This matters because the paper's bounds are worst-case.
\item \textbf{Stable identifiers.} Clique-tree nodes are slots in an array of capacity \(n+2\)
  with a free list. A chordal graph has at most \(n\) maximal cliques, and an update creates at
  most two new nodes before contracting, so \(n+2\) slots always suffice; node identifiers stay
  valid across updates, which A2 needs because \(W\) is indexed by them.
\item \textbf{Debug-only checkers.} Expensive invariant checkers (e.g.\ the \(O(n^2)\)
  \code{check_invariants} of the dynamic tree) are called only by tests, never by operations.
\end{enumerate}
%% ---- end expanded/05_arch.tex ----
%% ---- begin expanded/06_static.tex ----
%% =====================================================================
\chapter{Static Preprocessing: from a Graph to its Clique Tree}\label{ch:static}
%% =====================================================================
Every dynamic implementation starts from the clique tree of the initial graph. This chapter
describes how the program reads a graph and computes, in \(O(n+m)\) worst-case time, a perfect
elimination ordering, the decision ``chordal or not'', the maximal cliques and a clique tree
with dummy edges. \Cref{fig:pipeline} shows the pipeline.

\begin{figure}[ht]
\centering
\begin{tikzpicture}[node distance=7mm,
  s/.style={boxnode,minimum width=2.2cm,minimum height=11mm,text width=2.1cm}]
\node[s] (in) {edge list\\(file)};
\node[s,right=of in] (g) {\code{Graph}\\\(O(n+m)\)};
\node[s,right=of g] (mcs) {MCS order\\\(O(n+m)\)};
\node[s,right=of mcs] (peo) {PEO test\\\(O(n+m)\)};
\node[s,right=of peo] (bp) {Blair--Peyton\\clique tree\\\(O(n+m)\)};
\node[boxnode,below=8mm of peo,fill=red!8,draw=badred] (rej) {``not chordal''\\(exit code 2)};
\draw[arr] (in)--(g); \draw[arr] (g)--(mcs); \draw[arr] (mcs)--(peo);
\draw[arr] (peo)--node[above=1pt,font=\scriptsize]{PEO}(bp);
\draw[arr] (peo)--node[right,font=\scriptsize]{not a PEO}(rej);
\end{tikzpicture}
\caption[Static preprocessing pipeline]{Static preprocessing pipeline. Each stage is worst-case
linear in the size of the graph.}
\label{fig:pipeline}
\end{figure}

\section{The graph representation}\label{sec:graph-class}
The class \code{Graph} (\file{include/mtp/graph.hpp}, \file{src/graph.cpp}) stores, for every
vertex, two views of its neighbourhood:
\begin{itemize}
\item an adjacency \emph{vector} \code{adjacency_[v]}, used for iteration (MCS, BFS); and
\item a balanced search tree \code{neighbor_sets_[v]} (\code{std::set<int>}), used for the
  membership test \code{has_edge(u,v)} in worst-case \(O(\log n)\) time.
\end{itemize}
The second view exists because A2's \op{Insert-Query} must check that \(uv\) is not already an
edge in \(O(\log^2 n)\) time. A hash set would give \(O(1)\) only in expectation, and an
\(n\times n\) adjacency matrix would cost \(\Theta(n^2)\) to initialise; a balanced tree gives a
deterministic \(O(\log n)\) bound with \(O(n+m)\) space. \Cref{tab:graph-ops} lists the costs.

\begin{table}[ht]
\centering\small
\caption[Operations of the Graph class]{Operations of \code{Graph} and their worst-case costs
(\(d\) = degree).}\label{tab:graph-ops}
\begin{tabularx}{\textwidth}{@{}llX@{}}
\toprule
\textbf{Operation} & \textbf{Cost} & \textbf{Remark}\\
\midrule
\code{Graph::from_edges(n, edges)} & \(O(n+m)\) & bulk loader used by \code{mtp_dynamic_cli} (A1)\\
\code{neighbors(v)} & \(O(1)\) & returns a reference to the vector\\
\code{has_edge(u,v)} & \(O(\log n)\) & red--black tree search\\
\code{add_edge(u,v)} & \(O(\log n + d)\) & \(d\) only if the vector reallocates\\
\code{remove_edge(u,v)} & \(O(\log n + d(u)+d(v))\) \(=O(n)\) & erases from both vectors\\
\bottomrule
\end{tabularx}
\end{table}

\paragraph{Linear-time loading.} Building the sets by \(m\) separate insertions would cost
\(O(m\log n)\), more than the \(O(m+n)\) preprocessing bound of A1. The loader therefore first
produces every adjacency vector in \emph{sorted} order by a bucket sort, and then builds each
\code{std::set} from a sorted range, which the C++ standard guarantees to take linear time.
The trick is to bucket the edge endpoints by the \emph{neighbour} identifier: scanning the
buckets in increasing order appends the neighbours of every vertex in increasing order.
Duplicate edges are detected for free, because a duplicate would be appended twice in a row.

\begin{lstlisting}[style=cpp,firstnumber=8,caption={\file{src/graph.cpp}, lines 8--35: the linear-time bulk loader \code{Graph::from_edges}.},label={lst:from-edges}]
Graph Graph::from_edges(
        int n,const std::vector<std::pair<int,int>>& edges) {
    Graph graph(n);
    // Bucket by neighbor ID, not by source. Iterating buckets
    // in increasing order supplies each source's sorted neighbors
    // without an O(m log n) comparison sort.
    std::vector<std::vector<int>> by_neighbor(static_cast<std::size_t>(n));
    for(const auto& [u,v]:edges) {
        graph.validate_vertex(u);
        graph.validate_vertex(v);
        if(u==v)throw std::invalid_argument("self-loops are not allowed");
        by_neighbor[v].push_back(u);
        by_neighbor[u].push_back(v);
    }
    for(int v=0;v<n;++v) {
        for(int u:by_neighbor[v]) {
            auto& adj=graph.adjacency_[u];
            if(!adj.empty() && adj.back()==v)
                throw std::invalid_argument("duplicate input edge");
            adj.push_back(v);
        }
    }
    for(int v=0;v<n;++v)
        graph.neighbor_sets_[v]=std::set<int>(
            graph.adjacency_[v].begin(),graph.adjacency_[v].end());
    graph.edge_count_=edges.size();
    return graph;
}
\end{lstlisting}

\section{Maximum cardinality search}\label{sec:mcs-code}
\Cref{alg:mcs} is MCS with the standard bucket implementation. The vertices are kept in buckets
indexed by their current weight (number of numbered neighbours); a pointer \(\mathit{max}\)
tracks the highest non-empty bucket.

\begin{algorithm}[ht]
\caption{Maximum cardinality search (as implemented in \file{src/mcs.cpp})}\label{alg:mcs}
\begin{algorithmic}[1]
\Require graph \(G\) on vertices \(0,\dots,n-1\)
\Ensure array \(\mathit{ordering}[0..n-1]\); for chordal \(G\) it is a PEO
\State \(\mathit{weight}[v]\gets 0\) for all \(v\); put every vertex in bucket 0; \(\mathit{max}\gets 0\)
\For{\(\mathit{position}=n-1\) \textbf{down to} \(0\)}
  \While{\(\mathit{max}>0\) and bucket[\(\mathit{max}\)] is empty} \(\mathit{max}\gets\mathit{max}-1\) \EndWhile
  \State remove the first vertex \(c\) of bucket[\(\mathit{max}\)]
  \State \(\mathit{ordering}[\mathit{position}]\gets c\); mark \(c\) numbered
  \For{each unnumbered neighbour \(w\) of \(c\)}
    \State move \(w\) from bucket[\(\mathit{weight}[w]\)] to the front of bucket[\(\mathit{weight}[w]+1\)]
    \State \(\mathit{weight}[w]\gets\mathit{weight}[w]+1\); \(\mathit{max}\gets\max(\mathit{max},\mathit{weight}[w])\)
  \EndFor
\EndFor
\State \Return \(\mathit{ordering}\)
\end{algorithmic}
\end{algorithm}

Because positions are filled from \(n-1\) down to \(0\), the first vertex chosen by MCS ends
up last in \code{ordering}, so \code{ordering} is directly the PEO \(v_1,\dots,v_n\) of
\cref{thm:mcs} (with \(v_{i}={}\)\code{ordering}\([i-1]\)).

\begin{lstlisting}[style=cpp,firstnumber=10,caption={\file{src/mcs.cpp}, lines 10--52: maximum cardinality search with bucket lists.},label={lst:mcs}]
std::vector<int> maximum_cardinality_search(const Graph& graph) {
    const int n = graph.vertex_count();
    std::vector<int> ordering(static_cast<std::size_t>(n));
    if (n == 0) return ordering;

    std::vector<int> weight(static_cast<std::size_t>(n), 0);
    std::vector<bool> numbered(static_cast<std::size_t>(n), false);
    std::vector<std::list<int>> buckets(static_cast<std::size_t>(n));
    std::vector<std::list<int>::iterator> locations(static_cast<std::size_t>(n));

    // Use push_front to preserve deterministic bucket tie-breaking.
    // deterministic tie-breaking behavior for the supplied sample graph.
    for (int v = 0; v < n; ++v) {
        buckets[0].push_front(v);
        locations[static_cast<std::size_t>(v)] = buckets[0].begin();
    }

    int max_weight = 0;
    for (int position = n - 1; position >= 0; --position) {
        while (max_weight > 0 && buckets[static_cast<std::size_t>(max_weight)].empty()) {
            --max_weight;
        }
        auto& bucket = buckets[static_cast<std::size_t>(max_weight)];
        // Invariant: at least one unnumbered vertex remains, in some bucket.
        const int chosen = bucket.front();
        bucket.pop_front();
        ordering[static_cast<std::size_t>(position)] = chosen;
        numbered[static_cast<std::size_t>(chosen)] = true;

        for (const int neighbor : graph.neighbors(chosen)) {
            if (numbered[static_cast<std::size_t>(neighbor)]) continue;
            const int old_weight = weight[static_cast<std::size_t>(neighbor)];
            buckets[static_cast<std::size_t>(old_weight)].erase(
                locations[static_cast<std::size_t>(neighbor)]);
            const int new_weight = ++weight[static_cast<std::size_t>(neighbor)];
            buckets[static_cast<std::size_t>(new_weight)].push_front(neighbor);
            locations[static_cast<std::size_t>(neighbor)] =
                buckets[static_cast<std::size_t>(new_weight)].begin();
            if (new_weight > max_weight) max_weight = new_weight;
        }
    }
    return ordering;
}
\end{lstlisting}

\paragraph{Data structure details.} Each bucket is a \code{std::list<int>}, and
\code{locations[v]} stores an iterator to \(v\)'s list cell, so moving a vertex between buckets
is \(O(1)\). New entries are pushed to the \emph{front}; this fixes the tie-breaking rule and
makes the output deterministic (the program always prints the same ordering for the same input
file).

\begin{proposition}\label{prop:mcs-linear}
\code{maximum_cardinality_search} runs in \(O(n+m)\) worst-case time.
\end{proposition}
\begin{proof}
Each vertex is selected once. Each edge \(\{c,w\}\) is examined twice (once from each end) and
causes at most one bucket move, \(O(1)\) each. The variable \(\mathit{max}\) increases by at
most one per bucket move, so it increases at most \(m\) times in total, and hence the
\code{while} loop decrements it at most \(m\) times over the whole run. Total \(O(n+m)\).
\end{proof}

\paragraph{Trace on the running example.} \Cref{tab:mcs-trace} shows each step for the
running example, loaded as \code{mtp_cli} loads it (adjacency vectors in input order). The
result is the ordering \((4,5,1,2,0,3,6)\) printed by the program.

\begin{table}[ht]
\centering\small
\caption[MCS trace on the running example]{MCS on the running example. ``Weights after'' lists
the weights of the unnumbered vertices after the step.}\label{tab:mcs-trace}
\begin{tabular}{@{}cccl@{}}
\toprule
\textbf{Step} & \textbf{Position} & \textbf{Chosen} & \textbf{Weights after (vertex:weight)}\\
\midrule
1 & 6 & 6 & 0:0, 1:0, 2:0, 3:1, 4:0, 5:0\\
2 & 5 & 3 & 0:1, 1:1, 2:1, 4:0, 5:0\\
3 & 4 & 0 & 1:2, 2:2, 4:0, 5:0\\
4 & 3 & 2 & 1:3, 4:1, 5:1\\
5 & 2 & 1 & 4:1, 5:1\\
6 & 1 & 5 & 4:2\\
7 & 0 & 4 & ---\\
\bottomrule
\end{tabular}
\end{table}

\section{Testing for a perfect elimination ordering}\label{sec:peo-code}
\Cref{lem:peo-test} reduces the PEO test to checking \(L(u)\setminus\{p(u)\}\subseteq
N(p(u))\) for every \(u\). Checking each \(u\) separately against a scan of \(N(p(u))\) could
cost \(\Theta(n\cdot m)\) when one vertex is the parent of many others. The implementation
avoids this by grouping the vertices by their parent: for each parent \(p\) it marks
\(N(p)\) \emph{once} with the stamp \(p\), and then checks all children of \(p\) against the
marks. Each adjacency list is therefore scanned \(O(1)\) times.

\begin{lstlisting}[style=cpp,firstnumber=54,caption={\file{src/mcs.cpp}, lines 54--94: the linear-time PEO test.},label={lst:peo}]
bool is_perfect_elimination_ordering(const Graph& graph,
                                     const std::vector<int>& ordering) {
    const int n = graph.vertex_count();
    if (ordering.size() != static_cast<std::size_t>(n)) return false;

    std::vector<int> position(n, -1), parent(n, -1);
    for (int i = 0; i < n; ++i) {
        const int u = ordering[i];
        if (u < 0 || u >= n || position[u] != -1) return false;
        position[u] = i;
    }

    // Find the earliest later neighbor p(u).
    std::vector<std::vector<int>> children(n);
    for (int u = 0; u < n; ++u) {
        for (int v : graph.neighbors(u)) {
            if (position[v] > position[u] &&
                (parent[u] == -1 || position[v] < position[parent[u]])) {
                parent[u] = v;
            }
        }
        if (parent[u] != -1) children[parent[u]].push_back(u);
    }

    // Group vertices by their earliest later neighbor p. Mark N(p) once,
    // then check every other later neighbor for all children of p.
    // Each adjacency list is scanned O(1) times => deterministic O(n+m).
    std::vector<int> marked(n, -1);
    for (int p = 0; p < n; ++p) {
        if (children[p].empty()) continue;
        for (int v : graph.neighbors(p)) marked[v] = p;
        for (int u : children[p]) {
            for (int v : graph.neighbors(u)) {
                if (v != p && position[v] > position[u] && marked[v] != p) {
                    return false;
                }
            }
        }
    }
    return true;
}
\end{lstlisting}

\section{Maximal cliques (test oracle)}\label{sec:cliques-code}
\code{extract_maximal_cliques} implements \cref{lem:cliques,lem:subsumption} directly: it
computes the later neighbours and the parent of every vertex, marks a parent \(p\) as
non-maximal when some child \(u\) has \(|L(u)|-1=|L(p)|\), and outputs \(C(v)\) for every
remaining \(v\). It runs in \(O(n+m)\) time. The dynamic algorithms do not call it --- the
Blair--Peyton construction below produces the cliques and the tree together --- but every test
program uses it as an \emph{independent oracle} for the correct set of maximal cliques.

\section{The Blair--Peyton clique tree}\label{sec:bp}
Blair and Peyton~\cite[\S4.2]{blair1993} observe that MCS visits the maximal cliques of a
chordal graph one after another, and that the order in which it does so is the order in which
Prim's algorithm would build a maximum-weight spanning tree of \(W_G\). Consequently a clique
tree can be produced \emph{during} MCS, without ever building \(W_G\). The implementation
(\file{src/clique_tree_blair_peyton.cpp}; the code comments refer to the algorithm as
Blair--Peyton's ``Fig.~4.2'') replays the MCS order and maintains:
\begin{itemize}
\item \(\mathit{card}\): the number of already-numbered neighbours of the current vertex;
\item \(\mathit{prev}\): the value of \(\mathit{card}\) for the previous vertex;
\item \(\mathit{clique}[v]\): the index of the clique to which numbered vertex \(v\) was added.
\end{itemize}
A new clique starts exactly when \(\mathit{card}\le\mathit{prev}\) (the numbered neighbourhood
stopped growing, so the previous clique is complete). The new clique initially consists of the
current vertex's numbered neighbours --- these form the separator with the clique already built
that contains the most recently numbered of them --- and the tree edge to that clique gets
weight \(\mathit{card}\). If \(\mathit{card}=0\) the vertex starts a new connected component,
and, following Ibarra \S2, the new clique is attached to the first component's clique by a
weight-0 dummy edge.

\begin{algorithm}[ht]
\caption{Clique tree from MCS (after Blair and Peyton~\cite[\S4.2]{blair1993}) with dummy edges
(Ibarra \S2)}\label{alg:bp}
\begin{algorithmic}[1]
\Require chordal \(G\) and its MCS-derived PEO \(\sigma=(v_1,\dots,v_n)\)
\Ensure list of cliques and list of weighted tree edges
\State \(\mathit{prev}\gets 0\); \(\mathit{first}\gets\) none
\For{\(i=n\) \textbf{down to} \(1\)} \Comment{MCS visit order}
  \State \(\mathit{Lab}\gets N(v_i)\cap\{v_{i+1},\dots,v_n\}\); \(\mathit{card}\gets|\mathit{Lab}|\)
  \If{\(\mathit{card}\le\mathit{prev}\)} \Comment{start a new clique \(K_s\)}
    \State \(K_s\gets\mathit{Lab}\)
    \If{\(\mathit{card}>0\)}
      \State \(v_k\gets\) the vertex of \(\mathit{Lab}\) with the smallest index \(k\)
      \State add tree edge \(\{K_s,K_{\mathit{clique}[v_k]}\}\) with weight \(\mathit{card}\)
    \ElsIf{\(\mathit{first}\) is defined}
      \State add dummy edge \(\{K_s,K_{\mathit{first}}\}\) with weight 0
    \Else\ \(\mathit{first}\gets s\)
    \EndIf
  \EndIf
  \State \(K_s\gets K_s\cup\{v_i\}\) for the current (last) clique \(s\); \(\mathit{clique}[v_i]\gets s\)
  \State \(\mathit{prev}\gets\mathit{card}\)
\EndFor
\end{algorithmic}
\end{algorithm}

\begin{lstlisting}[style=cpp,firstnumber=9,caption={\file{src/clique_tree_blair_peyton.cpp}, lines 9--67: Blair--Peyton clique tree with dummy edges.},label={lst:bp}]
CliqueTree build_clique_tree_blair_peyton(
        const Graph& graph, const std::vector<int>& mcs_peo) {
    if (!is_perfect_elimination_ordering(graph, mcs_peo)) {
        throw std::invalid_argument("Blair-Peyton requires an MCS-derived valid PEO");
    }
    const int n = graph.vertex_count();
    CliqueTree tree;
    if (n == 0) return tree;

    std::vector<int> position(n), clique_of_vertex(n, -1);
    for (int i = 0; i < n; ++i) position[mcs_peo[i]] = i;
    int previous_card = 0;
    int first_component_clique = -1;

    // Blair & Peyton Fig. 4.2, i = n,...,1. MCS vertex selection
    // is performed by the caller before passing its PEO.
    for (int index = n - 1; index >= 0; --index) {
        const int u = mcs_peo[index];
        int current_card = 0;
        int minimum_labeled_neighbor = -1;
        for (int v : graph.neighbors(u)) {
            if (position[v] <= index) continue; // previously numbered
            ++current_card;
            if (minimum_labeled_neighbor < 0 ||
                position[v] < position[minimum_labeled_neighbor]) {
                minimum_labeled_neighbor = v;
            }
        }
        // Fig. 4.2: begin a new clique iff new_card <= prev_card.
        if (current_card <= previous_card) {
            const int new_clique = static_cast<int>(tree.cliques.size());
            tree.cliques.emplace_back();
            // K_s <- adj(v_i) intersect L_{i+1}.
            for (int v : graph.neighbors(u)) {
                if (position[v] > index) tree.cliques.back().push_back(v);
            }
            if (current_card != 0) {
                // k <- min{j : v_j in K_s}; p <- clique(v_k).
                const int parent_clique = clique_of_vertex[minimum_labeled_neighbor];
                if (parent_clique < 0 || parent_clique == new_clique)
                    throw std::logic_error("Blair-Peyton parent clique unavailable");
                tree.edges.push_back({new_clique, parent_clique, current_card, false});
            } else {
                // Ibarra Section 2: dummy edge joining components, weight 0.
                if (first_component_clique < 0)
                    first_component_clique = new_clique;
                else
                    tree.edges.push_back({new_clique, first_component_clique, 0, true});
            }
        }
        const int s = static_cast<int>(tree.cliques.size()) - 1;
        clique_of_vertex[u] = s;               // clique(v_i) <- s
        tree.cliques.back().push_back(u);       // K_s <- K_s union {v_i}
        previous_card = current_card;
    }
    if (tree.edges.size() + 1 != tree.cliques.size())
        throw std::logic_error("Blair-Peyton invalid clique-tree edge count");
    return tree;
}
\end{lstlisting}

The final check \code{tree.edges.size() + 1 != tree.cliques.size()} guards the tree property:
\(c\) cliques must be joined by exactly \(c-1\) edges.

\paragraph{Trace on the running example.} \Cref{tab:bp-trace} replays \cref{alg:bp} on the
PEO \((4,5,1,2,0,3,6)\). It reproduces exactly the output of \code{mtp_cli} shown in
\cref{fig:overview}: cliques \(K_0=\{6,3\}\), \(K_1=\{3,0,2,1\}\), \(K_2=\{2,5,4\}\) and edges
\(K_1K_0\) and \(K_2K_1\), both of weight 1.

\begin{table}[ht]
\centering\small
\caption[Blair--Peyton trace on the running example]{\Cref{alg:bp} on the running example.
``Labelled'' are the numbered neighbours of the current vertex.}\label{tab:bp-trace}
\begin{tabular}{@{}cccccl@{}}
\toprule
\(i\) & \(v_i\) & \textbf{Labelled} & \(\mathit{card}\) & \(\mathit{prev}\) & \textbf{Action}\\
\midrule
7 & 6 & \(\emptyset\) & 0 & 0 & new clique \(K_0=\{6\}\) (first component)\\
6 & 3 & \(\{6\}\) & 1 & 0 & add 3: \(K_0=\{6,3\}\)\\
5 & 0 & \(\{3\}\) & 1 & 1 & new \(K_1=\{3\}\), edge \(K_1K_0\) weight 1; add 0\\
4 & 2 & \(\{0,3\}\) & 2 & 1 & add 2: \(K_1=\{3,0,2\}\)\\
3 & 1 & \(\{0,3,2\}\) & 3 & 2 & add 1: \(K_1=\{3,0,2,1\}\)\\
2 & 5 & \(\{2\}\) & 1 & 3 & new \(K_2=\{2\}\), edge \(K_2K_1\) weight 1; add 5\\
1 & 4 & \(\{2,5\}\) & 2 & 1 & add 4: \(K_2=\{2,5,4\}\)\\
\bottomrule
\end{tabular}
\end{table}

\section{The quadratic baseline}\label{sec:baseline}
The repository also contains the ``textbook'' construction: build \(W_G\) explicitly
(\file{src/clique_graph.cpp}, all \(\binom{c}{2}\) pairwise intersections of sorted cliques)
and run Kruskal's maximum spanning tree algorithm with a disjoint-set union
(\file{src/clique_tree.cpp}), adding dummy edges between components. Its running time is
\(O(c^2 n + c^2\log c)\), which can be \(\Theta(n^3)\). It is compiled only into the test
library \code{mtp_clique_baseline}, where it serves as an independent cross-check of
\cref{thm:ct-equiv}(3): the tests verify that the Blair--Peyton tree has the same total weight
as a maximum-weight spanning tree computed independently with Kruskal's algorithm, that both
satisfy the induced subtree property, and (for graphs with at most five vertices) that the
Kruskal result agrees with an exhaustive enumeration of all spanning trees via Prüfer
sequences.

\section{Complexity summary}
\begin{table}[ht]
\centering\small
\caption[Worst-case cost of the static preprocessing]{Worst-case cost of static
preprocessing.}\label{tab:static-cost}
\begin{tabular}{@{}lll@{}}
\toprule
\textbf{Stage} & \textbf{Function} & \textbf{Worst-case time}\\
\midrule
load edge list & \code{Graph::from_edges} & \(O(n+m)\)\\
ordering & \code{maximum_cardinality_search} & \(O(n+m)\)\\
chordality test & \code{is_perfect_elimination_ordering} & \(O(n+m)\)\\
clique tree & \code{build_clique_tree_blair_peyton} & \(O(n+m)\); output size \(\sum|K|\le n+m\)\\
\midrule
total & & \(O(n+m)\)\\
\bottomrule
\end{tabular}
\end{table}

The above total refers to the bulk-loaded A1 preprocessing route.
The output size bound follows from \cref{lem:cliques}: the cliques are written vertex by vertex,
and their total size is at most \(n+m\). (The static program \code{mtp_cli} loads its input with
\(m\) calls of \code{add_edge}, which costs \(O(m\log n)\); this only affects that
demonstration program. Only the A1 dynamic CLI (\code{mtp_dynamic_cli}) uses the linear bulk loader.
The A2 and indexed CLIs (\code{mtp_a2_cli} and \code{mtp_indexed_cli})
ingest edges with repeated \code{Graph::add_edge} calls, requiring up to
\(O(m\log n)\) time for ingestion alone; this does not improve their
already larger initial preprocessing cost.)
%% ---- end expanded/06_static.tex ----
%% ---- begin expanded/07_a1.tex ----
%% =====================================================================
\chapter{Implementation of A1}\label{ch:a1}
%% =====================================================================
This chapter describes \code{SparseIbarraChordalGraphA1}
(\file{include/mtp/sparse_ibarra_a1.hpp}, \file{src/sparse_ibarra_a1.cpp}), the
implementation of Ibarra's first algorithm used by the program \code{mtp_dynamic_cli}. It
realises \cref{alg:delete-query,alg:delete,alg:insert-query,alg:insert} exactly, and attains
the bounds of the first row of \cref{tab:paper-bounds}: \(O(m+n)\) preprocessing and \(O(n)\)
worst-case time for each of the four operations.

\section{Why not characteristic vectors?}\label{sec:why-sparse}
Ibarra \S4 represents each maximal clique by a characteristic vector of length \(n\). Each
operation then costs \(O(n)\), as claimed. The preprocessing, however, must \emph{create} one
such vector per maximal clique, which costs \(\Theta(n)\) per clique and \(\Theta(cn)\) in
total. Since \(c\) can be close to \(n\) while \(m\) is small, this is \(\Theta(n^2)\) for
graphs with \(m=O(n)\):

\begin{example}
For the path \(P_n\) (or a star \(K_{1,n-1}\)) there are \(c=n-1\) maximal cliques (the edges),
\(m=n-1\), and the characteristic vectors occupy \((n-1)\cdot n\) entries --- quadratic in the
input size \(n+m=2n-1\). For an empty graph on \(n\) vertices there are \(n\) singleton cliques
and \(n^2\) vector entries for an input of size \(n\).
\end{example}

In our model initialisation is counted (\cref{sec:model}), so the literal \S4 representation
does not meet the \(O(m+n)\) preprocessing bound of \cref{tab:paper-bounds}. The repository
keeps the literal version as \code{DynamicChordalGraph} (\cref{sec:dense-a1}) and uses a
\emph{sparse} representation for A1, whose total size is \(\sum_{K}|K|\le n+m\)
(\cref{lem:cliques}).

\begin{caution}[Representation change (disclosed)]
The sparse incidence lists replace the paper's characteristic vectors. The clique tree, the
query criteria (\cref{thm:delete,thm:insert6}) and the update steps are exactly Ibarra's; only
the way clique membership is stored and tested is different. Membership tests that the paper
performs in \(O(1)\) on a vector are performed here in batches with a temporary mark array,
which keeps every operation within \(O(n)\). The paper does not describe this layout; it is a
standard data-structure refinement (inverse lists with cross pointers, cf.\
Briggs--Torczon~\cite{briggs1993}).
\end{caution}

\section{Data structures}\label{sec:a1-ds}
\Cref{fig:a1-ds} illustrates the representation and \cref{lst:a1-members} shows the member
variables.

\begin{figure}[ht]
\centering
\begin{tikzpicture}[
  cell/.style={draw=black!60,minimum width=6mm,minimum height=5.5mm,font=\scriptsize,inner sep=0pt},
  lab/.style={font=\small\bfseries,align=center},
  ptr/.style={-{Stealth[length=1.8mm]},draw=accent,thick}]
% clique nodes
\node[lab] at (0.9,3.9) {clique nodes\\\code{nodes_[x].members}};
\node[font=\small] at (-0.4,2.8) {\(x=1\)};
\foreach \v [count=\i] in {3,0,2,1}{ \node[cell,fill=teal!8] (a\i) at (0+0.6*\i,2.8) {\v}; }
\node[font=\small] at (-0.4,1.8) {\(x=2\)};
\foreach \v [count=\i] in {2,5,4}{ \node[cell,fill=teal!8] (b\i) at (0+0.6*\i,1.8) {\v}; }
\node[font=\small] at (-0.4,0.8) {\(x=0\)};
\foreach \v [count=\i] in {6,3}{ \node[cell,fill=teal!8] (c\i) at (0+0.6*\i,0.8) {\v}; }
% vertex lists
\node[lab] at (8.9,3.9) {per-vertex lists \code{nodes_of_[v]}\\(the node sets \(T_v\))};
\node[font=\small] at (7.9,2.8) {\(v=2\)};
\node[cell,fill=accentlight] (v2a) at (8.7,2.8) {1};
\node[cell,fill=accentlight] (v2b) at (9.3,2.8) {2};
\node[font=\small] at (7.9,1.8) {\(v=3\)};
\node[cell,fill=accentlight] (v3a) at (8.7,1.8) {0};
\node[cell,fill=accentlight] (v3b) at (9.3,1.8) {1};
\node[font=\small] at (7.9,0.8) {\(v=6\)};
\node[cell,fill=accentlight] (v6a) at (8.7,0.8) {0};
\draw[ptr] (a3.north) to[bend left=14] (v2a.north);
\node[font=\scriptsize,text=accent,fill=white,inner sep=1pt] at (4.9,3.35) {inverse iterators};
\draw[ptr] (b1.east) to[bend right=6] (v2b.west);
\draw[ptr] (a1.south) to[bend right=10] (v3b.south west);
\draw[ptr] (c2.east) to[bend right=8] (v3a.south west);
\draw[ptr] (c1.south) to[bend right=10] (v6a.south);
% tree edges
\node[font=\small,align=left,anchor=north west,text width=11cm] at (-0.8,0.0)
 {\textbf{tree edges} \code{edges_[e]}: \(e_0=\{1,0\}\), \(w=1\); \(e_1=\{2,1\}\), \(w=1\).
  Each node keeps a list \code{incident} of edge ids, and each edge stores iterators into
  both lists, so \code{link} and \code{cut} are \(O(1)\).};
\end{tikzpicture}
\caption[Sparse clique representation of A1]{The sparse representation of A1 for the running
example (only some inverse pointers are drawn). Each clique stores its members in a vector;
each vertex \(v\) stores the list of clique nodes containing it, which is exactly the node set
of the subtree \(T_v\). The \(i\)-th member of a clique keeps an iterator to the clique's cell
in that vertex's list, so a membership can be removed in \(O(1)\).}
\label{fig:a1-ds}
\end{figure}

\begin{lstlisting}[style=cpp,firstnumber=40,caption={\file{include/mtp/sparse_ibarra_a1.hpp}, lines 40--63: node, edge and member variables of A1.},label={lst:a1-members}]
    struct Node {
        bool active = false;
        std::vector<int> members;
        std::vector<InverseIterator> inverse_iterators;
        std::list<int> incident;
    };
    struct Edge {
        bool active = false;
        int a=-1, b=-1, weight=0;
        std::list<int>::iterator in_a, in_b;
    };
    struct EndpointPair {
        int x=-1, y=-1;
        std::vector<int> path_edges; // edge IDs, from y backwards towards x
    };

    Graph graph_;
    std::vector<Node> nodes_;
    std::vector<Edge> edges_;
    std::vector<std::list<int>> nodes_of_; // inverse list T_v
    std::vector<int> free_nodes_, free_edges_;
    int active_count_=0;
    std::size_t membership_count_=0;
    std::unique_ptr<FastCliqueIntersectionIndex> fast_;
\end{lstlisting}

\noindent The representation maintains the following invariants after every operation.
\begin{enumerate}[label=(I\arabic*)]
\item The active nodes are in bijection with the maximal cliques of the current graph, and
  \code{members} lists the vertices of the clique.
\item \code{nodes_of_[v]} lists exactly the active nodes whose clique contains \(v\); for the
  \(i\)-th member \(v\) of node \(x\), \code{inverse_iterators[i]} points to the cell of \(x\)
  in \code{nodes_of_[v]}.
\item The active edges form a spanning tree on the active nodes; each edge is listed in the
  \code{incident} lists of both of its endpoints, at the stored iterators.
\item Each active edge \(\{a,b\}\) stores \(\mathit{weight}=|K_a\cap K_b|\); weight 0 marks a
  dummy edge.
\item \code{membership_count_} \(=\sum_{x\text{ active}}|K_x|\).
\end{enumerate}
Inactive node and edge slots are kept in the free lists \code{free_nodes_} and
\code{free_edges_}. There are \(n+2\) node slots (\cref{ch:architecture}) and at most \(n+1\)
edge slots.

\section{Primitive operations}
The updates are composed of six primitives (lines 29--122 of \file{src/sparse_ibarra_a1.cpp}):
\begin{itemize}
\item \code{replace_members(x, vs)} removes \(x\) from the vertex lists of its old members
  (using the stored iterators) and inserts it into the lists of the new members:
  \(O(|K_x^{\text{old}}|+|K_x^{\text{new}}|)\).
\item \code{allocate_node(vs)} takes a free slot and calls \code{replace_members};
  \code{retire_node(x)} empties a node with no incident edges and frees its slot.
\item \code{link(a,b,w)} and \code{cut(e)} add and remove a tree edge in \(O(1)\) using the
  stored list iterators.
\item \code{contract(keep, drop)} moves every edge of \code{drop} (except the one to
  \code{keep}) to \code{keep}, then retires \code{drop}:
  \(O(\deg_T(\mathit{drop})+|K_{\mathit{drop}}|)\).
\end{itemize}

\begin{lstlisting}[style=cpp,firstnumber=109,caption={\file{src/sparse_ibarra_a1.cpp}, lines 109--122: contracting a tree edge.},label={lst:a1-contract}]
void SparseIbarraChordalGraphA1::contract(int keep,int drop) {
    if(keep==drop)throw std::logic_error("A1 self contraction");
    std::vector<std::pair<int,int>> neighbors;
    for(int edge_id:nodes_[drop].incident) {
        const auto& e=edges_[edge_id];
        neighbors.push_back({e.a==drop?e.b:e.a,edge_id});
    }
    for(const auto& [other,id]:neighbors) {
        const int weight=edges_[id].weight;
        cut(id);
        if(other!=keep)link(keep,other,weight);
    }
    retire_node(drop);
}
\end{lstlisting}

\section{Construction}
The constructor runs the static pipeline of \cref{ch:static} and stores its result:

\begin{lstlisting}[style=cpp,firstnumber=12,caption={\file{src/sparse_ibarra_a1.cpp}, lines 12--21: the A1 constructor.},label={lst:a1-ctor}]
SparseIbarraChordalGraphA1::SparseIbarraChordalGraphA1(Graph graph)
    : graph_(std::move(graph)), nodes_of_(graph_.vertex_count()) {
    const auto peo=maximum_cardinality_search(graph_);
    if(!is_perfect_elimination_ordering(graph_,peo))
        throw std::invalid_argument("initial A1 graph must be chordal");
    const auto tree=build_clique_tree_blair_peyton(graph_,peo);
    nodes_.resize(static_cast<std::size_t>(graph_.vertex_count()+2));
    for(const auto& clique:tree.cliques)allocate_node(clique);
    for(const auto& edge:tree.edges)link(edge.first,edge.second,edge.weight);
}
\end{lstlisting}

\begin{proposition}\label{prop:a1-pre}
The A1 constructor runs in \(O(n+m)\) worst-case time and space, given a \code{Graph}
built by \code{Graph::from_edges}.
\end{proposition}
\begin{proof}
MCS, the PEO test and Blair--Peyton take \(O(n+m)\) (\cref{tab:static-cost}). Allocating the
\(n+2\) empty node slots and \(n\) empty vertex lists is \(O(n)\). The calls of
\code{allocate_node} insert \(\sum_K|K|\le n+m\) memberships at \(O(1)\) each, and the
\(c-1\le n-1\) calls of \code{link} cost \(O(1)\) each. No structure of size \(\Theta(cn)\) or
\(\Theta(n^2)\) is created.
\end{proof}

\section{Delete-Query}\label{sec:a1-delete-query}
By \cref{thm:delete}, \op{Delete-Query}\((u,v)\) must decide whether exactly one node lies in
\(T_u\cap T_v\). With the vertex lists this is a set intersection of two lists of clique
identifiers, computed with a mark array:

\begin{lstlisting}[style=cpp,firstnumber=124,caption={\file{src/sparse_ibarra_a1.cpp}, lines 124--134: the unique clique containing an edge, or \(-1\).},label={lst:a1-clique-of-edge}]
int SparseIbarraChordalGraphA1::clique_of_edge(int u,int v) const {
    // Inverse incidence lists allow a batched, deterministic O(n) query.
    std::vector<unsigned char> marked(nodes_.size(),0);
    for(int id:nodes_of_[u])marked[id]=1;
    int found=-1;
    for(int id:nodes_of_[v])if(marked[id]) {
        if(found!=-1)return -1;
        found=id;
    }
    return found;
}
\end{lstlisting}

\begin{lstlisting}[style=cpp,firstnumber=185,caption={\file{src/sparse_ibarra_a1.cpp}, lines 185--191: \op{Delete-Query}.},label={lst:a1-delete-query}]
bool SparseIbarraChordalGraphA1::delete_query(int u,int v) const {
    check_vertices(u,v);
    if(u==v)return false;
    if(fast_)return fast_->count_common_cliques(u,v)==1;
    if(!graph_.has_edge(u,v))return false;
    return clique_of_edge(u,v)>=0;
}
\end{lstlisting}

\noindent\code{clique_of_edge} returns the unique node if there is exactly one and \(-1\)
otherwise (it stops at the second common node). The mark array has \(n+2\) entries and the two
lists have at most \(c\le n\) entries each, so \op{Delete-Query} costs \(O(n)\) (plus
\(O(\log n)\) for \code{has_edge}). If the A1\(^+\) index is enabled, the query is delegated to
it instead (\cref{ch:fastdelete}).

\section{Insert-Query: closest pair by multi-source BFS}\label{sec:a1-insert-query}
The insertion query needs the closest pair \(x\in T_u\), \(y\in T_v\), the path between them,
and \(|K_x\cap K_y|\). A1 finds the pair with a breadth-first search that starts
simultaneously from \emph{all} nodes of \(T_u\):

\begin{algorithm}[ht]
\caption{Closest pair of A1 (\code{closest_pair})}\label{alg:closest-bfs}
\begin{algorithmic}[1]
\Require \(T\), the node lists of \(T_u\) and \(T_v\) (disjoint)
\Ensure the closest pair \((x,y)\) and the edges of the \(x\)--\(y\) path
\State mark the nodes of \(T_v\) as targets
\State put every node of \(T_u\) in the queue with \(\mathit{prev}[s]=s\)
\While{the queue is not empty}
  \State \(a\gets\) dequeue; \textbf{if} \(a\) is a target \textbf{then} \(y\gets a\); \textbf{break}
  \For{each tree edge \(e=\{a,b\}\) with \(b\) unvisited}
    \State \(\mathit{prev}[b]\gets a\); \(\mathit{via}[b]\gets e\); enqueue \(b\)
  \EndFor
\EndWhile
\State follow \(\mathit{prev}\) from \(y\) until a node \(x\) with \(\mathit{prev}[x]=x\), collecting the edges \(\mathit{via}[\cdot]\)
\end{algorithmic}
\end{algorithm}

\begin{lemma}\label{lem:bfs-closest}
\Cref{alg:closest-bfs} returns the closest pair \((x,y)\) and the unique \(x\)--\(y\) path of
\(T\), in \(O(c)\subseteq O(n)\) time.
\end{lemma}
\begin{proof}
A multi-source BFS dequeues nodes in non-decreasing order of their distance from the set
\(T_u\). The first target dequeued, \(y\), is therefore a node of \(T_v\) at minimum distance
from \(T_u\). The \(\mathit{prev}\) chain from \(y\) is a shortest path to some source \(x\in
T_u\), and its interior nodes are not sources (a source has distance 0 and is never re-entered).
Hence \((x,y)\) is a pair at minimum distance; in a tree the closest pair between two disjoint
subtrees is unique, so it is \emph{the} closest pair. Every node and edge of \(T\) is handled
at most once.
\end{proof}

\begin{lstlisting}[style=cpp,firstnumber=136,caption={\file{src/sparse_ibarra_a1.cpp}, lines 136--143: intersection size with a vertex mark array.},label={lst:a1-overlap}]
int SparseIbarraChordalGraphA1::overlap(int a,int b) const {
    // One O(n) vertex-mark array, never a c-by-n characteristic matrix.
    std::vector<unsigned char> marked(graph_.vertex_count(),0);
    for(int v:nodes_[a].members)marked[v]=1;
    int count=0;
    for(int v:nodes_[b].members)count+=marked[v];
    return count;
}
\end{lstlisting}

\begin{lstlisting}[style=cpp,firstnumber=193,caption={\file{src/sparse_ibarra_a1.cpp}, lines 193--203: \op{Insert-Query} (\cref{alg:insert-query}).},label={lst:a1-insert-query}]
bool SparseIbarraChordalGraphA1::insert_query(int u,int v) const {
    check_vertices(u,v);
    if(u==v||graph_.has_edge(u,v))return false;
    const auto pair=closest_pair(u,v);
    if(pair.path_edges.empty())return false;
    const int common=overlap(pair.x,pair.y);
    int minimum=graph_.vertex_count();
    for(int id:pair.path_edges)
        minimum=std::min(minimum,edges_[id].weight);
    return minimum==common;
}
\end{lstlisting}

\noindent If \(x\) and \(y\) are adjacent, the path has a single edge whose weight \emph{is}
\(|K_x\cap K_y|\), so the comparison \code{minimum==common} correctly returns \yes; no special
case is needed. If \(T_u\) and \(T_v\) are in different components of \(G\), the path crosses a
dummy edge of weight 0 and \(|K_x\cap K_y|=0\), so the answer is \yes, as it must be.

\section{Delete}\label{sec:a1-delete}
\Cref{lst:a1-delete} implements \cref{alg:delete}. The correspondence is:
\begin{itemize}
\item lines 220--223: find the unique clique \(K_x\) containing \(uv\) (reject otherwise) and
  set \(k=|K_x|\);
\item lines 226--231: build \(K_{xu}=K_x\setminus\{v\}\) and \(K_{xv}=K_x\setminus\{u\}\);
\item lines 234--247: classify every neighbour of \(x\) into \(N_u\) (side 0), \(N_v\) (side 1)
  or \(N_0\) (side 2) by marking \(T_u\) with bit 1 and \(T_v\) with bit 2 --- one pass over two
  lists instead of a membership test per neighbour --- and remember a neighbour of weight
  \(k-1\) on each side (the absorption test \eqref{eq:delete-maximal});
\item lines 249--256: cut all edges of \(x\), turn \(x\) into \(K_{xu}\), create a new node
  \(x_2\) for \(K_{xv}\), reattach the neighbours (side 1 to \(x_2\), the others to \(x\)) and
  add the edge \(\{x,x_2\}\) of weight \(k-2\);
\item lines 257--258: contract \(x\) into the absorbing neighbour on the \(u\) side and \(x_2\)
  into the one on the \(v\) side, if any;
\item line 259: delete the edge from the graph (lines 260--264 update the optional A1\(^+\)
  index).
\end{itemize}

\begin{lstlisting}[style=cpp,firstnumber=219,caption={\file{src/sparse_ibarra_a1.cpp}, lines 219--266: \op{Delete} (\cref{alg:delete}).},label={lst:a1-delete}]
bool SparseIbarraChordalGraphA1::delete_edge(int u,int v) {
    if(!delete_query(u,v))return false;
    const int x=clique_of_edge(u,v);
    if(x<0)throw std::logic_error("A1 unique deletion clique absent");
    const int k=static_cast<int>(nodes_[x].members.size());
    if(fast_)fast_->remove_clique(x);

    std::vector<int> xu,xv;
    xu.reserve(k);xv.reserve(k);
    for(int w:nodes_[x].members) {
        if(w!=v)xu.push_back(w);
        if(w!=u)xv.push_back(w);
    }

    // O(n) batch node membership tests, including high-degree clique stars.
    std::vector<unsigned char> mark(nodes_.size(),0);
    for(int id:nodes_of_[u])mark[id]|=1;
    for(int id:nodes_of_[v])mark[id]|=2;
    struct Neighbor {int other,weight,side;};
    std::vector<Neighbor> neighbors;
    int absorb_u=-1,absorb_v=-1;
    for(int id:nodes_[x].incident) {
        const auto& e=edges_[id];
        const int other=e.a==x?e.b:e.a;
        const int side=(mark[other]&1)?0:((mark[other]&2)?1:2);
        neighbors.push_back({other,e.weight,side});
        if(e.weight==k-1 && side==0)absorb_u=other;
        if(e.weight==k-1 && side==1)absorb_v=other;
    }

    const std::vector<int> old_edges(
        nodes_[x].incident.begin(),nodes_[x].incident.end());
    for(int id:old_edges)cut(id);
    replace_members(x,std::move(xu));
    const int x2=allocate_node(std::move(xv));
    for(const auto& neighbor:neighbors)
        link(neighbor.side==1?x2:x,neighbor.other,neighbor.weight);
    link(x,x2,k-2);
    if(absorb_u>=0)contract(absorb_u,x);
    if(absorb_v>=0)contract(absorb_v,x2);
    graph_.remove_edge(u,v);
    if(fast_) {
        if(nodes_[x].active)add_to_fast_index(x);
        if(nodes_[x2].active)add_to_fast_index(x2);
        fast_->rebalance(static_cast<int>(graph_.edge_count()));
    }
    return true;
}
\end{lstlisting}

\section{Insert}\label{sec:a1-insert}
\Cref{lst:a1-insert} implements \cref{alg:insert}:
\begin{itemize}
\item lines 269--272: run the query (reject if \no), then recompute the closest pair and
  \(|I|=|K_x\cap K_y|\);
\item lines 273--279: if \(x,y\) are not adjacent, cut a minimum-weight edge of the path and
  link \(\{x,y\}\) with weight \(|I|\) (the exchange of \cref{thm:insert6});
\item lines 280--297: locate the edge \(\{x,y\}\) and build \(K_z=I\cup\{u,v\}\) with a vertex
  mark array;
\item lines 299--300: decide whether \(K_x\) and \(K_y\) are absorbed, by
  \eqref{eq:insert-maximal};
\item lines 305--310: subdivide \(\{x,y\}\) by the new node \(z\) (two edges of weight
  \(|I|+1\)) and contract absorbed cliques into \(z\);
\item line 311: insert the edge into the graph.
\end{itemize}

\begin{lstlisting}[style=cpp,firstnumber=268,caption={\file{src/sparse_ibarra_a1.cpp}, lines 268--317: \op{Insert} (\cref{alg:insert}).},label={lst:a1-insert}]
bool SparseIbarraChordalGraphA1::insert_edge(int u,int v) {
    if(!insert_query(u,v))return false;
    const auto pair=closest_pair(u,v);
    const int x=pair.x,y=pair.y;
    const int common=overlap(x,y);
    if(pair.path_edges.size()!=1) {
        int minimum=pair.path_edges.front();
        for(int id:pair.path_edges)
            if(edges_[id].weight<edges_[minimum].weight)minimum=id;
        cut(minimum);
        link(x,y,common);
    }
    int xy=-1;
    for(int id:nodes_[x].incident) {
        const auto& e=edges_[id];
        if((e.a==x&&e.b==y)||(e.b==x&&e.a==y)) {
            xy=id;
            break;
        }
    }
    if(xy<0)throw std::logic_error("A1 missing insertion tree edge");

    std::vector<unsigned char> marked(graph_.vertex_count(),0);
    for(int member:nodes_[x].members)marked[member]=1;
    std::vector<int> kz;
    kz.reserve(static_cast<std::size_t>(common+2));
    for(int member:nodes_[y].members)
        if(marked[member])kz.push_back(member);
    kz.push_back(u);
    kz.push_back(v);

    const bool drop_x=static_cast<int>(nodes_[x].members.size())==common+1;
    const bool drop_y=static_cast<int>(nodes_[y].members.size())==common+1;
    if(fast_) {
        if(drop_x)fast_->remove_clique(x);
        if(drop_y)fast_->remove_clique(y);
    }
    cut(xy);
    const int z=allocate_node(std::move(kz));
    link(x,z,common+1);
    link(z,y,common+1);
    if(drop_x)contract(z,x);
    if(drop_y)contract(z,y);
    graph_.add_edge(u,v);
    if(fast_) {
        add_to_fast_index(z);
        fast_->rebalance(static_cast<int>(graph_.edge_count()));
    }
    return true;
}
\end{lstlisting}

\begin{remark}
Ibarra \S4 describes a shortcut for the absorbed case: instead of creating \(z\) and
contracting, extend the absorbed clique's vector in place. The sparse A1 follows the general
description of \S3.4 (create \(z\), then contract); both take \(O(n)\) time. The dense class
used inside A2 implements the in-place shortcut, because there it saves dynamic-tree operations
(\cref{sec:a2-insert-on}).
\end{remark}

\section{Complexity of A1}\label{sec:a1-complexity}
\begin{theorem}[A1 bounds]\label{thm:a1-bounds}
After \(O(n+m)\) preprocessing, \code{SparseIbarraChordalGraphA1} performs each of
\op{Delete-Query}, \op{Insert-Query}, \op{Delete} and \op{Insert} in \(O(n)\) worst-case time.
\end{theorem}
\begin{proof}
Preprocessing: \cref{prop:a1-pre}. Every operation allocates at most a constant number of mark
arrays of size \(n\) or \(n+2\) (\(O(n)\)), scans at most two vertex lists (each of length
\(\le c\le n\)), at most a constant number of clique member lists (each of length \(\le n\)) and
the clique tree (\(c\le n\) nodes and \(c-1\) edges).
\begin{itemize}
\item \op{Delete-Query}: \code{has_edge} \(O(\log n)\) and \code{clique_of_edge} \(O(n)\).
\item \op{Insert-Query}: \cref{lem:bfs-closest} \(O(n)\); \code{overlap} \(O(n)\); the path
  minimum over at most \(c-1\) edges.
\item \op{Delete}: the query, two member lists of size \(k\le n\), the classification of at
  most \(c-1\) neighbours, \(O(c)\) cuts and links at \(O(1)\) each, \code{replace_members} and
  \code{allocate_node} \(O(k)\), at most two contractions \(O(c+n)\), and
  \code{Graph::remove_edge} \(O(n)\).
\item \op{Insert}: the query, at most one exchange \(O(1)\), \(K_z\) in \(O(n)\), one new node
  \(O(n)\), at most two contractions \(O(c+n)\), and \code{Graph::add_edge}
  \(O(\log n+n)\).
\end{itemize}
Each is \(O(n)\).
\end{proof}

\section{Worked examples with real program output}\label{sec:a1-traces}
All traces in this section were produced by \code{mtp_dynamic_cli} (A1) from the files listed
in \cref{app:traces}; the figures draw exactly the printed cliques and tree edges.

\paragraph{Running example.} \Cref{fig:trace-ex} shows the effect of the script
\begin{lstlisting}[style=shell]
Insert 0 4          => yes
Insert-Query 4 6    => no
Insert-Query 1 6    => yes
Delete-Query 0 1    => yes
Delete 0 1          => yes
\end{lstlisting}
\begin{itemize}
\item \op{Insert}\((0,4)\): \(T_0=\{K_1\}\), \(T_4=\{K_2\}\) are adjacent; \(I=\{2\}\);
  neither \(|K_1|=4\) nor \(|K_2|=3\) equals \(|I|+1=2\), so nothing is absorbed and the new
  clique \(K_3=\{2,0,4\}\) is inserted between them with edges of weight 2.
\item \op{Insert-Query}\((4,6)\): \(T_4=\{K_2,K_3\}\), \(T_6=\{K_0\}\); the closest pair is
  \((K_3,K_0)\) with path \(K_3\!-\!K_1\!-\!K_0\) of weights \(2,1\); \(|K_3\cap K_0|=0<1\),
  so \no.
\item \op{Insert-Query}\((1,6)\): \(K_1\ni 1\) and \(K_0\ni 6\) are adjacent: \yes.
\item \op{Delete}\((0,1)\): only \(K_1=\{0,1,2,3\}\) contains both, \(k=4\). The neighbours
  are \(K_0=\{3,6\}\) (neither side) and \(K_3=\{2,0,4\}\) (side of \(u=0\), weight \(2\ne
  k-1\)), so nothing is absorbed: \(K_1\) becomes \(\{3,2,0\}\), the new node
  \(K_4=\{3,2,1\}\) is attached to it with weight \(k-2=2\).
\end{itemize}

\begin{figure}[ht]
\centering
\begin{tikzpicture}
\begin{scope}
\node[cl] (K0) at (0,3.9) {\(K_0=\{6,3\}\)};
\node[cl] (K1) at (0,2.6) {\(K_1=\{3,2,1,0\}\)};
\node[clg] (K3) at (0,1.3) {\(K_3=\{2,0,4\}\)};
\node[cl] (K2) at (0,0) {\(K_2=\{2,5,4\}\)};
\draw[tedge] (K0)--node[wlab,right]{1}(K1);
\draw[tedge] (K1)--node[wlab,right]{2}(K3);
\draw[tedge] (K3)--node[wlab,right]{2}(K2);
\node[font=\small\bfseries] at (0,4.7) {after \op{Insert}\((0,4)\)};
\end{scope}
\begin{scope}[xshift=6.3cm]
\node[cl] (K0) at (-1.5,3.9) {\(K_0=\{6,3\}\)};
\node[clg] (K1) at (0,2.6) {\(K_1=\{3,2,0\}\)};
\node[clg] (K4) at (2.3,3.9) {\(K_4=\{3,2,1\}\)};
\node[cl] (K3) at (0,1.3) {\(K_3=\{2,0,4\}\)};
\node[cl] (K2) at (0,0) {\(K_2=\{2,5,4\}\)};
\draw[tedge] (K0)--node[wlab,left]{1}(K1);
\draw[tedge] (K1)--node[wlab,right]{2}(K4);
\draw[tedge] (K1)--node[wlab,right]{2}(K3);
\draw[tedge] (K3)--node[wlab,right]{2}(K2);
\node[font=\small\bfseries] at (0.4,4.7) {after \op{Delete}\((0,1)\)};
\end{scope}
\end{tikzpicture}
\caption[A1 trace on the running example]{Clique trees printed by \code{mtp_dynamic_cli} for
the running example after \op{Insert}\((0,4)\) and after \op{Delete}\((0,1)\). New or changed
cliques are green.}
\label{fig:trace-ex}
\end{figure}

\paragraph{Double subsumption on a star.} For the star with centre 0 and leaves 1, 2, 3 the
program builds the tree \(\{0,2\}\!-\!\{3,0\}\!-\!\{0,1\}\). \op{Insert}\((1,3)\): the closest
pair is \(x=\{0,1\}\), \(y=\{3,0\}\) (adjacent), \(I=\{0\}\), and \(|K_x|=|K_y|=2=|I|+1\): both
cliques are absorbed into \(K_z=\{0,1,3\}\) (\cref{fig:trace-star}). Deleting \(\{1,3\}\) again
splits \(K_z\) into \(\{0,1\}\) and \(\{0,3\}\), restoring a star. On a star with \(n-1\) leaves
the centre clique of the tree can have degree \(\Theta(n)\); this is the case that makes A2's
insertion delicate (\cref{sec:a2-insert-on}).

\begin{figure}[ht]
\centering
\begin{tikzpicture}
\begin{scope}
\node[cl] (a) at (0,1.2) {\(\{0,2\}\)};
\node[clh] (b) at (0,0) {\(\{3,0\}\)};
\node[clh] (c) at (0,-1.2) {\(\{0,1\}\)};
\draw[tedge] (a)--node[wlab,right]{1}(b);
\draw[tedge] (b)--node[wlab,right]{1}(c);
\node[font=\small\bfseries] at (0,2.0) {before};
\end{scope}
\draw[arr,accent] (1.4,0)--node[above,font=\small]{\op{Insert}\((1,3)\)}(3.4,0);
\begin{scope}[xshift=4.8cm]
\node[cl] (a) at (0,0.6) {\(\{0,2\}\)};
\node[clg] (z) at (0,-0.6) {\(\{0,1,3\}\)};
\draw[tedge] (a)--node[wlab,right]{1}(z);
\node[font=\small\bfseries] at (0,2.0) {after};
\end{scope}
\draw[arr,accent] (6.2,0)--node[above,font=\small]{\op{Delete}\((1,3)\)}(8.2,0);
\begin{scope}[xshift=9.6cm]
\node[cl] (a) at (0,1.2) {\(\{0,2\}\)};
\node[cl] (b) at (0,0) {\(\{0,1\}\)};
\node[cl] (c) at (0,-1.2) {\(\{0,3\}\)};
\draw[tedge] (a)--node[wlab,right]{1}(b);
\draw[tedge] (b)--node[wlab,right]{1}(c);
\node[font=\small\bfseries] at (0,2.0) {after};
\end{scope}
\end{tikzpicture}
\caption[Double subsumption on a star]{Star \(K_{1,3}\): \op{Insert}\((1,3)\) absorbs both end
cliques (orange) into the new clique (green); \op{Delete}\((1,3)\) splits it again.}
\label{fig:trace-star}
\end{figure}

\paragraph{Path.} On the path \(0\!-\!1\!-\!2\!-\!3\!-\!4\) the tree is the path of edge
cliques with all weights 1. \op{Insert-Query}\((0,2)\) answers \yes\ (adjacent cliques
\(\{1,0\}\), \(\{2,1\}\)), and \op{Insert-Query}\((0,3)\) answers \no\ (path weights \(1,1\) but
\(\{1,0\}\cap\{3,2\}=\emptyset\)): the cycle \(0\!-\!1\!-\!2\!-\!3\!-\!0\) would be chordless.
After \op{Insert}\((0,2)\) the cliques \(\{1,0\}\) and \(\{2,1\}\) are absorbed into
\(\{1,0,2\}\), and now \op{Insert-Query}\((0,3)\) answers \yes.

\paragraph{Fan with an isolated vertex.} For the graph of \cref{fig:dummy} the program answers
\op{Insert-Query}\((0,4)\) \(=\) \no\ (\cref{fig:insert-query-no}) and
\op{Insert-Query}\((0,3)\) \(=\) \yes; after \op{Insert}\((0,3)\) the clique \(\{0,1,2\}\) is
absorbed into \(\{2,1,0,3\}\), and the singleton \(\{5\}\) stays attached by its dummy edge.

\section{The dense A1 used inside A2}\label{sec:dense-a1}
\code{DynamicChordalGraph} (\file{src/dynamic_ibarra.cpp}) is the literal \S4 implementation:
each node stores \code{std::vector<unsigned char> contains} of length \(n\) and its size. Its
queries and updates have the same structure as above and run in \(O(n)\); its constructor runs
in \(\Theta(n+m+cn)\). It differs from the sparse class in two ways that matter for A2:
\begin{enumerate}
\item Every node has \(O(1)\) membership tests \code{contains[v]}, which A2's
  \(W\)-update formulas use for \emph{all} nodes (\cref{sec:w-updates}).
\item It exposes two \emph{hooks}, \code{tree_link_hook_} and \code{tree_cut_hook_}, which
  are called on every \code{link} and \code{cut}; A2 uses them to mirror every change of the
  clique tree into the Sleator--Tarjan forest. When the hooks are installed, \op{Insert} uses
  Ibarra's in-place shortcut for absorbed cliques (\file{src/dynamic_ibarra.cpp}, lines
  274--311) so that a single absorption costs only \(O(1)\) link/cut operations.
\end{enumerate}
The test programs \file{tests/test_ibarra_dynamic.cpp} and
\file{tests/test_sparse_ibarra_a1.cpp} run the same exhaustive oracle checks on both classes.
%% ---- end expanded/07_a1.tex ----
%% ---- begin expanded/08_a2.tex ----
%% =====================================================================
\chapter{Implementation of A2}\label{ch:a2}
%% =====================================================================
A1 spends \(O(n)\) time on every \op{Insert-Query}, because it searches the clique tree for the
closest pair and computes \(|K_x\cap K_y|\) from scratch. Ibarra's second implementation (A2,
\S5) reduces this to \(O(\log^2 n)\) with two ingredients: a \emph{dynamic tree} that answers
path questions about the clique tree in \(O(\log n)\) time each, and a matrix \(W\) that stores
\(|K_x\cap K_y|\) for \emph{every} pair of cliques. The price is that \op{Delete} becomes
\(O(n\log n)\) and that \(W\) must be kept up to date. This chapter first describes the
dynamic tree (\cref{sec:st83,sec:bst85}), which is the largest single component of the project,
and then the clique-level algorithms of A2 (\cref{sec:a2-structure} onwards).

\section{Overview of the A2 class}\label{sec:a2-structure}
\code{SleatorTarjanChordalGraphA2} (\file{include/mtp/dynamic_ibarra_a2.hpp},
\file{src/dynamic_ibarra_a2.cpp}) combines four parts (\cref{fig:a2-parts}):

\begin{lstlisting}[style=cpp,firstnumber=44,caption={\file{include/mtp/dynamic_ibarra_a2.hpp}, lines 44--55: the members of the A2 class.},label={lst:a2-members}]
private:
    DynamicChordalGraph data_;
    SleatorTarjanWorstCaseForest tree_;
    std::vector<int> vertex_witness_;
    std::vector<std::vector<int>> W_;
    int link_count_last_=0,cut_count_last_=0;
    bool rebuilt_last_insert_=false;
    long long rebuild_steps_last_insert_=0;
    [[nodiscard]] std::pair<int,int> closest(int u,int v);
    void initialize_weights();
    void attach_tree_hooks();
};
\end{lstlisting}

\begin{itemize}
\item \code{data_}: a dense A1 (\cref{sec:dense-a1}) that owns the clique tree, the
  characteristic vectors and the graph, and performs Ibarra's tree updates;
\item \code{tree_}: a Sleator--Tarjan worst-case dynamic forest with one node per clique slot
  (\(n+2\) nodes; free slots are isolated nodes). Every clique-tree edge \(\{x,y\}\) is an edge of
  the forest with cost \(w(x,y)\). The two structures are kept identical through the link/cut
  hooks of \code{data_};
\item \code{vertex_witness_[v]}: for every vertex \(v\), some clique node containing \(v\)
  (Ibarra \S5 keeps exactly such a pointer for every vertex);
\item \code{W_}: the \((n+2)\times(n+2)\) matrix of \S5.1 with
  \(W[x][y]=|K_x\cap K_y|\) for all active nodes \(x,y\).
\end{itemize}

\begin{figure}[ht]
\centering
\begin{tikzpicture}[b/.style={boxnode,minimum width=4.3cm,minimum height=12mm}]
\node[boxnode,fill=okgreen!12,draw=okgreen!70!black,minimum width=7cm] (api) at (0,2.7)
  {public operations: \op{Insert-Query}, \op{Insert}, \op{Delete-Query}, \op{Delete}};
\node[b] (data) at (-3.6,1.0) {\code{data_}: dense A1\\clique tree + vectors + graph};
\node[b] (st) at (3.6,1.0) {\code{tree_}: Sleator--Tarjan\\worst-case forest};
\node[b] (wit) at (-3.6,-1.2) {\code{vertex_witness_}\\one clique per vertex};
\node[b] (W) at (3.6,-1.2) {\code{W_}: \((n{+}2)\times(n{+}2)\)\\intersection sizes};
\draw[arr] ([xshift=6mm]api.south west) -- (data.north);
\draw[arr] ([xshift=-6mm]api.south east) -- (st.north);
\draw[arr,dashed] (data.east) -- node[above,font=\scriptsize]{link/cut hooks} (st.west);
\draw[arr] (data.south) -- node[left,font=\scriptsize]{repaired, \(O(n)\)} (wit.north);
\draw[arr] (data.south east) -- node[pos=0.5,below,sloped,font=\scriptsize]{rows updated, \(O(n)\)} (W.north west);
\end{tikzpicture}
\caption[Components of A2]{The four components of the A2 class. \op{Insert-Query} uses only
the witnesses, the dynamic tree and \(W\); updates are carried out by \code{data_} and
mirrored into the other three.}
\label{fig:a2-parts}
\end{figure}

\Cref{tab:st-interface} lists the dynamic-tree operations that A2 needs. The next two sections
explain how the Sleator--Tarjan structure provides them with worst-case \(O(\log n)\) bounds.

\begin{table}[ht]
\centering\small
\caption[Dynamic-tree interface used by A2]{Interface of \code{SleatorTarjanWorstCaseForest}
(\file{include/mtp/worst_case_dynamic_tree.hpp}) for a forest with \(N\) nodes. All bounds are
worst-case.}\label{tab:st-interface}
\begin{tabularx}{\textwidth}{@{}lXl@{}}
\toprule
\textbf{Operation} & \textbf{Meaning} & \textbf{Time}\\
\midrule
\code{link(u,v,c)} & add edge \(\{u,v\}\) of cost \(c\) (different trees) & \(O(\log N)\)\\
\code{cut(u,v)} & delete edge \(\{u,v\}\) & \(O(\log N)\)\\
\code{evert(u)} & make \(u\) the root of its tree & \(O(\log N)\)\\
\code{root(u)} & the root of \(u\)'s tree & \(O(\log N)\)\\
\code{min_on_path(u,v)} & minimum edge cost on the \(u\)--\(v\) path & \(O(\log N)\)\\
\code{kth_on_path(u,v,k)} & the \(k\)-th node (0-based) on the path from \(u\) to \(v\) & \(O(\log N)\)\\
\code{path_length(u,v)} & number of nodes on the \(u\)--\(v\) path & \(O(\log N)\)\\
\code{build_linear_forest(N,E)} & build the whole structure for a given forest (static) & \(O(N)\)\\
\bottomrule
\end{tabularx}
\end{table}

\section{Sleator--Tarjan worst-case dynamic trees}\label{sec:st83}
This section summarises the data structure of Sleator and Tarjan~\cite{sleator1983} (ST83) and
how \file{src/worst_case_dynamic_tree.cpp} implements it. ST83 builds the structure in three
steps, and the implementation follows the last one:
\begin{itemize}
\item \textbf{ST83 \S3} (``Dynamic trees as sets of paths''): every tree is partitioned into
  vertex-disjoint \emph{solid paths} joined by \emph{dashed} edges; all tree operations are
  written in terms of path operations and the two composite operations \code{splice} and
  \code{expose}. This section also defines \(\size(v)\) and heavy and light edges (Lemma~1).
\item \textbf{ST83 \S4} (``Dynamic paths as biased binary trees''): each solid path is
  represented by a binary tree, at first balanced and then \emph{biased}, with vertex weights
  \(\wt(v)\) (Lemmas~2 and~5). With the partition of \S3 this gives \emph{amortised}
  \(O(\log n)\) bounds (Theorem~5).
\item \textbf{ST83 \S5} (``Partitioning by size''): the solid edges are defined to be exactly
  the heavy edges; every operation ends with \code{conceal}, which restores this property;
  solid paths and the new \emph{path sets} are stored in \emph{globally} biased binary trees;
  and every operation takes \(O(\log n)\) time in the \emph{worst case} (Theorems~6--8).
\end{itemize}
The forest is kept rooted (\code{evert} changes the root); \(\size(v)\) is the number of
descendants of \(v\), including \(v\). \Cref{tab:st83-map} maps each ingredient of ST83 to the
code. (Two comments in the code --- \file{src/worst_case_dynamic_tree.cpp}, line 26, and
\file{src/dynamic_ibarra_a2.cpp}, line 245 --- attribute the heavy-edge rule to ST83 \S4; in
ST83 it is defined in \S3 and becomes the rule for solid edges in \S5. This affects only the
comments, not the code.)

\begin{table}[ht]
\centering\small
\caption[ST83 ingredients and their implementation]{Ingredients of the worst-case dynamic tree
of ST83 and where they are implemented (\code{SleatorTarjanWorstCaseForest} and
\code{GloballyBiasedBinaryTree}).}\label{tab:st83-map}
\begin{tabularx}{\textwidth}{@{}>{\raggedright\arraybackslash}p{5.2cm}>{\raggedright\arraybackslash}p{1.9cm}X@{}}
\toprule
\textbf{ST83 ingredient} & \textbf{Where in ST83} & \textbf{Implementation}\\
\midrule
size, heavy/light edges, at most \(\lfloor\lg n\rfloor\) light edges on a root path & \S3, Lemma~1 & heavy-child test \code{2LL*subtree[v]>subtree[u]} in \code{initialize_linear_forest}\\
solid paths with head (bottom) and tail (top); \code{dparent}, \code{dcost} of dashed edges & \S3 & \code{path_root_}, \code{dashed_parent_}, \code{dashed_cost_}\\
\code{splice}, \code{expose} & \S3; path-set updates \S5 & \code{splice}, \code{expose}\\
path as binary tree: external nodes = vertices head to tail, internal nodes = edges; reversal bits; minimum edge cost & \S4 & leaves/internal nodes of \code{GloballyBiasedBinaryTree}; \code{reversed}; \code{edge_cost}, \code{minimum_cost}\\
biased-tree rank properties (i), (ii); depth bound & \S4, Lemma~2 & checks (B1)--(B3) in \code{check_invariants}\\
vertex weight \(\wt(v)\), \(1\le\wt(v)\le n\) & \S4, Lemma~5 & leaf weights; \code{set_detached_leaf_weight}\\
solid \(\iff\) heavy; \code{conceal}, \code{slice} & \S5 & \code{conceal}, \code{slice}\\
\code{lefttilt}, \code{leftmin}, \code{righttilt}, \code{rightmin}; \code{light}\((p)\) & \S5 & \code{left_tilt_min}, \code{right_tilt_min}; \code{light}\\
globally biased trees; concatenate and split bounds & \S5, Lemmas~6--7 (citing BST85) & \code{global_join}, \code{split}\\
path sets ordered by decreasing weight; \code{maxwt}, \code{insert}, \code{delete} & \S5 & \code{path_set_root_}; \code{max_path}, \code{insert_path_set}, \code{delete_path_set}\\
\code{root}, \code{mincost}, \code{link}, \code{cut}, \code{evert} & \S5 & \code{root}, \code{min_on_path}, \code{link}, \code{cut}, \code{evert}\\
worst-case \(O(\log n)\) per operation & \S5, Theorems~6--8 & (the bound relied on by A2)\\
\bottomrule
\end{tabularx}
\end{table}

\subsection{The size partition}
\begin{definition}[Heavy and light edges; ST83 \S3]
An edge from a node \(v\) to its parent \(p\) is \emph{heavy} if \(2\,\size(v)>\size(p)\), and
\emph{light} otherwise.
\end{definition}
Every node has at most one heavy child, so the heavy edges form vertex-disjoint paths. If the
edge from \(v\) to its parent is light, then \(\size(p)\ge 2\,\size(v)\); hence any path from a
node to the root contains at most \(\lfloor\lg N\rfloor\) light edges (ST83 Lemma~1). In the
worst-case structure of ST83 \S5 an edge is \emph{solid} (stored inside a path representation)
if and only if it is heavy, except temporarily inside an operation; the other edges are
\emph{dashed}. A path from a node to the root therefore meets at most \(\lfloor\lg N\rfloor+1\)
solid paths. \Cref{fig:st-partition} shows an example.

\begin{figure}[ht]
\centering
\begin{tikzpicture}[level distance=11mm,
  tn/.style={circle,draw=accent!85!black,fill=accentlight,thick,minimum size=6.4mm,inner sep=0pt,font=\small},
  sz/.style={font=\scriptsize,text=black!60},
  heavy/.style={line width=1.8pt,draw=accent},
  light/.style={thick,dashed,draw=black!55}]
\node[tn] (r) at (0,0) {r};
\node[tn] (a) at (-1.4,-1.1) {a};
\node[tn] (b) at (1.6,-1.1) {b};
\node[tn] (c) at (-2.2,-2.2) {c};
\node[tn] (d) at (-0.6,-2.2) {d};
\node[tn] (h) at (1.1,-2.2) {h};
\node[tn] (i) at (2.1,-2.2) {i};
\node[tn] (e) at (-3.0,-3.3) {e};
\node[tn] (f) at (-1.4,-3.3) {f};
\node[tn] (g) at (-3.6,-4.4) {g};
\node[tn] (j) at (-2.4,-4.4) {j};
\draw[heavy] (r)--(a); \draw[heavy] (a)--(c); \draw[heavy] (c)--(e);
\draw[light] (r)--(b); \draw[light] (a)--(d); \draw[light] (c)--(f);
\draw[light] (e)--(g); \draw[light] (e)--(j); \draw[light] (b)--(h); \draw[light] (b)--(i);
\foreach \n/\s in {r/11,a/7,b/3,c/5,d/1,e/3,f/1,g/1,h/1,i/1,j/1}{\node[sz,right=0.5mm of \n]{\s};}
\node[font=\small,align=left,anchor=west] at (3.3,-0.4) {\textbf{solid path} (head \(\to\) tail):};
\node[font=\small,anchor=west] at (3.3,-1.0) {\(e\,[3]\ \ c\,[2]\ \ a\,[2]\ \ r\,[4]\)\quad(\(\wt\) in brackets)};
\node[font=\small,anchor=west] at (3.3,-1.6) {prefix sums \(3,5,7,11\) \(=\) the sizes};
\node[font=\small,align=left,anchor=west] at (3.3,-2.5) {\textbf{path sets} (dashed children, by weight):};
\node[font=\small,anchor=west] at (3.3,-3.1) {\(r:\{b\,[3]\}\), \(a:\{d\,[1]\}\), \(c:\{f\,[1]\}\),};
\node[font=\small,anchor=west] at (3.3,-3.6) {\(e:\{g\,[1],j\,[1]\}\), \(b:\{h\,[1],i\,[1]\}\)};
\end{tikzpicture}
\caption[Size partition of a rooted tree]{The size partition of an 11-node tree rooted at
\(r\) (subtree sizes in small print). Thick edges are heavy (\(2\size(v)>\size(p)\)) and form
the solid path \(e\!-\!c\!-\!a\!-\!r\), whose head (bottom) is \(e\) and tail (top) is \(r\);
dashed edges are light. Every other node is a solid path of length zero. Each node keeps its
dashed children's paths in a ``path set'' (ST83 \S5).}
\label{fig:st-partition}
\end{figure}

\subsection{Representing a solid path (ST83 \S4)}
Each solid path is represented by a full binary tree whose external nodes (leaves), in
left-to-right order, are the vertices of the path from head to tail, and whose internal nodes,
in symmetric order, are the edges of the path. An internal node therefore stands for the
boundary between the rightmost leaf of its left subtree and the leftmost leaf of its right
subtree, and stores the cost of that edge. Every node stores aggregates of its subtree:
\begin{itemize}
\item \code{minimum_cost}: the minimum edge cost in the subtree --- ST83's \emph{grossmin};
  it answers \code{min_on_path} at the root in \(O(1)\);
\item \code{first}, \code{last}: the leftmost and rightmost leaves --- ST83's \emph{bhead}
  and \emph{btail};
\item \code{reversed}: ST83's reversal bit, so that \code{evert} can reverse a whole path in
  \(O(1)\) (the bit is pushed down when a node is visited);
\item \code{sum_weight}: the total weight of the leaves;
\item \code{left_tilt_min}, \code{right_tilt_min}: ST83's \emph{leftmin} and \emph{rightmin}
  (below);
\item \code{leaf_count}: the number of leaves --- an augmentation not in ST83, used for
  \code{path_length} and \code{kth_on_path} (below).
\end{itemize}

\begin{caution}[Two representation details that differ from ST83 (disclosed)]
ST83 stores costs and tilt minima in a difference form (\emph{netcost}, \emph{netmin},
\emph{netleftmin}, \emph{netrightmin}) so that \emph{pupdate} --- adding a value to every edge
cost on a path --- takes \(O(1)\). A2 never adds a value to the costs of a path, so the
implementation stores the plain minima and recomputes them in \(O(1)\) per node whenever a node
is created or rotated; this changes no bound. The \code{leaf_count} field is an addition needed
for Ibarra's ``\(i\)-th node on a path'' operation, which ST83 does not list.
\end{caution}

\paragraph{Weights.} ST83 \S4 gives every vertex the weight
\[ \wt(v)=\begin{cases}\size(v) & \text{if no solid edge enters } v,\\
\size(v)-\size(w) & \text{if the solid edge } (w,v) \text{ enters } v,\end{cases} \]
i.e.\ one plus the sizes of the subtrees hanging from \(v\) by dashed edges, and proves
\(1\le\wt(v)\le n\) (Lemma~5). For a solid path \(v_1,\dots,v_k\) listed from head to tail,
\(\size(v_i)=\sum_{j\le i}\wt(v_j)\) (\cref{fig:st-partition}). In a biased binary tree an
external node \(u\) of a tree with root \(r\) has depth at most
\(2\lg(\wt(r)/\wt(u))+4\) (ST83 Lemma~2). Along a root path these ratios telescope, which is
what makes an \code{expose} that meets \(O(\log N)\) solid paths cost \(O(\log N)\) in total
rather than \(O(\log^2 N)\) (ST83 Theorems~5 and~6).

\paragraph{Detecting light solid edges with tilts (ST83 \S5).} For the internal node \(u\) of
the edge \((v_i,v_{i+1})\), ST83 defines \(\mathit{lefttilt}(u)\) as the total weight of the
vertices from the head up to \(v_i\) minus \(\wt(v_{i+1})\), and observes that the edge is heavy
iff \(\mathit{lefttilt}(u)>0\). Indeed, the edge is heavy iff
\(2\size(v_i)>\size(v_{i+1})=\size(v_i)+\wt(v_{i+1})\), i.e.\ iff
\[ \textstyle \sum_{j\le i}\wt(v_j)\;-\;\wt(v_{i+1})\;>\;0 . \]
The aggregate \code{left_tilt_min} (ST83's \emph{leftmin}) is the minimum of this quantity over
the internal nodes of a subtree (\code{right_tilt_min} is the mirror image, needed after
reversals). The path operation \code{light}\((p)\) of ST83 \S5 --- ``the vertex \(w\) closest
to the tail of \(p\) such that the edge entering \(w\) on \(p\) is light'' --- is found by one
root-to-leaf descent guided by these minima (function \code{light} in
\file{src/globally_biased_binary_tree.cpp}).

\subsection{Path sets (ST83 \S5)}
For every vertex \(v\), the solid paths whose tails hang from \(v\) by dashed edges form the
\emph{path set} of \(v\). ST83 \S5 represents it by a globally biased binary tree whose external
nodes are these paths, in left-to-right order by \emph{decreasing weight} (the weight of a path
being the sum of its vertex weights, i.e.\ the size of the dashed subtree), with a pointer to
the leftmost leaf. Hence \code{maxwt}\((v)\) --- the heaviest dashed child path --- takes
\(O(1)\) (function \code{max_path}), and inserting or deleting a path \(p\) takes
\(O(\log(W/\wt(p)))\), where \(W\) is the total weight of the path set.

\subsection{Expose, conceal and the operations}
All operations are built from two procedures:
\begin{itemize}
\item \code{expose(u)} (ST83 \S3, with the path-set updates of \S5) makes the whole path from
  \(u\) to the root a single solid path with head \(u\). It splits \(u\)'s solid path at \(u\),
  turns the part below \(u\) into a dashed child (inserting it into \(u\)'s path set), and then
  repeatedly \code{splice}s: the tail of the current path hangs from some vertex \(v\) by a
  dashed edge; split \(v\)'s solid path at \(v\), move the part below \(v\) into \(v\)'s path
  set, remove the current path from it, and concatenate. There is one splice per light edge on
  the path to the root, hence \(O(\log N)\) splices, and the whole \code{expose} takes
  \(O(\log N)\) time (ST83 Theorem~6).
\item \code{conceal(p)} (ST83 \S5) is ``an expose run backwards'': it converts to dashed every
  light edge on \(p\) by \code{slice} operations (each located with \code{light}), and converts
  to solid every heavy edge incident to \(p\), found as the heaviest path of a path set
  (``\(2\,\wt(q)>\wt(v)\)'' in ST83, \code{2LL*weight(heavy)>vertex_[v]->weight} in the code). It
  takes \(O(\log N)\) time (ST83 Theorem~7).
\end{itemize}

\begin{lstlisting}[style=cpp,firstnumber=297,caption={\file{src/worst_case_dynamic_tree.cpp}, lines 297--313: \code{expose}.},label={lst:expose}]
// ST83 Section 5, published expose procedure, with path-set changes.
SleatorTarjanWorstCaseForest::Node*
SleatorTarjanWorstCaseForest::expose(int u) {
    auto parts=split_vertex(u);
    int revised=vertex_[u]->weight+weight(parts.before);
    paths_.set_detached_leaf_weight(vertex_[u],revised);
    if(parts.before) {
        const int low_tail=tail(parts.before);
        dashed_parent_[low_tail]=u;
        dashed_cost_[low_tail]=parts.before_cost;
        insert_path_set(parts.before,u);
    }
    Node* p=vertex_[u];
    if(parts.after)p=concatenate(p,parts.after,parts.after_cost);
    while(dashed_parent_[tail(p)]>=0)p=splice(p);
    return p;
}
\end{lstlisting}

The public operations are then short compositions of ST83 \S5's procedures
(\file{src/worst_case_dynamic_tree.cpp}, lines 368--436):
\begin{itemize}
\item \code{evert(u)}: \code{expose(u)}, reverse the exposed path, clear \(u\)'s dashed parent,
  \code{conceal} --- exactly ST83's \code{evert};
\item \code{root(u)}: the tail of \code{expose(u)}, then \code{conceal} --- exactly ST83's
  \code{root};
\item \code{link(u,v,c)}: ST83's \code{link}\((v,w,x)\) requires \(v\) to be a tree root; it
  exposes \(w\), concatenates the path of \(v\) with the exposed path using an edge of cost \(x\),
  and conceals the result. The implementation first checks that
  \(u\) and \(v\) are in different trees and calls \code{evert(u)}, so that \(u\) is a root, and
  then performs exactly that sequence;
\item \code{cut(u,v)}: ST83's \code{cut}\((v)\) removes the edge from \(v\) to its parent; the
  implementation calls \code{evert(u)} so that \(u\) is the parent of \(v\), then follows
  ST83's program (expose, split, conceal both parts);
\item \code{min_on_path(u,v)}: \code{evert(u)} and then ST83's \code{mincost}\((v)\), which
  exposes \(v\) and reads the minimum at the root of the exposed path;
\item \code{path_length(u,v)} and \code{kth_on_path(u,v,k)}: \code{evert(u)}, \code{expose(v)},
  read \code{leaf_count} at the root or descend to the \(k\)-th leaf, \code{conceal}.
\end{itemize}

\begin{lstlisting}[style=cpp,firstnumber=383,caption={\file{src/worst_case_dynamic_tree.cpp}, lines 383--415: \code{link}, \code{cut} and \code{min_on_path}.},label={lst:link-cut}]
void SleatorTarjanWorstCaseForest::link(int u,int v,int w) {
    verify_vertex(u);verify_vertex(v);
    if(u==v || root(u)==root(v))
        throw std::invalid_argument("ST link would form a cycle");
    evert(u);
    Node* p=expose(v);
    Node* united=concatenate(path(u),p,w);
    conceal(united);
}
void SleatorTarjanWorstCaseForest::cut(int u,int v) {
    verify_vertex(u);verify_vertex(v);
    if(u==v)throw std::invalid_argument("invalid self cut");
    evert(u);
    Node* p=expose(v);
    auto parts=split_vertex(v);
    if(parts.before || !parts.after || parts.after->leaf_count!=1 ||
       first(parts.after)!=u)
        throw std::invalid_argument("ST cut requires an existing edge");
    dashed_parent_[v]=-1;
    conceal(vertex_[v]);
    conceal(parts.after);
}
int SleatorTarjanWorstCaseForest::min_on_path(int u,int v) {
    verify_vertex(u);verify_vertex(v);
    if(u==v)return 1000000000;
    if(root(u)!=root(v))
        throw std::invalid_argument("ST path between disconnected nodes");
    evert(u);
    Node* p=expose(v);
    const int result=paths_.min_cost(p);
    conceal(p);
    return result;
}
\end{lstlisting}

\begin{theorem}[Sleator and Tarjan~\cite{sleator1983}, \S5, Theorem~8]\label{thm:st83}
With partitioning by size and a representation of solid paths as globally biased binary
trees, any dynamic tree operation takes \(O(\log n)\) time.
\end{theorem}

\begin{corollary}\label{cor:st83-aug}
In the implementation, \code{link}, \code{cut}, \code{evert}, \code{root},
\code{min_on_path}, \code{path_length} and \code{kth_on_path} each take \(O(\log N)\)
worst-case time on an \(N\)-node forest.
\end{corollary}
\begin{proof}
Each is a constant number of ST83 operations (evert, expose, conceal, concatenate, split), which
are \(O(\log N)\) by \cref{thm:st83}, plus, for \code{kth_on_path}, a descent from the root of
the exposed path to one leaf. By ST83 Lemmas~2 and~5 every leaf of a tree with root \(r\) has
depth at most \(2\lg(\wt(r)/1)+4\le 2\lg N+4\), so the descent is \(O(\log N)\).
\end{proof}

\begin{caution}[What is proved and what is tested]
The \(O(\log N)\) bound is Sleator and Tarjan's theorem; it holds for this implementation
provided the implementation maintains the structural invariants of ST83 and of globally biased
trees. The project does not contain a machine-checked proof of that correspondence. Instead,
the invariants are checked explicitly by debug checkers (\code{check_invariants}) after every
operation in the tests: vertex weights equal one plus the dashed path weights, every dashed edge
has exactly one path-set entry, path sets are sorted by decreasing weight, and every biased tree
satisfies the rank conditions of \cref{sec:bst85} (\cref{sec:tests-st}).
\end{caution}

\section{Globally biased binary trees}\label{sec:bst85}
ST83 \S4 uses \emph{biased binary trees} and lists the two properties of ranks that its analysis
needs; ST83 \S5 replaces them by \emph{globally biased} binary trees, ``which have the same
static properties'' but worst-case concatenate and split bounds (ST83 Lemmas~6 and~7), citing
Bent, Sleator and Tarjan~\cite{bent1985} (BST85) for both. The implementation
(\file{src/globally_biased_binary_tree.cpp}) provides these trees.

\begin{caution}[Source availability]
The BST85 paper itself was not available during the writing of this report. The ST83
properties (B1)--(B3) below and ST83 Lemmas~6--7 were checked against ST83. The global-bias
condition (B4), the case analysis of the global join and the linear-time bottom-up construction
are described as implemented, with the BST85 figure and theorem numbers that the source-code
comments cite (Figure~14; Figure~15 and Theorem~12); these citations could not be checked
against BST85.
\end{caution}

\paragraph{Ranks.} Every leaf of weight \(w\ge 1\) has rank \(r=\lfloor\log w\rfloor\) (ST83
property~(i)). This is computed in \(O(1)\) by a fixed sequence of five shifts (a binary search
over the bits of a 32-bit word):

\begin{lstlisting}[style=cpp,firstnumber=11,caption={\file{src/globally_biased_binary_tree.cpp}, lines 11--25: constant-time rank of a weight.},label={lst:floorlog}]
int GloballyBiasedBinaryTree::floor_log2(int weight) {
    if(weight<=0)throw std::invalid_argument("nonpositive leaf weight");
    // Word-RAM O(1) (not a repeated division loop):
    // weight is a 32-bit signed positive count <= number of vertices.
    // Each branch halves the search field by a fixed power of 2.
    // Five constant-time shift steps establish the exact bit rank.
    unsigned int x=static_cast<unsigned int>(weight);
    int result=0;
    if(x>=(1u<<16)) {x>>=16;result+=16;}
    if(x>=(1u<<8)) {x>>=8;result+=8;}
    if(x>=(1u<<4)) {x>>=4;result+=4;}
    if(x>=(1u<<2)) {x>>=2;result+=2;}
    if(x>=2u)++result;
    return result;
}
\end{lstlisting}

\paragraph{Invariants.} The checker \code{check_invariants} enforces, for every node \(x\) with
children \(c\) and grandchildren \(g\):
\begin{enumerate}[label=(B\arabic*)]
\item a leaf of weight \(w\) has rank \(\lfloor\log w\rfloor\), and a node whose leaves have
  total weight \(W\) has rank at most \(1+\lfloor\log W\rfloor\) (ST83 \S4, property~(i));
\item \(\mathrm{rank}(c)\le\mathrm{rank}(x)\), with strict inequality if \(c\) is a leaf
  (ST83 \S4, property~(ii));
\item \(\mathrm{rank}(g)<\mathrm{rank}(x)\) (ST83 \S4, property~(ii));
\item (\emph{global bias}, from BST85 as cited in the code) if
  \(\mathrm{rank}(x)-\mathrm{rank}(c)\ge 2\), then the leaves immediately to the left and right
  of \(c\)'s subtree (in symmetric order) have rank at least \(\mathrm{rank}(x)-1\).
\end{enumerate}
Properties (B1)--(B3) give the depth bound of ST83 Lemma~2; (B4) is the additional condition
that, according to the code's citation of BST85, makes worst-case concatenation and splitting
possible.

\paragraph{Operations.}
\begin{itemize}
\item \code{global_join(x, y, c)}: concatenate two trees (all leaves of \(x\) before those of
  \(y\)), with \(c\) the cost of the new boundary --- ST83's \code{concatenate}, whose worst-case
  time is bounded by ST83 Lemma~6. The implementation follows the case analysis that its comments
  attribute to BST85 Figure~14 (cases 1--3 for unequal root ranks, case 4 with its sub-cases for
  equal ranks), using rotations called \emph{tilts}.
\item \code{split(leaf)}: return the trees of the leaves before it, the leaf itself and the
  leaves after it, together with the two boundary costs --- ST83's \code{split}, bounded by ST83
  Lemma~7.
\item \code{build_linear(leaves, costs)}: the \emph{bottom-up} construction that the code
  attributes to BST85 (Figure~15, Theorem~12), which builds a globally biased tree on \(k\)
  ordered leaves in \(O(k)\) worst-case time: leaves are added left to right; for each new leaf \(y\) the algorithm
  climbs the right spine from the previous leaf to the highest node whose parent has rank at
  most \(\mathrm{rank}(y)\), inserts a new internal node there, and repairs the spine with
  tilt-left rotations that stop at the first rotation without a rank promotion. The aggregates
  are computed afterwards in one post-order pass, so they do not spoil the linear bound.
\end{itemize}

\begin{lstlisting}[style=cpp,firstnumber=173,caption={\file{src/globally_biased_binary_tree.cpp}, lines 173--228: the bottom-up construction attributed in the code to BST85 (main loop).},label={lst:buildlinear}]
GloballyBiasedBinaryTree::build_linear(
    const std::vector<Node*>& leaves,
    const std::vector<int>& boundary_cost) {
    if(leaves.empty()) {
        if(!boundary_cost.empty())
            throw std::invalid_argument("empty leaf sequence has costs");
        return nullptr;
    }
    if(boundary_cost.size()+1!=leaves.size())
        throw std::invalid_argument("one cost per leaf boundary required");
    linear_steps_=0;
    for(Node* l:leaves)
        if(!l||!l->is_leaf()||l->parent||!l->live)
            throw std::invalid_argument("linear builder requires detached live leaves");
    Node* root=leaves.front();
    Node* rightmost=root;
    for(std::size_t i=1;i<leaves.size();++i) {
        Node* y=leaves[i];
        Node* x=rightmost;
        // Figure 15: find the maximal right-path node whose
        // parent's rank does not exceed rank(y).
        while(x->parent && x->parent->rank<=y->rank) {
            x=x->parent;
            ++linear_steps_;
        }
        Node* old_parent=x->parent;
        Node* added=internal(x,y,std::max(x->rank,y->rank)+1,
                             boundary_cost[i-1]);
        ++linear_steps_;
        if(!old_parent) {
            root=added;
        } else {
            // x was on the rightmost spine of the old search tree.
            if(old_parent->right!=x)
                throw std::logic_error("linear builder lost rightmost spine");
            set_right(old_parent,added);
            Node* current=old_parent;
            for(;;) {
                Node* ancestor=current->parent;
                if(ancestor && ancestor->right!=current)
                    throw std::logic_error("linear tilt left spine broken");
                const int previous_rank=current->rank;
                Node* revised=tilt_left(current);
                if(ancestor)set_right(ancestor,revised);
                else root=revised;
                ++linear_steps_;
                // The paper stops after the first tilt without a rank
                // promotion. Rotations preserve the rank of that parent.
                if(revised!=current || current->rank==previous_rank ||
                   !ancestor)
                    break;
                current=ancestor;
            }
        }
        rightmost=y;
    }
\end{lstlisting}

\section{Insert-Query in \texorpdfstring{\(O(\log^2 n)\)}{O(log\string^2 n)}}\label{sec:a2-insert-query}
Ibarra's \op{Insert-Query} for A2 (\S5, p.\,40:12) finds the closest pair by binary search on
the path between the two witnesses. Its correctness rests on the following observation.

\begin{lemma}[Prefix--suffix structure]\label{lem:prefix}
Let \(w\in T_u\), \(z\in T_v\) and let \(P=(p_0=w,p_1,\dots,p_\ell=z)\) be the \(w\)--\(z\) path
in \(T\). Then \(P\cap T_u\) is a prefix \(p_0,\dots,p_a\) of \(P\), \(P\cap T_v\) is a suffix
\(p_b,\dots,p_\ell\) with \(a<b\), and \((p_a,p_b)\) is the closest pair \((x,y)\).
\end{lemma}
\begin{proof}
\(T_u\) is a subtree containing \(p_0\), and the intersection of a path with a subtree is a
subpath; hence it is a prefix. Symmetrically \(P\cap T_v\) is a suffix, and the two are disjoint
because no clique contains both \(u\) and \(v\). In a tree, every path from a node of \(T_u\) to
a node of \(T_v\) contains the unique shortest path between the two subtrees, whose interior
avoids both; so this shortest path is \(p_a,\dots,p_b\).
\end{proof}

\begin{figure}[ht]
\centering
\begin{tikzpicture}[x=1.35cm]
\foreach \i/\s in {0/u,1/u,2/u,3/,4/,5/v,6/v}{
  \ifx\s u \node[clh,minimum width=8mm] (q\i) at (\i,0) {\(p_{\i}\)};
  \else\ifx\s v \node[clg,minimum width=8mm] (q\i) at (\i,0) {\(p_{\i}\)};
  \else \node[cl,minimum width=8mm] (q\i) at (\i,0) {\(p_{\i}\)};\fi\fi
}
\foreach \i [count=\j from 1] in {0,...,5}{\draw[tedge] (q\i)--(q\j);}
\draw[decorate,decoration={brace,amplitude=5pt}] (-0.3,0.5)--node[above=5pt,font=\small]{contain \(u\) (prefix)}(2.3,0.5);
\draw[decorate,decoration={brace,amplitude=5pt}] (4.7,0.5)--node[above=5pt,font=\small]{contain \(v\) (suffix)}(6.3,0.5);
\node[font=\small,below=2mm of q0] {\(w\)}; \node[font=\small,below=2mm of q6] {\(z\)};
\node[font=\small,below=2mm of q2] {\(x\)}; \node[font=\small,below=2mm of q5] {\(y\)};
\end{tikzpicture}
\caption[Binary search for the closest pair in A2]{The path between the witnesses \(w\) of \(u\)
and \(z\) of \(v\). The cliques containing \(u\) form a prefix and those containing \(v\) a
suffix, so \(x\) (last of the prefix) and \(y\) (first of the suffix) can be found by binary
search with \code{kth_on_path}.}
\label{fig:binsearch}
\end{figure}

\begin{algorithm}[ht]
\caption{A2 \op{Insert-Query}\((u,v)\) (Ibarra \S5, p.\,40:12)}\label{alg:a2-iq}
\begin{algorithmic}[1]
\If{\(u=v\) or \(uv\in E\)} \Return \no \EndIf \Comment{\(O(\log n)\) with \code{std::set}}
\State \(w\gets\mathit{witness}[u]\); \(z\gets\mathit{witness}[v]\); \(\ell\gets\) \code{path_length}\((w,z)\)
\State \(x\gets p_a\) where \(a=\max\{i: u\in K_{p_i}\}\), by binary search over \(i\in[0,\ell-1]\) using \code{kth_on_path}\((w,z,i)\)
\State \(y\gets p_b\) where \(b=\min\{i: v\in K_{p_i}\}\), by binary search likewise
\State \Return \yes\ if \code{min_on_path}\((x,y)=W[x][y]\), else \no \Comment{\cref{thm:insert6}}
\end{algorithmic}
\end{algorithm}

\begin{lstlisting}[style=cpp,firstnumber=67,caption={\file{src/dynamic_ibarra_a2.cpp}, lines 67--100: closest pair by binary search and \op{Insert-Query}.},label={lst:a2-iq}]
std::pair<int,int> SleatorTarjanChordalGraphA2::closest(int u,int v) {
    const int w=vertex_witness_.at(u),z=vertex_witness_.at(v);
    if(w<0||z<0||!data_.nodes_[w].active||!data_.nodes_[z].active||
       !data_.nodes_[w].contains[u]||!data_.nodes_[z].contains[v])
        throw std::logic_error("stale Ibarra A2 clique witnesses");

    // Ibarra (2008), Section 5, paragraph Insert-Query:
    // binary search the nearest u/v clique nodes on the w--z path.
    const int count=tree_.path_length(w,z);
    int left=0,right=count-1;
    while(left<right) {
        const int mid=(left+right+1)/2;
        const int candidate=tree_.kth_on_path(w,z,mid);
        if(data_.nodes_[candidate].contains[u])left=mid;
        else right=mid-1;
    }
    const int x=tree_.kth_on_path(w,z,left);
    left=0;right=count-1;
    while(left<right) {
        const int mid=(left+right)/2;
        const int candidate=tree_.kth_on_path(w,z,mid);
        if(data_.nodes_[candidate].contains[v])right=mid;
        else left=mid+1;
    }
    const int y=tree_.kth_on_path(w,z,left);
    return {x,y};
}
bool SleatorTarjanChordalGraphA2::insert_query(int u,int v) {
    if(u==v||data_.graph_.has_edge(u,v))return false;
    const auto [x,y]=closest(u,v);
    if(x==y)return false;
    // Ibarra, Theorem 6, constant-time stored clique intersection W.
    return tree_.min_on_path(x,y)==W_[x][y];
}
\end{lstlisting}

\begin{proposition}\label{prop:a2-iq}
A2's \op{Insert-Query} takes \(O(\log^2 n)\) worst-case time.
\end{proposition}
\begin{proof}
The adjacency test costs \(O(\log n)\). Each binary search makes \(O(\log n)\) probes, each a
\code{kth_on_path} call of cost \(O(\log n)\) (\cref{cor:st83-aug}) plus an \(O(1)\) vector lookup
\code{contains[u]}. \code{path_length} and \code{min_on_path} are \(O(\log n)\) (\cref{cor:st83-aug}), and
\(W[x][y]\) is \(O(1)\). Total \(O(\log^2 n)\).
\end{proof}

Note that when \(x\) and \(y\) are adjacent, \code{min_on_path}\((x,y)=w(x,y)=W[x][y]\), so the
adjacent case of \cref{alg:insert-query} needs no special treatment, exactly as in A1.

\section{Maintaining the intersection matrix \texorpdfstring{\(W\)}{W}}\label{sec:w-updates}
\(W\) is indexed by clique-node slots. After every update only the rows (and columns) of
\emph{new or changed} cliques change, and Ibarra \S5.1 shows that each changed row can be
computed in \(O(n)\) from old rows. We state the formulas as lemmas, with proofs, because they
are the part of A2 most easily implemented wrongly. Below, \([v\in K_i]\) is 1 or 0, read in
\(O(1)\) from the characteristic vector of node \(i\).

\begin{lemma}[\(W\) after \op{Delete}; Ibarra \S5.1]\label{lem:w-delete}
Let \op{Delete}\((u,v)\) split \(K_x\) (\(|K_x|=k\)) into \(K_{xu}=K_x\setminus\{v\}\), kept in
node \(x\), and \(K_{xv}=K_x\setminus\{u\}\), in a new node \(x_2\). For every other active node
\(i\):
\[ W[x][i]=W_{\mathrm{old}}[x][i]-[v\in K_i],\qquad W[x_2][i]=W_{\mathrm{old}}[x][i]-[u\in K_i],
\qquad W[x][x_2]=k-2. \]
If \(K_{xu}\) or \(K_{xv}\) is absorbed into a neighbour, that node's row is unchanged and the
absorbed node disappears.
\end{lemma}
\begin{proof}
Since \(v\in K_x\), \(|(K_x\setminus\{v\})\cap K_i|=|K_x\cap K_i|-[v\in K_i]\); similarly for
\(u\). \(K_{xu}\cap K_{xv}=K_x\setminus\{u,v\}\) has \(k-2\) elements. An absorbing clique is an
old clique that does not change.
\end{proof}

\begin{lemma}[\(W\) after \op{Insert}; Ibarra \S5.1]\label{lem:w-insert}
Let \op{Insert}\((u,v)\) create \(K_z=I\cup\{u,v\}\) with \(I=K_x\cap K_y\). For every other
active node \(i\):
\begin{enumerate}
\item if \(K_x\) is absorbed (\(K_x=I\cup\{u\}\)): \(W[z][i]=W[x][i]+[v\in K_i]\);
\item else if \(K_y\) is absorbed: \(W[z][i]=W[y][i]+[u\in K_i]\);
\item otherwise, let \(T_x,T_y\) be the two components of \(T\) after removing the exchanged edge
  (so \(x\in T_x\), \(y\in T_y\)): \(W[z][i]=W[y][i]+[u\in K_i]\) for \(i\in T_x\), and
  \(W[z][i]=W[x][i]+[v\in K_i]\) for \(i\in T_y\).
\end{enumerate}
\end{lemma}
\begin{proof}
(1) \(K_z=K_x\cup\{v\}\) with \(v\notin K_x\). (2) is symmetric. (3) Let \(i\in T_x\). In the
new clique tree \(T'=T-e+\{x,y\}\), the path from \(i\) to \(y\) passes through \(x\), so by the
clique intersection property \(K_i\cap K_y\subseteq K_x\), i.e.\ \(K_i\cap K_y=K_i\cap I\).
Moreover \(v\notin K_i\): otherwise \(T_v\) (a subtree of \(T'\) containing \(i\) and \(y\))
would contain \(x\), but \(v\notin K_x\). Therefore
\(|K_z\cap K_i|=|I\cap K_i|+[u\in K_i]=W[y][i]+[u\in K_i]\). The case \(i\in T_y\) is
symmetric.
\end{proof}

\paragraph{Implementation.} \code{delete_edge} (lines 116--169) records the old row of \(x\),
lets \code{data_} perform the deletion, identifies the new node \(x_2\) as the only slot that
became active, and applies \cref{lem:w-delete}:

\begin{lstlisting}[style=cpp,firstnumber=150,caption={\file{src/dynamic_ibarra_a2.cpp}, lines 150--167: \(W\) update after \op{Delete} (\cref{lem:w-delete}).},label={lst:a2-wdel}]
    if(data_.nodes_[x].active) {
        for(int i=0;i<cap;++i) {
            if(i==x||i==x2||!data_.nodes_[i].active)continue;
            const int value=old_row[i]-data_.nodes_[i].contains[v];
            W_[x][i]=W_[i][x]=value;
        }
        W_[x][x]=k-1;
    }
    if(x2>=0) {
        for(int i=0;i<cap;++i) {
            if(i==x||i==x2||!data_.nodes_[i].active)continue;
            const int value=old_row[i]-data_.nodes_[i].contains[u];
            W_[x2][i]=W_[i][x2]=value;
        }
        W_[x2][x2]=k-1;
        if(data_.nodes_[x].active)
            W_[x2][x]=W_[x][x2]=k-2;
    }
\end{lstlisting}

\noindent \code{insert_edge} (lines 173--292) computes the side \(T_x\) by a depth-first search
that does not cross the exchanged edge (lines 199--220) and the new row \emph{before} the tree
is modified, while the old rows and vectors are still available:

\begin{lstlisting}[style=cpp,firstnumber=222,caption={\file{src/dynamic_ibarra_a2.cpp}, lines 222--232: the new row of \(W\) after \op{Insert} (\cref{lem:w-insert}).},label={lst:a2-wins}]
    // Precompute all new weights while the original row values and
    // characteristic vectors are still available.
    std::vector<int> new_row(cap,0);
    for(int i=0;i<cap;++i) {
        if(!old_active[i])continue;
        if(drop_x) new_row[i]=old_x[i]+data_.nodes_[i].contains[v];
        else if(drop_y) new_row[i]=old_y[i]+data_.nodes_[i].contains[u];
        else if(on_x_side[i])
            new_row[i]=old_y[i]+data_.nodes_[i].contains[u];
        else new_row[i]=old_x[i]+data_.nodes_[i].contains[v];
    }
\end{lstlisting}

\noindent As a safeguard, \code{insert_edge} also recomputes \(|K_x\cap K_y|\) directly (an
\(O(n)\) scan, allowed in \op{Insert} but never in \op{Insert-Query}) and throws an exception if
it disagrees with \(W[x][y]\).

\begin{example}[\(W\) on the running example]
After \op{Insert}\((0,4)\) the cliques are \(K_0=\{3,6\}\), \(K_1=\{0,1,2,3\}\),
\(K_2=\{2,4,5\}\), \(K_3=\{0,2,4\}\) (A1 numbering). Here \(x=K_1\), \(y=K_2\), nothing is
absorbed, and \(T_x=\{K_0,K_1\}\). \Cref{lem:w-insert}(3) gives
\(W[K_3][K_0]=W[K_2][K_0]+[0\in K_0]=0+0=0\),
\(W[K_3][K_1]=W[K_2][K_1]+[0\in K_1]=1+1=2\) and
\(W[K_3][K_2]=W[K_1][K_2]+[4\in K_2]=1+1=2\), which are indeed \(|K_3\cap K_i|\).
After \op{Delete}\((0,1)\), \(K_1\) becomes \(\{0,2,3\}\) and \(K_4=\{1,2,3\}\);
\cref{lem:w-delete} gives \(W[K_1][K_3]=2-[1\in K_3]=2\), \(W[K_4][K_3]=2-[0\in K_3]=1\) and
\(W[K_1][K_4]=4-2=2\), again the true intersection sizes.
\end{example}

The formulas were additionally checked by the independent script
\file{audit/verify_w_updates.py}, which applies them on random chordal graphs and compares with
directly computed intersections (\cref{sec:audit-scripts}).

\section{Keeping the dynamic tree and the clique tree identical}
Every \code{link}/\code{cut} performed by \code{data_} on the clique tree calls the hooks
installed by A2, which perform the same \code{link}/\code{cut} on \code{tree_}:

\begin{lstlisting}[style=cpp,firstnumber=47,caption={\file{src/dynamic_ibarra_a2.cpp}, lines 47--56: mirroring clique-tree changes into the dynamic tree.},label={lst:a2-hooks}]
void SleatorTarjanChordalGraphA2::attach_tree_hooks() {
    data_.tree_link_hook_=[this](int a,int b,int cost) {
        tree_.link(a,b,cost);
        ++link_count_last_;
    };
    data_.tree_cut_hook_=[this](int a,int b) {
        tree_.cut(a,b);
        ++cut_count_last_;
    };
}
\end{lstlisting}

\noindent The witnesses are repaired after each update in \(O(n)\): after \op{Delete}, every
vertex whose witness was \(x\) is pointed to \(x\), \(x_2\) or the absorbing neighbour, as
appropriate; after \op{Insert}, every vertex whose witness became inactive is pointed to the
new clique \(K_z\) (an absorbed clique is contained in \(K_z\)).

\section{Delete in \texorpdfstring{\(O(n\log n)\)}{O(n log n)}}
\op{Delete} detaches all \(\deg_T(x)\) edges of \(x\), reattaches them to \(x\) or \(x_2\),
adds \(\{x,x_2\}\) and possibly contracts: \(O(\deg_T(x))=O(n)\) link/cut operations, each
\(O(\log n)\) by \cref{thm:st83}. The \(W\) and witness updates are \(O(n)\), and so is the work
of \code{data_}. Hence \op{Delete} costs \(O(n\log n)\), matching Ibarra \S5 (p.\,40:12). A2's
\op{Delete-Query} is that of \code{data_}: \(O(n)\).

\section{Insert in \texorpdfstring{\(O(n)\)}{O(n)}}\label{sec:a2-insert-on}
Ibarra states that A2's \op{Insert} takes \(O(n)\). The dynamic-tree work of an insertion
depends on which cliques are absorbed:
\begin{description}
\item[Case 1: nothing absorbed.] At most one exchange (one cut, one link), then the edge
  \(\{x,y\}\) is replaced by \(\{x,z\},\{z,y\}\) (one cut, two links): \(O(1)\) operations,
  \(O(\log n)\) time.
\item[Case 2: exactly one absorbed, say \(K_x\).] Ibarra \S4 extends \(K_x\)'s vector in place
  (\(K_x\leftarrow K_x\cup\{v\}\)) instead of creating \(z\); apart from the possible exchange,
  the only tree change is that the weight of \(\{x,y\}\) grows by one (one cut, one link). The dense class does exactly this when
  A2's hooks are installed (\file{src/dynamic_ibarra.cpp}, lines 279--311). \(O(\log n)\) time.
\item[Case 3: both absorbed (double subsumption).] \(K_x\) and \(K_y\) merge into one node, which
  inherits the neighbours of both. If both have \(\Theta(n)\) tree neighbours, mirroring the
  contraction edge by edge costs \(\Theta(n)\) link/cut operations, i.e.\ \(\Theta(n\log n)\)
  time --- more than the claimed \(O(n)\).
\end{description}

Case 3 really occurs: take a vertex \(r\) adjacent to \(u\) and \(v\), and give \(u\) and \(v\)
each \(k\) pendant neighbours. Inserting \(uv\) absorbs \(\{r,u\}\) and \(\{r,v\}\) into
\(\{r,u,v\}\), and for suitable vertex labellings the clique tree built by Blair--Peyton attaches
\(\Theta(k)\) pendant cliques to each of them. The benchmark \code{mtp_a2_two_hub_stress}
searches for such labellings (\cref{tab:two-hub}).

\begin{keyidea}[The \(O(n)\) repair for case 3]
In case 3 the implementation does not mirror the contraction edge by edge. It temporarily
removes the hooks, lets \code{data_} perform the whole clique-tree update in \(O(n)\), and then
rebuilds the \emph{same} Sleator--Tarjan structure from scratch for the new clique tree with
\code{build_linear_forest}, which takes \(O(n)\) worst-case time (\cref{lem:build-forest}).
\end{keyidea}

\begin{lstlisting}[style=cpp,firstnumber=253,caption={\file{src/dynamic_ibarra_a2.cpp}, lines 253--262: case 3: detach the hooks before the clique-tree update \ldots},label={lst:a2-bulk1}]
    const bool bulk_rebuild=drop_x && drop_y;
    if(bulk_rebuild) {
        data_.tree_link_hook_={};
        data_.tree_cut_hook_={};
    }
    const bool changed=data_.insert_edge(u,v);
    if(!changed) {
        if(bulk_rebuild)attach_tree_hooks();
        throw std::logic_error("insert precondition changed");
    }
\end{lstlisting}

\begin{lstlisting}[style=cpp,firstnumber=281,caption={\file{src/dynamic_ibarra_a2.cpp}, lines 281--290: \ldots{} and rebuild the Sleator--Tarjan forest in linear time.},label={lst:a2-bulk2}]
    if(bulk_rebuild) {
        std::vector<SleatorTarjanWorstCaseForest::Edge> live;
        live.reserve(static_cast<std::size_t>(data_.active_count_));
        for(const auto& edge:data_.edges_)
            if(edge.active)live.push_back({edge.a,edge.b,edge.weight});
        tree_=SleatorTarjanWorstCaseForest::build_linear_forest(cap,live);
        rebuild_steps_last_insert_=tree_.static_build_steps();
        rebuilt_last_insert_=true;
        attach_tree_hooks();
    }
\end{lstlisting}

\begin{algorithm}[ht]
\caption{\code{build_linear_forest}: static construction of the ST83 structure}\label{alg:build-forest}
\begin{algorithmic}[1]
\Require a forest on \(N\) nodes with \(\le N-1\) weighted edges
\State root every tree; compute parents, subtree sizes and the heavy child of every node
  (\(2\size(c)>\size(v)\)) by one DFS \Comment{\(O(N)\)}
\State \(\wt(v)\gets\size(v)-\size(\text{heavy child})\) for every \(v\) \Comment{\(O(N)\)}
\For{every node \(b\) without a heavy child} \Comment{bottom of a heavy path}
  \State collect the heavy path from \(b\) upwards and build its globally biased tree with
    \code{build_linear} \Comment{\(O(\text{path length})\)}
  \State if the top of the path has a parent, record the dashed edge
\EndFor
\State bucket-sort all dashed paths by weight (integers in \([1,N]\)), in decreasing order \Comment{\(O(N)\)}
\For{every node \(v\) with dashed children}
  \State build \(v\)'s path set with \code{build_linear} on its sorted dashed paths
\EndFor
\end{algorithmic}
\end{algorithm}

\begin{lemma}\label{lem:build-forest}
Assume that the bottom-up construction \code{build_linear} builds a globally biased tree on
\(k\) ordered leaves in \(O(k)\) worst-case time (the bound the code attributes to BST85
Theorem~12; see \cref{sec:bst85}). Then \code{build_linear_forest} runs in \(O(N)\) worst-case
time and produces the ST83 \S5 structure, in which the solid edges are exactly the heavy edges.
\end{lemma}
\begin{proof}
The DFS, sizes, heavy children and weights take \(O(N)\). Every node lies on exactly one heavy
path, so the solid-path builds take \(O(\sum\text{path lengths})=O(N)\) by the assumption.
Every dashed edge contributes one leaf to exactly one path set, so there are at most \(N-1\)
path-set leaves; the bucket sort over weights in \([1,N]\) (ST83 Lemma~5) is \(O(N)\), and the
path-set builds take \(O(N)\) in total, again by the assumption. Ranks are computed in \(O(1)\)
each. Discarding the old structure is \(O(N)\). By construction the solid edges are the heavy
edges, the weights are those of ST83 \S4, and the path sets are ordered by decreasing weight.
\end{proof}

Replacing the structure by a freshly built one does not affect the cost of later operations:
the bounds of ST83 \S5 (Lemmas~6--7, Theorems~6--8) are worst-case bounds that depend only on
the invariants of the current structure, unlike the amortised analysis of ST83 \S\S3--4, which
relies on credits saved by earlier operations.

\begin{theorem}[A2 \op{Insert} in \(O(n)\)]\label{thm:a2-insert}
Under the assumption of \cref{lem:build-forest}, every \op{Insert} of A2 takes \(O(n)\)
worst-case time.
\end{theorem}
\begin{proof}
The preceding \op{Insert-Query} is \(O(\log^2 n)\). The clique-tree update by \code{data_} is
\(O(n)\) (\cref{thm:a1-bounds} for the dense class), the \(W\) row and the witnesses are
\(O(n)\) (\cref{lem:w-insert}). The dynamic-tree work is \(O(\log n)\) in cases 1 and 2
(\cref{cor:st83-aug}), and \(O(N)=O(n)\) in case 3 by \cref{lem:build-forest}, with
\(N=n+2\).
\end{proof}

\begin{caution}[Deviation from the literal text of the paper (disclosed)]
Ibarra does not describe rebuilding the dynamic tree. The rebuild is a composition of
published results --- Ibarra's clique-tree update (\S3.4, \S4), the ST83 \S5 size-partitioned
worst-case structure, and a linear-time construction of globally biased trees that the code
attributes to BST85 (\S4, Figure~15, Theorem~12) --- and the represented data structure is
exactly the ST83 structure, not a different kind of dynamic tree. What changes is only the
\emph{schedule} of maintenance in case 3. The \(O(n)\) bound for case 3 rests on the linear-time
bound of that construction, which could not be checked against BST85 itself; the measured step
counts (\cref{tab:two-hub}, \cref{fig:rebuild}) grow linearly, as it predicts. The means of
meeting the paper's \(O(n)\) bound in case 3 is ours.
\end{caution}

\begin{table}[ht]
\centering\small
\caption[Double-subsumption insertions in A2]{Output of \code{mtp_a2_two_hub_stress}:
double-subsumption insertions on two-hub trees with \(n=2k+3\). ``ST calls'' are individual
link/cut calls during the insertion; ``rebuild steps'' counts the elementary steps of
\code{build_linear_forest}.}\label{tab:two-hub}
\begin{tabular}{@{}rrrrrr@{}}
\toprule
\(n\) & \(\deg(x)\) & \(\deg(y)\) & ST calls & rebuilt & rebuild steps\\
\midrule
27 & 13 & 13 & 0 & yes & 358\\
51 & 25 & 25 & 0 & yes & 692\\
99 & 49 & 2 & 0 & yes & 1349\\
\bottomrule
\end{tabular}
\end{table}

The rebuild steps grow linearly (about \(13.3n\)--\(13.6n\)); \cref{ch:experiments} shows the
same for larger \(n\). These numbers illustrate, but do not replace, \cref{lem:build-forest}.

\section{Preprocessing of A2 as implemented}\label{sec:a2-preprocessing}
Following Ibarra \S5.1 (``begin with an empty graph \ldots\ and insert each edge of \(G\) using
Insert''), the constructor
\begin{enumerate}
\item computes an MCS ordering of the input and checks that it is a PEO (\(O(n+m)\));
\item creates \code{data_} on the \emph{empty} graph with \(n\) vertices: \(n\) singleton
  cliques, each with a characteristic vector of length \(n\) (\(\Theta(n^2)\)), joined by
  \(n-1\) dummy edges;
\item creates the ST forest with \(n+2\) nodes and links the dummy edges (\(O(n\log n)\));
\item sets the witnesses (\(\Theta(n^2)\) as written: one scan of every vector) and initialises
  the \((n+2)\times(n+2)\) matrix \(W\) (\(\Theta(n^2)\));
\item inserts the \(m\) edges of the input, processing \(i=n,\dots,1\) and inserting the edges
  from \(v_i\) to its later neighbours; every intermediate graph has the PEO \(v_1,\dots,v_n\)
  and is therefore chordal, so every insertion is accepted (\(m\) insertions, \(O(n)\) each).
\end{enumerate}

\begin{lstlisting}[style=cpp,firstnumber=11,caption={\file{src/dynamic_ibarra_a2.cpp}, lines 11--45: the A2 constructor (replay of the edges in reverse PEO order).},label={lst:a2-ctor}]
SleatorTarjanChordalGraphA2::SleatorTarjanChordalGraphA2(Graph graph)
    : data_(Graph(graph.vertex_count())), tree_(graph.vertex_count()+2) {
    // Ibarra Sec.5.1: construct W using chordality-preserving edge insertions
    // beginning with n isolated vertices in reverse PEO order.
    const auto peo=maximum_cardinality_search(graph);
    if(!is_perfect_elimination_ordering(graph,peo))
        throw std::invalid_argument("A2 initial graph must be chordal");

    // Ibarra Sec.2: the initial clique tree contains weight-zero dummy
    // edges joining disconnected singleton cliques.
    for(const auto& e:data_.edges_)
        if(e.active)tree_.link(e.a,e.b,e.weight);
    const int n=graph.vertex_count();
    vertex_witness_.assign(n,-1);
    for(int c=0;c<static_cast<int>(data_.nodes_.size());++c)
        if(data_.nodes_[c].active)
            for(int v=0;v<n;++v)
                if(data_.nodes_[c].contains[v])vertex_witness_[v]=c;
    for(int c:vertex_witness_)
        if(c<0)throw std::logic_error("missing singleton clique witness");
    initialize_weights();
    attach_tree_hooks();

    std::vector<int> position(n,-1);
    for(int i=0;i<n;++i)position[peo[i]]=i;
    for(int i=n-1;i>=0;--i) {
        const int u=peo[i];
        for(int v:graph.neighbors(u)) {
            if(position[v]<=i)continue;
            if(!insert_edge(u,v))
                throw std::logic_error("paper-specified PEO insertion rejected");
        }
    }
    link_count_last_=cut_count_last_=0;
}
\end{lstlisting}

\begin{proposition}\label{prop:a2-pre}
The A2 constructor runs in \(\Theta(n^2+mn)\) worst-case time; in particular it takes
\(\Theta(n^2)\) time even when \(m=0\).
\end{proposition}
This does \emph{not} meet the \(O(m+n\log n)\) preprocessing bound in Ibarra's summary table.
\Cref{ch:preprocessing} discusses this discrepancy in detail and explains why it cannot be
removed by a better implementation of the same method.

\section{Complexity summary for A2}
\begin{table}[ht]
\centering\small
\caption[A2: claimed versus implemented bounds]{A2: bounds claimed by Ibarra versus bounds of
the implementation (all worst-case).}\label{tab:a2-summary}
\begin{tabular}{@{}llll@{}}
\toprule
\textbf{Operation} & \textbf{Ibarra \S7} & \textbf{Implementation} & \textbf{Reason}\\
\midrule
\op{Insert-Query} & \(O(\log^2 n)\) & \(O(\log^2 n)\) & \cref{prop:a2-iq}\\
\op{Insert} & \(O(n)\) & \(O(n)\) & \cref{thm:a2-insert}\\
\op{Delete-Query} & \(O(n)\) & \(O(n)\) & dense clique scan\\
\op{Delete} & \(O(n\log n)\) & \(O(n\log n)\) & \(O(n)\) ST operations\\
preprocessing & \(O(m+n\log n)\) & \(\Theta(n^2+mn)\) & \cref{prop:a2-pre}, \cref{ch:preprocessing}\\
\bottomrule
\end{tabular}
\end{table}

\section{The indexed comparison variant}\label{sec:indexed}
The repository also contains \code{IndexedDynamicChordalGraph}
(\file{src/indexed_dynamic_ibarra.cpp}). It answers \op{Insert-Query} like A2 (witnesses,
binary search, matrix of intersections) but, instead of a dynamic tree, it rebuilds a static
heavy-path decomposition of the clique tree with a segment tree for path minima after every
update, in \(O(n)\). Its \op{Insert-Query} is \(O(\log^2 n)\) and its updates are \(O(n)\), but
it is \emph{not} the data structure that Ibarra specifies, so it is used only as a baseline in
tests and benchmarks and is clearly labelled as such in its header.
%% ---- end expanded/08_a2.tex ----
%% ---- begin expanded/09_fast.tex ----
%% =====================================================================
\chapter{Faster Deletion Queries: A1\texorpdfstring{\(^+\)}{+} and A2\texorpdfstring{\(^+\)}{+}}\label{ch:fastdelete}
%% =====================================================================
Both A1 and A2 answer \op{Delete-Query}\((u,v)\) in \(O(n)\) time by scanning cliques. Ibarra
\S6 adds an optional data structure that answers it in \(O(\sqrt m)\) worst-case time, which is
faster on sparse graphs (e.g.\ planar chordal graphs, where \(m=O(n)\)). This chapter describes
that structure and its implementation, \code{FastCliqueIntersectionIndex}
(\file{include/mtp/fast_delete_ibarra.hpp}, \file{src/fast_delete_ibarra.cpp}). It is enabled
in the three dynamic programs by the command-line flag \code{--fast-delete}.

\section{Reformulation as a set-intersection problem}
By \cref{thm:delete}, \(G-uv\) is chordal iff exactly one maximal clique contains both \(u\)
and \(v\), i.e.\ iff
\[ |\KG(u)\cap\KG(v)| = 1 . \]
So it suffices to maintain the \(n\) sets \(S_v=\KG(v)\) (one per vertex, whose elements are
maximal cliques) under the changes caused by updates, and to answer ``how large is
\(S_u\cap S_v\)?''. For a non-edge the answer is 0, so the query does not even need an
adjacency test. The total size of the sets is \(s=\sum_v|\KG(v)|=\sum_K|K|\), which is at most
\(2m\) plus the number of isolated vertices (Ibarra Lemma~7; \cref{sec:cliques-background}).

\section{The modified Yellin structure}
Yellin~\cite{yellin1994} gave a dynamic set-intersection structure with amortised bounds.
Ibarra \S6.2 modifies it, using that both the universe of elements (cliques) and the set
indices (vertices) are at most \(n\), to obtain worst-case bounds. With a parameter \(p\) (set to
the current \(m\)):

\begin{paperbox}[Ibarra \S6.2 (pp.\,40:14--40:15) --- the modified structure]
\begin{itemize}
\item For every element \(x\) (a maximal clique), the set indices containing it (its vertices)
  are partitioned into \(x.\mathit{member}\) and \(x.\mathit{overflow}\), each stored as a
  characteristic vector plus an unsorted doubly linked list. Invariant:
  \(|x.\mathit{member}|\le\sqrt p\), and \(x.\mathit{overflow}\ne\emptyset\) only if
  \(|x.\mathit{member}|=\sqrt p\).
\item A list \(\mathit{Overflows}\) of the elements with non-empty overflow; it has at most
  \(s/\sqrt p\) entries.
\item An \(n\times n\) array \(\mathit{COMMON}\), where \(\mathit{COMMON}[i,j]\) is the number of
  elements \(x\) with \(i,j\in x.\mathit{member}\).
\end{itemize}
\(\mathit{intersect}(i,j)\) returns \(\mathit{COMMON}[i,j]\) plus the number of elements \(x\)
on \(\mathit{Overflows}\) that contain both \(i\) and \(j\) with at least one of them in
\(x.\mathit{overflow}\); this takes \(O(s/\sqrt p)\) time.
\end{paperbox}

With \(p=m\) and \(s\le 2m+n_0\) (where \(n_0\) counts isolated vertices, whose singleton cliques
never overflow when \(\sqrt m\ge 1\)), at most \(2\sqrt m\) elements can overflow, since each
overflowing element has at least \(\sqrt m\) members. Hence \op{Delete-Query} costs
\(O(\sqrt m)\).

\begin{figure}[ht]
\centering
\begin{tikzpicture}[
  cell/.style={draw=black!60,minimum width=6mm,minimum height=5.5mm,font=\scriptsize,inner sep=0pt}]
\node[font=\small\bfseries,anchor=west] at (-0.4,2.3) {element \(x\) (a maximal clique \(K_x\), here \(|K_x|=6\), \(\lceil\sqrt m\rceil=4\))};
\node[font=\small,anchor=east] at (1.6,1.4) {member:};
\foreach \v [count=\i] in {a,b,c,d}{\node[cell,fill=teal!10] at (1.6+0.6*\i,1.4) {\(\v\)};}
\node[font=\small,anchor=east] at (1.6,0.6) {overflow:};
\foreach \v [count=\i] in {e,f}{\node[cell,fill=orange!15] at (1.6+0.6*\i,0.6) {\(\v\)};}
\node[font=\small,align=left,anchor=west] at (5.0,1.0) {\(\mathit{COMMON}[i,j]\) counts \(x\) for all\\pairs \(i,j\) in \(\{a,b,c,d\}\); \(x\) is on the\\list \(\mathit{Overflows}\) because \(e,f\) overflow};
\node[font=\small,align=left,anchor=west] at (-0.4,-0.5) {\(\mathit{intersect}(a,e)\): not in \(\mathit{COMMON}\); found by scanning \(\mathit{Overflows}\) (\(e\in x.\mathit{overflow}\)).\\
\(\mathit{intersect}(a,b)\): counted in \(\mathit{COMMON}[a,b]\); \(x\) is skipped in the scan (both are members).};
\end{tikzpicture}
\caption[The modified Yellin structure]{One element of the modified Yellin structure. Only
pairs of \emph{member} vertices are counted in \(\mathit{COMMON}\); pairs involving an overflow
vertex are counted by scanning the short list of overflowing elements.}
\label{fig:yellin}
\end{figure}

\section{Implementation}
The class keeps \code{elements_} (one per clique slot, \(n+2\) in total), the matrix
\code{common_} (\(n\times n\)), the list \code{overflows_} and the threshold
\(\lceil\sqrt m\rceil\) (at least 1). Each element stores its member and overflow lists and,
for \(O(1)\) membership tests and removals, two bit vectors and two iterator vectors indexed by
vertex.

\begin{lstlisting}[style=cpp,firstnumber=100,caption={\file{src/fast_delete_ibarra.cpp}, lines 100--113: the intersection count used by \op{Delete-Query}.},label={lst:yellin-query}]
int FastCliqueIntersectionIndex::count_common_cliques(int u,int v) const {
    if(u<0||v<0||u>=n_||v>=n_)
        throw std::out_of_range("invalid vertex");
    int count=common_[u][v];
    // Ibarra Section 6.2: only overflow elements may contribute
    // intersections not already included in COMMON[u,v].
    for(int id:overflows_) {
        const auto& e=elements_[id];
        if((e.member_bit[u]||e.overflow_bit[u])&&
           (e.member_bit[v]||e.overflow_bit[v])&&
           (e.overflow_bit[u]||e.overflow_bit[v]))++count;
    }
    return count;
}
\end{lstlisting}

\noindent The three conditions in the loop are exactly Ibarra's three cases ((member, overflow),
(overflow, member), (overflow, overflow)). Adding a member vertex increments \code{common_} for
all current members (\(O(\sqrt m)\)):

\begin{lstlisting}[style=cpp,firstnumber=16,caption={\file{src/fast_delete_ibarra.cpp}, lines 16--25: adding a member vertex to an element.},label={lst:yellin-add}]
void FastCliqueIntersectionIndex::add_member(int id,int v) {
    auto& e=elements_[id];
    if(e.member_bit[v]||e.overflow_bit[v])
        throw std::logic_error("duplicate clique member");
    for(int w:e.member){++common_[v][w];++common_[w][v];}
    ++common_[v][v];
    e.member.push_back(v);
    e.member_iterator[v]=std::prev(e.member.end());
    e.member_bit[v]=1;
}
\end{lstlisting}

\paragraph{Integration with the updates.} The dynamic classes call the index whenever a clique
changes: before an update they \code{remove_clique} every clique that will be modified or
absorbed, after the update they \code{add_clique} every new or modified clique, and finally they
call \code{rebalance(m)} for the new edge count (see the \code{if(fast_)} blocks in
\cref{lst:a1-delete,lst:a1-insert}). \code{rebalance} moves vertices between member and overflow
lists so that the invariant holds for the new threshold \(\lceil\sqrt m\rceil\).

\begin{caution}[Implementation choice]
Ibarra \S6.3 describes the updates as three cases per operation that rename elements and move
single indices. The implementation instead removes the changed cliques and re-adds the new
ones. This is simpler and has the same worst-case cost, \(O(m+n)\) per update, as shown below.
\end{caution}

\section{Complexity}
\begin{proposition}\label{prop:yellin}
With the index enabled: (i) \op{Delete-Query} takes \(O(\sqrt m+1)\) worst-case time;
(ii) each \op{Delete} or \op{Insert} spends an additional \(O(m+n)\) worst-case time on the
index; (iii) enabling the index takes \(O(n^2+mn)\) time.
\end{proposition}
\begin{proof}
(i) One array lookup plus a scan of \code{overflows_}, which has at most \(2\sqrt m+O(1)\)
entries, each tested in \(O(1)\) with the bit vectors.
(ii) An update changes \(O(1)\) cliques. Removing a clique costs
\(O(|\mathit{member}|^2+|K|)=O(m+n)\); adding one costs \(O(n)\) to scan its characteristic
vector plus \(O(|\mathit{member}|^2)=O(m)\) for \code{common_}. Since \(m\) changes by one,
\(\lceil\sqrt m\rceil\) changes by at most one, so \code{rebalance} moves at most one vertex per
element, and only elements with at least \(\sqrt m-1\) vertices can need a move; there are
\(O(\sqrt m)\) of them, each move costing \(O(\sqrt m)\). With the \(O(n)\) loop over all
element slots this is \(O(m+n)\).
(iii) Allocating \code{common_} costs \(\Theta(n^2)\); each of at most \(n\) cliques is added in
\(O(n+m)\).
\end{proof}

Combined with \cref{thm:a1-bounds,tab:a2-summary} this gives exactly the A1\(^+\) and A2\(^+\)
rows of Ibarra's table (\cref{tab:paper-bounds}): \op{Delete-Query} \(O(\sqrt m)\),
\op{Delete} \(O(m+n)\) resp.\ \(O(m+n\log n)\), \op{Insert} \(O(m+n)\), preprocessing
\(O(mn+n^2)\).

\section{Testing}
The tests \file{tests/test_ibarra_fast.cpp} and \file{tests/test_sparse_ibarra_a1_fast.cpp}
repeat the exhaustive oracle tests of A1 with the index enabled, and
\file{tests/test_ibarra_fast_stress.cpp} exercises the two extreme regimes: stars with
\(n\in\{128,256\}\) (many small cliques, no overflow) and complete graphs with
\(n\in\{100,200\}\) (one huge clique with many overflow vertices), under legal deletions and
insertions that split and merge cliques. A2\(^+\) is tested by
\file{tests/test_ibarra_a2.cpp} and by the CLI test \code{ibarra_a2_fast_delete_sample_cli}.
%% ---- end expanded/09_fast.tex ----
%% ---- begin expanded/10_split.tex ----
%% =====================================================================
\chapter{Dynamic Split Graphs}\label{ch:split}
%% =====================================================================
Ibarra \S8 gives a strikingly simple fully dynamic algorithm for split graphs: after
\(O(m+n)\) preprocessing it maintains, under arbitrary edge insertions and deletions of an
\emph{arbitrary} graph, whether the graph is split, in \(O(1)\) worst-case time per update and
per query. This chapter explains the method and its implementation in
\file{include/mtp/dynamic_split_ibarra.hpp} and \file{src/dynamic_split_ibarra.cpp}.

\section{The degree-sequence test}
By the Hammer--Simeone theorem (\cref{thm:hs}), whether \(G\) is split depends only on its
degree sequence \(d_1\ge\dots\ge d_n\): with \(j=\max\{i:d_i\ge i-1\}\), \(G\) is split iff
\(\sum_{i\le j}d_i=j(j-1)+\sum_{i>j}d_i\), and then \(\omega(G)=j\).
Note that the condition \(d_i\ge i-1\) holds for a prefix of the sorted sequence (if it holds for
\(d_l\) it holds for all \(d_k\), \(k<l\); Ibarra p.\,40:18). So the sequence splits into a left
part \(L_l=(d_1,\dots,d_j)\) and a right part \(L_r=(d_{j+1},\dots,d_n)\).

\begin{example}[After Ibarra Fig.~9]
The graph of \cref{fig:split-example} has degree sequence \(4,3,3,2,1,1\) and \(j=3\);
\(4+3+3=10=3\cdot 2+(2+1+1)\), so it is split with clique number 3. Inserting \(\{d,e\}\) gives
\(4,3,3,3,2,1\) with \(j=4\): \(13\ne 12+3\), not split. Then inserting \(\{b,d\}\) gives
\(4,4,4,3,2,1\) with \(j=4\): \(15=12+3\), split again with clique number 4 (clique
\(\{a,b,c,d\}\)). These are exactly the three degree sequences of Ibarra's Figure~9 example; we
checked all three cases independently with NetworkX (split \(\iff\) \(G\) and \(\bar G\) chordal).
\end{example}

\begin{figure}[ht]
\centering
\begin{tikzpicture}
\begin{scope}
\node[vtxh] (a) at (0,1.2) {a};
\node[vtxh] (b) at (1.4,1.2) {b};
\node[vtxh] (c) at (0.7,0.1) {c};
\node[vtx] (d) at (-1.2,-1.0) {d};
\node[vtx] (e) at (-1.4,1.9) {e};
\node[vtx] (f) at (2.7,0.4) {f};
\draw[gedge] (a)--(b) (a)--(c) (b)--(c) (a)--(d) (a)--(e) (b)--(f) (c)--(d);
\node[font=\small,align=center] at (0.6,-2.0) {degrees \(4,3,3,2,1,1\)\\\(j=3\); split};
\end{scope}
\begin{scope}[xshift=6.6cm]
\node[vtxh] (a) at (0,1.2) {a};
\node[vtxh] (b) at (1.4,1.2) {b};
\node[vtxh] (c) at (0.7,0.1) {c};
\node[vtxh] (d) at (-1.2,-1.0) {d};
\node[vtx] (e) at (-1.4,1.9) {e};
\node[vtx] (f) at (2.7,0.4) {f};
\draw[gedge] (a)--(b) (a)--(c) (b)--(c) (a)--(d) (a)--(e) (b)--(f) (c)--(d) (d)--(e);
\draw[newedge] (b)--(d);
\node[font=\small,align=center] at (0.6,-2.0) {after \(+de\), \(+bd\):\\degrees \(4,4,4,3,2,1\); \(j=4\); split};
\end{scope}
\end{tikzpicture}
\caption[A split graph and an update sequence]{Left: a split graph (clique \(\{a,b,c\}\) in
orange, independent set \(\{d,e,f\}\)). Right: after inserting \(\{d,e\}\) (not split) and then
\(\{b,d\}\) (split again, clique \(\{a,b,c,d\}\)).}
\label{fig:split-example}
\end{figure}

\section{Data structure}
The algorithm never needs to know \emph{which} vertex has which entry of the sorted degree
sequence: when \(u\)'s degree grows from \(d\) to \(d+1\), any entry equal to \(d\) may be
taken to represent \(u\). The structure (Ibarra p.\,40:18) is:
\begin{itemize}
\item a doubly linked list \(L\) of the degrees in non-increasing order (\code{sequence_}, a
  \code{std::list<Entry>} with \code{Entry\{degree, left\}});
\item for every degree value \(d\), pointers to the \emph{first} and \emph{last} list entries
  equal to \(d\) (\code{first_[d]}, \code{last_[d]}, \code{present_[d]});
\item the boundary \(j\), an iterator to the entry \(d_j\) (\code{boundary_}), and the sums
  \(\mathit{sum}_l=\sum_{i\le j}d_i\), \(\mathit{sum}_r=\sum_{i>j}d_i\);
\item the actual degree of every vertex (\code{degree_}), to know which degree classes to touch.
\end{itemize}
\op{Is-Split} is then the \(O(1)\) comparison \(\mathit{sum}_l=j(j-1)+\mathit{sum}_r\).
Preprocessing computes the degrees (\(O(n+m)\)), sorts them by counting sort (\(O(n)\)), builds
the list and the pointer arrays, and finds \(j\) and the sums in one pass.

\section{Updates in constant time}
\paragraph{Insert\((u,v)\).} Two entries must be incremented. If \(d(u)=d(v)=d\), increment the
first two entries equal to \(d\); otherwise increment the first entry equal to \(d(u)\) and the
first entry equal to \(d(v)\). Incrementing the \emph{first} entry of a class keeps the list
sorted (it becomes the last entry of class \(d+1\)), and only the class pointers of \(d\) and
\(d+1\) change. Then \(j\) can increase by at most one: before the insertion \(d_{j+1}\le j-1\),
so only \(d_{j+1}\) can start to satisfy \(d_{j+1}\ge j\), and Ibarra shows that \(j\) is
incremented at most once. The sums are adjusted in \(O(1)\).

\paragraph{Delete\((u,v)\).} Symmetric: decrement the \emph{last} entries of the classes, and
decrease \(j\) by one if \(d_j<j-1\).

\begin{lstlisting}[style=cpp,firstnumber=68,caption={\file{src/dynamic_split_ibarra.cpp}, lines 68--79: incrementing the first entry of a degree class.},label={lst:split-inc}]
void DynamicSplitGraph::increase(Iterator it) {
    const int d=it->degree;
    if(d+1>=static_cast<int>(present_.size()))
        throw std::invalid_argument("invalid duplicate insertion");
    const auto next=std::next(it);
    if(next!=sequence_.end()&&next->degree==d) first_[d]=next;
    else present_[d]=false;
    if(!present_[d+1]) {present_[d+1]=true;first_[d+1]=it;}
    last_[d+1]=it;
    ++it->degree;
    if(it->left)++sum_left_;else ++sum_right_;
}
\end{lstlisting}

\begin{lstlisting}[style=cpp,firstnumber=91,caption={\file{src/dynamic_split_ibarra.cpp}, lines 91--126: boundary update and \op{Insert}.},label={lst:split-insert}]
void DynamicSplitGraph::update_boundary_after_insert() {
    if(j_==static_cast<int>(degree_.size()))return;
    const auto next=std::next(boundary_);
    if(next!=sequence_.end()&&next->degree>=j_) {
        next->left=true;
        sum_left_+=next->degree;sum_right_-=next->degree;
        boundary_=next;++j_;
    }
}
void DynamicSplitGraph::update_boundary_after_delete() {
    if(j_<=1||boundary_->degree>=j_-1)return;
    auto old=boundary_;
    boundary_=std::prev(boundary_);
    old->left=false;
    sum_left_-=old->degree;sum_right_+=old->degree;
    --j_;
}
void DynamicSplitGraph::insert_edge(int u,int v) {
    validate(u,v);
    const int du=degree_[u],dv=degree_[v];
    // Ibarra Section 8: choose first one/two elements of degree classes.
    // Incrementing those entries preserves descending degree order.
    if(du==dv) {
        auto first=first_[du];
        const auto second=std::next(first);
        if(second==sequence_.end()||second->degree!=du)
            throw std::logic_error("invalid degree multiplicity");
        increase(first);increase(second);
    } else if(du>dv) {
        increase(first_[du]);increase(first_[dv]);
    } else {
        increase(first_[dv]);increase(first_[du]);
    }
    ++degree_[u];++degree_[v];
    update_boundary_after_insert();
}
\end{lstlisting}

\section{The four optimisation numbers}
Since chordal graphs are perfect, a split graph's chromatic number equals its clique number
\(j\). Running the same structure on the complement \(\bar G\) (whose degrees are \(n-1-d(v)\),
and where every insertion into \(G\) is a deletion from \(\bar G\) and vice versa) gives the
clique number of \(\bar G\), which is the independence number of \(G\), and hence also the
clique cover number. \code{DynamicSplitStatistics} maintains both structures and exposes
\code{clique_number()}, \code{chromatic_number()}, \code{independence_number()} and
\code{clique_cover_number()}, each \(O(1)\) (Ibarra p.\,40:18, last paragraph).

\section{Precondition}
Because only degrees are stored, the structure cannot itself detect an invalid update (inserting
an existing edge or deleting a non-edge); some such updates are detected indirectly (e.g.\
decrementing a zero degree), but not all. The caller must therefore pass only valid updates, as
stated in the header. This matches the paper, whose operations are edge insertions and deletions
on a graph whose edge set is known to the caller.

\section{Complexity and testing}
Preprocessing is \(O(n+m)\) and every update and query is \(O(1)\) worst-case: each update
changes two list entries, \(O(1)\) class pointers and at most one boundary position. The test
\file{tests/test_ibarra_split.cpp} performs 120{,}000 random update steps on graphs with up to
seven vertices; after every step it checks every maintained degree and compares \op{Is-Split}
with a from-scratch evaluation of \eqref{eq:hs} (sorting the degrees anew).
\file{tests/test_ibarra_split_stats.cpp} compares the clique and independence numbers (and
hence all four numbers) of every split graph reached with brute-force values obtained by
enumerating all vertex subsets, for \(n\le 7\).
%% ---- end expanded/10_split.tex ----
%% ---- begin expanded/11_preproc.tex ----
%% =====================================================================
\chapter{The A2 Preprocessing Claim}\label{ch:preprocessing}
%% =====================================================================
This chapter examines the one bound of Ibarra's summary table that the implementation does not
achieve: \(O(m+n\log n)\) preprocessing for A2. We first show that the paper's own text does not
support it, then that no implementation of the paper's method can achieve it, and finally that
achieving it by \emph{any} method, together with the claimed \(O(\log^2 n)\) insertion query,
would imply a breakthrough for Boolean matrix multiplication.

\begin{caution}[Provenance of this chapter]
The statements of the paper quoted in \cref{sec:preproc-paper} are Ibarra's. The analysis in
\cref{sec:preproc-explicit,sec:preproc-needed,sec:preproc-bmm} was carried out during this
project. It is not taken from a published source; in particular \cref{thm:bmm} is a
\emph{conditional} statement (it does not prove that the claim is impossible), and it should be
read as such.
\end{caution}

\section{What the paper says}\label{sec:preproc-paper}
\begin{table}[ht]
\centering\small
\caption[Statements of the paper about A2 preprocessing]{The paper's statements about the
preprocessing of A2 and A2\(^+\).}\label{tab:preproc-statements}
\begin{tabularx}{\textwidth}{@{}llX@{}}
\toprule
\textbf{Where} & \textbf{Page} & \textbf{Statement (paraphrased)}\\
\midrule
\S4 & 40:11 & A1: cliques and clique tree in \(O(m+n)\).\\
\S5 & 40:12 & Building the initial dynamic tree requires \(O(n\log n)\) time.\\
\S5.1 & 40:13 & \(W\) can be computed in \(O(n^3)\) by comparing the characteristic vectors of all pairs of cliques, or in \(O(mn)\) by starting from the empty graph and inserting the edges of \(G\) with \op{Insert} in reverse PEO order (``an improvement for sparse graphs'').\\
\S6.2 & 40:15 & A\(^+\) structure: preprocessing \(O(s_t+n^2+pn)\), i.e.\ \(O(mn+n^2)\).\\
\S7 & 40:17 & Table: A2 preprocessing \(O(m+n\log n)\); A2\(^+\) preprocessing \(O(mn+n^2)\).\\
\bottomrule
\end{tabularx}
\end{table}

The \S7 entry for A2 equals the A1 preprocessing plus the dynamic tree (\(O(m+n)+O(n\log n)\)),
but it does not include the cost of initialising \(W\), which \S5.1 itself gives as \(O(n^3)\)
or \(O(mn)\). Moreover, the \(O(mn)\) replay of \S5.1 starts from \(n\) singleton cliques, each
with a characteristic vector of length \(n\), which already costs \(\Theta(n^2)\) --- more than
\(O(mn)\) when \(m<n\). So the paper's own method gives \(O(n^2+mn)\), and the A2\(^+\) row
(\(O(mn+n^2)\)), which includes the A\(^+\) structure, is consistent with that.

\section{What the implementation does}
The implementation follows \S5.1 literally (\cref{sec:a2-preprocessing}) and therefore needs
\(\Theta(n^2+mn)\) time (\cref{prop:a2-pre}). The benchmark \code{mtp_a2_preprocessing_audit}
measures the cost on empty graphs (\(m=0\)), where the claimed bound would be
\(O(n\log n)\):

\begin{table}[ht]
\centering\small
\caption[A2 preprocessing on empty graphs]{A2 construction on the empty graph with \(n\)
vertices (\code{mtp_a2_preprocessing_audit}; machine of \cref{sec:exp-setup}).}\label{tab:preproc-empty}
\begin{tabular}{@{}rrrrr@{}}
\toprule
\(n\) & cliques & vector entries & \(W\) entries & time (ms)\\
\midrule
100 & 100 & 10{,}000 & 10{,}404 & 0.19\\
200 & 200 & 40{,}000 & 40{,}804 & 0.39\\
400 & 400 & 160{,}000 & 161{,}604 & 0.90\\
800 & 800 & 640{,}000 & 643{,}204 & 2.65\\
1600 & 1600 & 2{,}560{,}000 & 2{,}566{,}404 & 8.50\\
\bottomrule
\end{tabular}
\end{table}

The number of matrix entries is exactly \((n+2)^2\) and the time grows roughly fourfold when
\(n\) doubles at the larger sizes, as expected for \(\Theta(n^2)\) work.

\section{An explicit \texorpdfstring{\(W\)}{W} needs quadratic time}\label{sec:preproc-explicit}
\begin{proposition}\label{prop:w-lower}
Any algorithm that stores \(W\) as an explicit array with one written entry per pair of maximal
cliques needs \(\Omega(n^2)\) time on the star \(K_{1,n-1}\), although this graph has
\(m=n-1\), so that \(m+n\log n=O(n\log n)\).
\end{proposition}
\begin{proof}
The star has the \(n-1\) maximal cliques \(\{0,i\}\), and every two of them intersect in
\(\{0\}\). Hence all \((n-1)(n-2)\) off-diagonal entries of \(W\) equal 1. None of them is zero,
so even memory that is initialised to zero for free does not help: each entry must be written,
which takes \(\Omega(n^2)\) time.
\end{proof}

So no implementation of A2 \emph{as described in \S5.1} --- with an explicit matrix \(W\) ---
can be preprocessed in \(O(m+n\log n)\) time. Any method that meets that bound must avoid
storing \(W\) explicitly.

\section{Which intersections does \op{Insert-Query} need?}\label{sec:preproc-needed}
Could A2 avoid \(W\) by computing \(|K_x\cap K_y|\) from information on the tree path? Two
observations show that this is not straightforward.

\paragraph{The value is not a function of the path weights.}\label{sec:path-min-counterexample}
For the cliques \(\{0,1,2\},\{1,2,3\},\{2,3,4\}\) of \cref{fig:insert-query-no} (a path in the
clique tree), both tree edges have weight 2 but the end cliques intersect in one vertex. If the
end cliques were instead \(\{0,1,2\},\{1,2,3\},\{1,2,4\}\), the weights would again be \(2,2\) but
the end cliques would intersect in two vertices. So the path weights alone do not determine
\(|K_x\cap K_y|\), and in particular \(|K_x\cap K_y|\) cannot be replaced by the path minimum
(that replacement would answer \yes\ for every query).

\paragraph{Every pair of cliques can be queried.} For every two nodes \(x,y\) at distance at
least two in a clique tree, choose \(u\in K_x\setminus K_{x'}\) and \(v\in K_y\setminus K_{y'}\),
where \(x'\) and \(y'\) are the neighbours of \(x\) and \(y\) on the \(x\)--\(y\) path; then
\(uv\notin E\) and \((x,y)\) is the closest pair of \op{Insert-Query}\((u,v)\). (This was checked
for all 565 such pairs occurring in the 406 test graphs of
\file{audit/verify_chordal_identities.py}, check C6.) Hence a
structure answering every insertion query must, in effect, be able to compare
\(|K_x\cap K_y|\) with the path minimum for \emph{every} pair of cliques.

\section{A conditional barrier via Boolean matrix multiplication}\label{sec:preproc-bmm}
The Boolean product \(C=AB\) of two Boolean \(u\times u\) matrices is
\(C_{ij}=\bigvee_k(A_{ik}\wedge B_{kj})\). Fast algebraic matrix multiplication provides subcubic bounds. A 2026 preprint
reports \(\omega\le 9/4\) for matrix multiplication over the complex numbers
~\cite{openai2026,shpilka2026}; this concerns arithmetic complexity over a field,
not an \(O(u^2\polylog u)\) Boolean-matrix multiplication algorithm.
No general \(O(u^2\polylog u)\) Boolean-matrix multiplication algorithm is
known. Thus the reduction below remains a conditional barrier rather than an
impossibility theorem.

\paragraph{The gadget.} Given \(A,B\), build the graph \(G_{A,B}\) on \(n=5u\) vertices
(\cref{fig:bmm}):
\begin{itemize}
\item three groups \(E=\{e_1,\dots,e_u\}\), \(D=\{d_1,\dots,d_u\}\), \(F=\{f_1,\dots,f_u\}\)
  that together form one clique \(L\) of \(3u\) vertices;
\item vertices \(p_1,\dots,p_u\), where \(p_i\) is adjacent to
  \(R_i=\{e_k: A_{ik}=1\}\) padded with vertices of \(D\) to exactly \(u\) neighbours;
\item vertices \(q_1,\dots,q_u\), where \(q_j\) is adjacent to
  \(Q_j=\{e_k: B_{kj}=0\}\cup D\), padded with vertices of \(F\) to exactly \(2u\) neighbours.
\end{itemize}
\(G_{A,B}\) is a split graph (clique \(L\), independent set \(\{p_i\}\cup\{q_j\}\)), hence
chordal, with \(m=\binom{3u}{2}+u\cdot u+u\cdot 2u=7.5u^2-1.5u\) edges, and it can be built in
\(O(u^2)\) time.

\begin{figure}[ht]
\centering
\begin{tikzpicture}
\draw[rounded corners=6pt,draw=accent,fill=accentlight!60,thick] (-0.3,-0.6) rectangle (6.3,0.6);
\node[font=\small] at (1.0,0) {\(E\) (\(u\))};
\node[font=\small] at (3.0,0) {\(D\) (\(u\))};
\node[font=\small] at (5.0,0) {\(F\) (\(u\))};
\draw[black!40] (2.0,-0.6)--(2.0,0.6) (4.0,-0.6)--(4.0,0.6);
\node[font=\small\bfseries,accent] at (3.0,0.95) {clique \(L=E\cup D\cup F\)};
\node[vtx] (p) at (1.2,-2.2) {\(p_i\)};
\node[vtx] (q) at (4.6,-2.2) {\(q_j\)};
\draw[gedge,draw=orange!80!black] (p)--(0.7,-0.6) (p)--(1.4,-0.6) (p)--(2.4,-0.6);
\node[font=\scriptsize,align=center,text=orange!60!black] at (-0.4,-1.7) {\(\{e_k:A_{ik}=1\}\)\\+ pad from \(D\)};
\draw[gedge,draw=okgreen!70!black] (q)--(1.7,-0.6) (q)--(2.6,-0.6) (q)--(3.5,-0.6) (q)--(4.6,-0.6);
\node[font=\scriptsize,align=center,text=okgreen!50!black] at (6.4,-1.7) {\(\{e_k:B_{kj}=0\}\cup D\)\\+ pad from \(F\)};
\node[font=\small,align=center] at (3.0,-3.3) {\(G_{A,B}+p_iq_j\) is chordal \(\iff\) \(C_{ij}=0\)};
\end{tikzpicture}
\caption[Reduction from Boolean matrix multiplication]{The graph \(G_{A,B}\): \(p_i\) has \(u\)
neighbours and \(q_j\) has \(2u\) neighbours in the clique \(L\). Each insertion query
\((p_i,q_j)\) decides one entry of the Boolean product.}
\label{fig:bmm}
\end{figure}

\begin{lemma}\label{lem:bmm-gadget}
\(G_{A,B}+p_iq_j\) is chordal if and only if \(C_{ij}=0\).
\end{lemma}
\begin{proof}
The maximal cliques are \(L\), \(P_i=\{p_i\}\cup R_i\) and \(Q'_j=\{q_j\}\cup Q_j\), and the star
with centre \(L\) is a clique tree (each \(P_i\) and \(Q'_j\) has its heaviest possible edge,
of weight \(u\) resp.\ \(2u\), to \(L\)). The closest pair for \((p_i,q_j)\) is
\((P_i,Q'_j)\), the path is \(P_i\!-\!L\!-\!Q'_j\) with minimum weight \(u\), and
\[ |P_i\cap Q'_j| = |\{k: A_{ik}=1, B_{kj}=0\}| + (u-|\{k:A_{ik}=1\}|) = u - |\{k: A_{ik}=B_{kj}=1\}| \]
(the padding vertices of \(p_i\) lie in \(D\subseteq Q_j\)). By \cref{thm:insert6} the insertion
is allowed iff this equals \(u\), i.e.\ iff no \(k\) has \(A_{ik}=B_{kj}=1\).
\end{proof}

\begin{theorem}[Conditional barrier]\label{thm:bmm}
Suppose some data structure, for every chordal graph with \(n\) vertices and \(m\) edges, can be
built in \(O(m+n\log n)\) time and then answers \op{Insert-Query} in \(O(\log^2 n)\) worst-case
time. Then the Boolean product of two \(u\times u\) matrices can be computed in
\(O(u^2\log^2 u)\) time.
\end{theorem}
\begin{proof}
Build \(G_{A,B}\) (\(O(u^2)\)), build the structure (\(O(m+n\log n)=O(u^2)\)), and ask the
\(u^2\) queries \((p_i,q_j)\) (\(O(u^2\log^2 u)\)); by \cref{lem:bmm-gadget} the answers are the
entries of \(C\).
\end{proof}

Since Ibarra's queries do not modify the graph, the statement applies to the \S7 claims for A2
taken together (preprocessing \emph{and} \op{Insert-Query}). It shows that meeting both would
yield near-quadratic Boolean matrix multiplication, far beyond the best known algorithms. We
therefore regard the \S7 preprocessing entry for A2 as not achievable with known techniques,
and the implementation keeps the method that \S5.1 actually describes. The reduction was
verified on 150 random instances with 1{,}776 queries by \file{audit/verify_chordal_identities.py}
(check C7).

\section{Research prototypes that were not merged}
Two alternatives were explored on separate branches and deliberately \emph{not} merged into the
frozen implementation, because neither achieves the claimed bound:
\begin{itemize}
\item \textbf{R1} (branch \code{research/a2-r1-constructor-audit}): build the clique tree
  directly with Blair--Peyton, compute \(W\) from the identity
  \(W[x][y]=\sum_v[v\in K_x][v\in K_y]\) by incrementing \(W[x][y]\) for all pairs
  \(x,y\in T_v\) for each vertex \(v\), and build the dynamic tree with
  \code{build_linear_forest}. This avoids the \(O(mn)\) replay, but costs
  \(O(n^2+\sum_v|T_v|^2)\), which is still \(\Theta(n^2)\) (e.g.\ on the star).
\item \textbf{R2} (design note only): drop \(W\) and compute \(|K_x\cap K_y|\) on demand during
  \op{Insert-Query}. Preprocessing becomes fast, but the query becomes \(O(n)\) in the worst
  case, which is no longer A2.
\end{itemize}

\begin{summarybox}[Conclusion of this chapter]
The implementation achieves every A2 operation bound of the paper, and \(\Theta(n^2+mn)\)
preprocessing, which is what the paper's own construction in \S5.1 gives. The
\(O(m+n\log n)\) entry of the summary table is \emph{not achieved}: it is impossible for any
explicit-\(W\) implementation (\cref{prop:w-lower}), and achieving it together with
\(O(\log^2 n)\) queries by any method would imply \(O(u^2\log^2 u)\) Boolean matrix
multiplication (\cref{thm:bmm}).
\end{summarybox}
%% ---- end expanded/11_preproc.tex ----
%% ---- begin expanded/12_verif.tex ----
%% =====================================================================
\chapter{Verification}\label{ch:verification}
%% =====================================================================
An implementation of a dynamic graph algorithm can be wrong in ways that small hand-made
examples do not reveal: a clique tree can lose the induced subtree property after a particular
sequence of updates, an edge weight can drift, or a dynamic-tree invariant can break only after
many operations. The project therefore relies on systematic testing against \emph{independent
oracles} --- simpler programs that compute the same answer in a completely different way --- on
exhaustive families of small graphs and on long random sequences of legal updates. This chapter
describes the oracles, the test suite and its results.

\section{Oracles}
\begin{description}
\item[Chordality oracle.] For a graph \(H\), compute an MCS ordering and test whether it is a
  PEO (\cref{thm:mcs,lem:peo-test}). Every \op{Insert-Query} and \op{Delete-Query} answer is
  compared with this oracle applied to \(G+uv\) or \(G-uv\), which are rebuilt from scratch.
\item[Maximal-clique oracle.] \code{extract_maximal_cliques} (\cref{sec:cliques-code}); after
  every update the set of cliques stored in the clique tree must equal the oracle's set.
\item[Clique-tree checker.] For the maintained tree: exactly \(c-1\) edges; every edge weight
  equals the true intersection size; dummy edges are exactly the weight-0 edges; and, for every
  vertex \(v\), the nodes containing \(v\) induce a connected subtree (checked by BFS
  restricted to those nodes).
\item[Spanning-tree oracles.] Kruskal on the explicit \(W_G\), and exhaustive enumeration of
  all spanning trees by Prüfer sequences for tiny graphs (\cref{sec:baseline}).
\item[Differential oracle.] A2, A2\(^+\), A1\(^+\) and the indexed variant are compared query
  by query and update by update with A1.
\item[Dynamic-tree oracle.] A naive forest (explicit adjacency, paths by search) against which
  \code{root}, \code{kth_on_path}, \code{path_length} and \code{min_on_path} are compared for
  \emph{all} pairs of nodes after every link or cut.
\item[Structural checkers.] \code{check_invariants} of the ST forest and of every globally
  biased tree (\cref{sec:st83,sec:bst85}), run after every operation in the tests.
\end{description}

\section{Exhaustive and randomised testing}
\paragraph{All small chordal graphs.} Several tests enumerate \emph{every} graph on
\(n\le 6\) labelled vertices (all \(2^{\binom n2}\) edge subsets) and keep the chordal ones.
There are \(1+2+8+61+822+18{,}154=19{,}048\) labelled chordal graphs with 1 to 6 vertices, and
the static algorithms are checked on all of them. The dynamic tests enumerate all chordal
graphs with \(n\le 5\) and, from each one, perform random sequences of legal updates; before
every update they compare the answers of \op{Insert-Query} / \op{Delete-Query} for
\emph{every} pair of vertices with the chordality oracle.

\paragraph{Random graphs and long sequences.} Random chordal graphs are generated by adding
vertices one at a time and joining each new vertex to a random subset of a random maximal clique
(so the new vertex is simplicial and the graph stays chordal). On graphs with 6--12 vertices the
tests then run 20--25 legal updates each, again comparing every possible query with the oracle
and validating the full clique tree after each update.

\section{The test suite}
\Cref{tab:tests} lists all 23 CTest entries of the repository. On the machine of
\cref{sec:exp-setup} the complete suite runs in about 45 seconds in the Release configuration,
and all 23 tests pass (last run: 23 November 2026, commit \code{651b236}).

\begin{longtable}{@{}>{\raggedright\arraybackslash}p{0.33\textwidth}>{\raggedright\arraybackslash}p{0.62\textwidth}@{}}
\caption[The CTest suite]{The 23 CTest entries (\code{ctest --test-dir build}). The numbers
are those printed by the test programs.}\label{tab:tests}\\
\toprule
\textbf{Test (program)} & \textbf{What is checked}\\
\midrule
\endfirsthead
\toprule
\textbf{Test (program)} & \textbf{What is checked}\\
\midrule
\endhead
\bottomrule
\endfoot
\code{chordal_algorithms} (\code{mtp_tests}) & MCS, PEO test and maximal cliques on the sample, edge cases, all 19{,}048 chordal graphs with \(n\le6\), and 600 random chordal graphs with 7--12 vertices\\
\code{clique_tree_algorithms} & baseline \(W_G\) + Kruskal: running intersection on all 19{,}048 graphs; maximum weight against Prüfer enumeration for \(n\le5\); 600 random graphs\\
\code{linear_clique_tree} & Blair--Peyton tree: all 19{,}048 graphs, 800 random graphs, a large star and a large clique; weight equals an independent maximum spanning tree\\
\code{ibarra_first_dynamic_implementation} & dense A1: oracle comparison of every query and full tree validation after every update (all chordal graphs \(n\le5\) + 200 random sequences)\\
\code{ibarra_a1_linear_preprocessing} & sparse A1: the same checks\\
\code{ibarra_a1_sparse_fast_index} & sparse A1\(^+\): the same checks with the \S6 index\\
\code{ibarra_a1_sparse_adversarial} & sparse A1 versus dense A1: stored memberships \(\le n+m\); exactly \(2(n-1)\) on stars and \(n\) on empty graphs (not \(\Theta(n^2)\)); repeated deletions/insertions inside a large clique with node recycling; the bulk loader (adjacency, duplicate rejection)\\
\code{ibarra_split_graph_implementation} & split graphs: 120{,}000 random update steps against a from-scratch Hammer--Simeone evaluation\\
\code{ibarra_split_four_numbers} & clique, chromatic, independence and clique-cover numbers against brute force (\(n\le7\))\\
\code{deterministic_graph_membership} & \code{Graph} membership under 12{,}000 random updates and queries\\
\code{ibarra_section5_actual_ST_dynamic_tree} & A2 versus A1: all chordal graphs \(n\le5\) (with and without A2\(^+\)), 60 random sequences, tree and ST invariants after every update, repeated two-hub insert/delete cycles\\
\code{sleator_tarjan_1983_dynamic_forest} & ST forest against the naive forest: \(n=3..12\), 30 trials \(\times\) 120 random link/cut steps, all-pairs path queries after every step\\
\code{ibarra_a2_linear_insert} & double-subsumption insertions on seven two-hub graphs (\(n=7..259\)), three insert/delete cycles each: zero individual ST calls, exact linear rebuild, A1 oracle\\
\code{st1983_linear_forest_construction} & \code{build_linear_forest} on stars, paths, random trees and double stars, followed by dynamic link/cut/evert and BFS path/min oracles\\
\code{bst85_linear_global_bias} & \code{build_linear}: 840 weighted configurations; rank and global-bias invariants, order, path minima, step counts\\
\code{published_globally_biased_binary_tree} & global join and split: 892{,}800 verified operations\\
\code{ibarra_indexed_path_variant} & indexed comparison variant versus A1\\
\code{ibarra_fast_delete_query} & dense A1\(^+\) versus oracle\\
\code{ibarra_fast_large_sparse_dense} & modified Yellin index on stars (\(n=128,256\)) and complete graphs (\(n=100,200\)) under legal updates\\
\code{ibarra_dynamic_sample_cli} & \code{mtp_dynamic_cli} on the sample files\\
\code{ibarra_indexed_sample_cli} & \code{mtp_indexed_cli} on the sample files\\
\code{ibarra_a2_sample_cli} & \code{mtp_a2_cli} on the sample files\\
\code{ibarra_a2_fast_delete_sample_cli} & \code{mtp_a2_cli --fast-delete} on the sample files\\
\end{longtable}

\section{Testing the dynamic tree}\label{sec:tests-st}
The Sleator--Tarjan forest is the most intricate component, so it is tested at three levels:
\begin{enumerate}
\item the globally biased trees alone (join, split, linear build) against simple sequence
  models, with the rank invariants (B1)--(B4) checked after every operation;
\item the forest alone against a naive forest, with \emph{all-pairs} path queries after every
  random link or cut, and with \code{check_invariants} verifying that every vertex weight equals
  one plus the weights of its dashed paths, that every dashed edge has exactly one path-set
  entry, and that every path set is sorted by decreasing weight;
\item A2 as a whole against A1.
\end{enumerate}

\section{Memory safety and continuous integration}\label{sec:ci}
The repository has five GitHub Actions workflows (\file{.github/workflows/}):
\begin{itemize}
\item \file{ci.yml} (every push and pull request): Release build with benchmarks, the full
  CTest suite, the sample CLI and the benchmark programs;
\item \file{a1-sanitizers.yml} (pushes that touch the A1 or \code{Graph} sources or tests): the
  A1 tests built with AddressSanitizer and UndefinedBehaviorSanitizer;
\item \file{a2-linear-research.yml} (pushes to \code{main} and the research branch) and
  \file{a2-research.yml} (research branch, or on demand): the dynamic-tree and A2 tests, the
  operation-count audits, and a sanitizer build of the dynamic-tree tests;
\item \file{final-mtp-audit.yml} (changes to the audit scripts or the earlier report, or on
  demand): checks that the algorithm sources are unchanged since the frozen commit, runs the
  Python audit scripts and builds the earlier LaTeX report.
\end{itemize}
(The sanitizer builds use smaller test parameters, selected with \code{__SANITIZE_ADDRESS__},
because instrumented code is much slower.)

\section{Independent audit scripts}\label{sec:audit-scripts}
Two Python scripts in \file{audit/} re-derive the key mathematical facts with NetworkX, without
using any of the C++ code:
\begin{itemize}
\item \file{verify_chordal_identities.py} checks, on 406 random and named chordal graphs:
  \(\sum_K|K|\le n+m\) (C1); that \(|K_x\cap K_y|\) equals the intersection of the separators
  on the tree path (C2); \(\sum_v|T_v|^2=\sum_{x,y}W[x][y]\) (C3); that Ibarra's insertion
  criterion agrees with chordality of \(G+uv\) for all 4{,}466 non-edges (C4, C5); that every
  pair of nodes at distance \(\ge2\) is the closest pair of some query (C6, 565 pairs); the BMM
  reduction of \cref{sec:preproc-bmm} on 150 instances with 1{,}776 queries (C7); and a
  common-neighbourhood form of the insertion criterion (C8). Result of the run on
  23 November 2026: \emph{all checks passed}.
\item \file{verify_w_updates.py} applies the formulas of \cref{lem:w-delete,lem:w-insert} to
  1{,}171 insertions and 5{,}895 deletions and compares them with directly computed
  intersections. Result: \emph{all formulas correct}.
\end{itemize}

\section{Limits of the verification}
Testing cannot prove the absence of errors. The evidence is strongest for small graphs, where
it is exhaustive, and for the logical correctness of the clique-tree maintenance, which is
checked against oracles that share no code with the implementation. The \emph{complexity}
claims are established by the proofs in \cref{ch:static,ch:a1,ch:a2,ch:fastdelete,ch:split},
not by the tests; for A2 they rely on the published theorems of ST83 and BST85 together with
the invariant checks described above, and there is no machine-checked proof that the C++ code
maintains every invariant on every input.
%% ---- end expanded/12_verif.tex ----
%% ---- begin expanded/13_exp.tex ----
%% =====================================================================
\chapter{Experiments}\label{ch:experiments}
%% =====================================================================
\definecolor{series1}{HTML}{2A78D6}
\definecolor{series2}{HTML}{EB6834}
\definecolor{refgray}{HTML}{8A8984}
\pgfplotsset{
  mtpaxis/.style={
    width=0.47\textwidth,height=5.4cm,
    axis lines=left,
    tick style={black!45},
    axis line style={black!55},
    grid=major,major grid style={black!10},
    tick label style={font=\scriptsize},
    label style={font=\small},
    title style={font=\small\bfseries},
    legend style={font=\scriptsize,draw=none,fill=white,fill opacity=0.85,text opacity=1},
    legend cell align=left,
  },
  s1/.style={color=series1,line width=1.4pt,mark=*,mark size=2.2pt,mark options={fill=series1,draw=white,line width=0.6pt}},
  s2/.style={color=series2,line width=1.4pt,mark=square*,mark size=2.1pt,mark options={fill=series2,draw=white,line width=0.6pt}},
  refline/.style={color=refgray,dashed,line width=0.8pt},
}

The experiments in this chapter illustrate the behaviour of the implementation; they are
\emph{not} evidence for the asymptotic bounds, which are proved in the previous chapters.
Wall-clock times on a shared virtual machine are noisy, and constant factors dominate at the
sizes that can be measured quickly, so we mostly look at how times \emph{grow} with \(n\).

\section{Setup}\label{sec:exp-setup}
All numbers were produced on 23 November 2026 by the benchmark programs of the repository
(\cref{tab:benchmarks}), built with \code{-DCMAKE_BUILD_TYPE=Release
-DMTP_BUILD_BENCHMARKS=ON} (flags \code{-O3 -DNDEBUG}) with g++~13.3.0 and CMake~3.28.3 on a
virtual machine with 2 vCPUs of an Intel Xeon processor at 2.10\,GHz, running Ubuntu~24.04. Each
reported operation time is the average over the repetitions stated; the repository also
archives earlier runs of the same programs in \file{benchmarks/results/}.
The seven first-run CSV files supporting the tables and figures here, a second
measurement run, the benchmark script and machine/toolchain information are now
archived under \file{benchmarks/results/2026-11-23/}. The figures and tables
use the first-run files. These user-provided experimental logs were checked for
consistency with the printed numbers; no independent rerun on that VM is claimed.
The second run reproduced deterministic count columns but displayed normal
timing variation. Construction timings are individual measurements, not
means of multiple trials; operation timings average the repetitions shown.
The machine was a shared two-vCPU VM without CPU pinning or frequency control,
so absolute timing comparisons require care.

\begin{table}[ht]
\centering\small
\caption[Benchmark programs]{Benchmark programs (\file{benchmarks/}).}\label{tab:benchmarks}
\begin{tabularx}{\textwidth}{@{}lX@{}}
\toprule
\textbf{Program} & \textbf{Measurement}\\
\midrule
\code{mtp_benchmarks} & sparse A1 on stars \(K_{1,n-1}\): construction, and 200 insert/delete cycles of a leaf--leaf edge\\
\code{mtp_benchmark_a2_st} & A1 (dense) versus A2 on paths \(P_n\): construction, and 700 repetitions of the rejected query \op{Insert-Query}\((0,n-1)\)\\
\code{mtp_a2_preprocessing_audit} & A2 construction on empty graphs (\cref{tab:preproc-empty})\\
\code{mtp_a2_two_hub_stress} & double-subsumption insertion on two-hub trees (\cref{tab:two-hub})\\
\code{mtp_a2_operation_audit} & dynamic-tree calls of A2 insertions on stars and double stars\\
\code{mtp_benchmark_indexed} & A1 (dense) versus the indexed comparison variant on paths\\
\bottomrule
\end{tabularx}
\end{table}

\section{A1 on stars}
A star is the extreme case for the dense representation (\(c=n-1\) cliques) and for clique-tree
degrees (the centre's cliques form a tree in which one node can have degree \(\Theta(n)\)).
\Cref{tab:a1-star,fig:a1-star} show the sparse A1.

\begin{table}[ht]
\centering\small
\caption[A1 on stars]{Sparse A1 on stars (\code{mtp_benchmarks}); times in microseconds, except
construction in milliseconds. ``Memberships'' is \(\sum_K|K|\).}\label{tab:a1-star}
\begin{tabular}{@{}rrrrrrrr@{}}
\toprule
\(n\) & \(m\) & memberships & build (ms) & \op{Ins-Q} (\(\mu\)s) & \op{Insert} (\(\mu\)s) & \op{Del-Q} (\(\mu\)s) & \op{Delete} (\(\mu\)s)\\
\midrule
50 & 49 & 98 & 0.033 & 0.79 & 1.79 & 0.05 & 0.33\\
100 & 99 & 198 & 0.046 & 1.17 & 2.81 & 0.05 & 0.49\\
200 & 199 & 398 & 0.073 & 2.20 & 5.09 & 0.06 & 0.40\\
400 & 399 & 798 & 0.144 & 3.76 & 7.92 & 0.05 & 0.37\\
800 & 799 & 1598 & 0.286 & 5.93 & 12.91 & 0.06 & 0.47\\
1600 & 1599 & 3198 & 0.538 & 10.30 & 21.69 & 0.24 & 0.48\\
\bottomrule
\end{tabular}
\end{table}

\begin{figure}[ht]
\centering
\begin{tikzpicture}
\begin{loglogaxis}[mtpaxis,title={Construction time},xlabel={\(n\)},ylabel={time (ms)},
  log basis x=2,xtick={50,100,200,400,800,1600},xticklabels={50,100,200,400,800,1600},
  ytick={0.03,0.1,0.3,0.6},yticklabels={0.03,0.1,0.3,0.6}]
\addplot[s1] coordinates {(50,0.03304)(100,0.04622)(200,0.07301)(400,0.144)(800,0.28616)(1600,0.53819)};
\addplot[refline,domain=200:1600,samples=2] {0.0727*x/200};
\node[font=\scriptsize,text=refgray,anchor=north west] at (axis cs:300,0.085) {slope 1 (linear)};
\end{loglogaxis}
\end{tikzpicture}\hfill
\begin{tikzpicture}
\begin{loglogaxis}[mtpaxis,title={Insertion query and insertion},xlabel={\(n\)},ylabel={time (\(\mu\)s)},
  log basis x=2,xtick={50,100,200,400,800,1600},xticklabels={50,100,200,400,800,1600},
  ytick={1,3,10,20},yticklabels={1,3,10,20},legend pos=north west]
\addplot[s1] coordinates {(50,0.79)(100,1.17)(200,2.20)(400,3.76)(800,5.93)(1600,10.30)};
\addlegendentry{\op{Insert-Query}}
\addplot[s2] coordinates {(50,1.79)(100,2.81)(200,5.09)(400,7.92)(800,12.91)(1600,21.69)};
\addlegendentry{\op{Insert}}
\end{loglogaxis}
\end{tikzpicture}
\caption[A1 on stars: growth of construction and insertion time]{Sparse A1 on stars
(log--log axes). The observed construction and operation timings are consistent with the worst-case \(O(n+m)\) and \(O(n)\) bounds proved by the data-structure analysis; five timing points do not establish asymptotic worst-case behaviour.}
\label{fig:a1-star}
\end{figure}

The number of stored memberships is exactly \(2(n-1)=n+m-1\), confirming that the sparse
representation stays within \(\sum_K|K|\le n+m\), while the dense representation would store
\(n(n-1)\) vector entries (2{,}558{,}400 for \(n=1600\)).

\section{Insert-Query: A1 versus A2}\label{sec:exp-a1-a2}
On the path \(P_n\) the query \op{Insert-Query}\((0,n-1)\) is the worst case for A1: the BFS
traverses the whole clique tree. A2 answers it with two binary searches on the dynamic tree.

\begin{table}[ht]
\centering\small
\caption[Insert-Query on paths: dense A1 versus A2]{Dense A1 versus A2, \op{Insert-Query}\((0,n-1)\) on paths, average of
700 repetitions (\code{mtp_benchmark_a2_st}); construction times for reference.}\label{tab:a1-a2}
\begin{tabular}{@{}rrrrr@{}}
\toprule
\(n\) & dense A1 construct (ms) & A2 construct (ms) & dense A1 query (\(\mu\)s) & A2 query (\(\mu\)s)\\
\midrule
64 & 0.043 & 3.65 & 1.08 & 53.95\\
128 & 0.057 & 9.92 & 2.21 & 74.99\\
256 & 0.131 & 31.55 & 4.23 & 85.30\\
512 & 0.221 & 64.35 & 8.95 & 110.42\\
1024 & 0.615 & 175.29 & 16.96 & 125.02\\
\bottomrule
\end{tabular}
\end{table}

\begin{figure}[ht]
\centering
\begin{tikzpicture}
\begin{loglogaxis}[mtpaxis,width=0.62\textwidth,height=6cm,title={\op{Insert-Query}\((0,n-1)\) on paths},
  xlabel={\(n\)},ylabel={time per query (\(\mu\)s)},log basis x=2,
  xtick={64,128,256,512,1024},xticklabels={64,128,256,512,1024},
  legend style={at={(1.03,0.5)},anchor=west}]
\addplot[s1] coordinates {(64,1.07927)(128,2.21047)(256,4.23004)(512,8.94706)(1024,16.95669)};
\addlegendentry{Dense A1 (BFS, \(O(n)\))}
\addplot[s2] coordinates {(64,53.9549)(128,74.9954)(256,85.29872)(512,110.41675)(1024,125.02314)};
\addlegendentry{A2 (dynamic tree, \(O(\log^2 n)\))}
\end{loglogaxis}
\end{tikzpicture}
\caption[Insert-Query growth: dense A1 versus A2]{Growth of the insertion query (dense A1 baseline versus A2). From \(n=64\) to
\(n=1024\) (16\(\times\)), A1's time grows by a factor of about 16 (linear), A2's by about 2.3
(polylogarithmic). At these sizes A1 is still faster in absolute terms, because each of A2's
\(O(\log^2 n)\) steps is an expose/conceal sequence on biased trees with a large constant
factor.}
\label{fig:a1-a2}
\end{figure}

This comparison deliberately uses the dense \code{DynamicChordalGraph} A1 baseline, not the sparse A1 production CLI. It illustrates the dynamic-tree query complexity but does not compare the two production loading routes. A2 trades a much larger constant for a slower growth rate on this path family.
(The A2 construction times include the \(\Theta(n^2+mn)\) preprocessing discussed in
\cref{ch:preprocessing}.)

\section{A2 preprocessing on empty graphs}
\Cref{fig:preproc} plots \cref{tab:preproc-empty}. On the empty graph Ibarra's table would
predict \(O(n\log n)\); the measured growth approaches the quadratic reference line, matching
\cref{prop:a2-pre}.

\begin{figure}[ht]
\centering
\begin{tikzpicture}
\begin{loglogaxis}[mtpaxis,title={A2 construction, \(m=0\)},xlabel={\(n\)},ylabel={time (ms)},
  log basis x=2,xtick={100,200,400,800,1600},xticklabels={100,200,400,800,1600},
  ytick={0.2,0.5,1,2,5,10},yticklabels={0.2,0.5,1,2,5,10}]
\addplot[s1] coordinates {(100,0.186137)(200,0.391003)(400,0.90046)(800,2.64715)(1600,8.49934)};
\addplot[refline,domain=400:1600,samples=2] {0.9*(x/400)^2};
\node[font=\scriptsize,text=refgray,anchor=south east] at (axis cs:900,4.0) {slope 2};
\end{loglogaxis}
\end{tikzpicture}\hfill
\begin{tikzpicture}
\begin{axis}[mtpaxis,title={Entries of \(W\)},xlabel={\(n\)},ylabel={entries (millions)},
  xtick={100,400,800,1200,1600},ymin=0]
\addplot[s1] coordinates {(100,0.010404)(200,0.040804)(400,0.161604)(800,0.643204)(1600,2.566404)};
\end{axis}
\end{tikzpicture}
\caption[A2 preprocessing on empty graphs]{A2 construction on empty graphs. Left: time
(log--log) with a dashed slope-2 reference. Right: the matrix \(W\) has exactly \((n+2)^2\)
entries.}
\label{fig:preproc}
\end{figure}

\section{Linear rebuild in double-subsumption insertions}
The test \code{ibarra_a2_linear_insert} reports the number of elementary steps of
\code{build_linear_forest} for double-subsumption insertions on two-hub graphs of growing size.

\begin{figure}[ht]
\centering
\begin{tikzpicture}
\begin{axis}[mtpaxis,width=0.6\textwidth,height=5.6cm,title={Rebuild steps per double-subsumption insertion},
  xlabel={\(n\)},ylabel={steps},xmin=0,ymin=0]
\addplot[s1] coordinates {(7,80)(11,132)(19,240)(35,460)(67,904)(131,1796)(259,3584)};
\addplot[refline,domain=0:260,samples=2] {13.84*x};
\node[font=\scriptsize,text=refgray,anchor=south east] at (axis cs:250,3700) {\(13.84\,n\)};
\end{axis}
\end{tikzpicture}
\caption[Linear dynamic-tree rebuild]{Steps of the Sleator--Tarjan rebuild against \(n\)
(two-hub graphs, \(n=7,\dots,259\)). The points lie on a straight line through the origin
(\(\approx 11.4n\) to \(13.8n\)); no individual link/cut calls are made.}
\label{fig:rebuild}
\end{figure}

\section{The indexed comparison variant}
For completeness, \cref{tab:indexed} compares dense A1 with the indexed variant of
\cref{sec:indexed} on paths. Its query is very fast (a static heavy-path index), but its
construction and every update rebuild the index, which is why it is not a substitute for the
Sleator--Tarjan tree required by the paper.

\begin{table}[ht]
\centering\small
\caption[A1 versus the indexed variant]{Dense A1 versus the indexed variant on paths
(\code{mtp_benchmark_indexed}, 2000 repetitions of \op{Insert-Query}\((0,n-1)\)).}\label{tab:indexed}
\begin{tabular}{@{}rrrrr@{}}
\toprule
\(n\) & dense A1 build (ms) & indexed build (ms) & dense A1 query (\(\mu\)s) & indexed query (\(\mu\)s)\\
\midrule
100 & 0.093 & 1.41 & 1.64 & 0.11\\
200 & 0.099 & 5.28 & 3.34 & 0.14\\
400 & 0.186 & 20.10 & 6.68 & 0.16\\
800 & 0.391 & 90.84 & 15.45 & 0.32\\
1600 & 1.823 & 343.88 & 31.88 & 0.19\\
\bottomrule
\end{tabular}
\end{table}

\section{Discussion}
\begin{itemize}
\item The sparse A1 is the practical choice for all operations at the sizes tested; its
  construction is linear, as proved.
\item A2's insertion query grows far more slowly than A1's, as the theory predicts, but its
  constant factor is large; the break-even point lies beyond the sizes measured here.
\item A2's construction is quadratic even for empty graphs, which is the experimental face of
  \cref{ch:preprocessing}.
\item The double-subsumption rebuild performs work proportional to \(n\) and no individual
  dynamic-tree operations, consistent with \cref{lem:build-forest}.
\end{itemize}
%% ---- end expanded/13_exp.tex ----
%% ---- begin expanded/14_concl.tex ----
%% =====================================================================
\chapter{Conclusions and Future Work}\label{ch:conclusion}
%% =====================================================================
\section{Summary}
This project implemented, in about 4{,}250 lines of C++17, every algorithm of Ibarra's paper
\emph{Fully Dynamic Algorithms for Chordal Graphs and Split Graphs}: the static clique-tree
preprocessing, the two implementations A1 and A2 of the dynamic chordal graph algorithm, their
faster deletion-query variants A1\(^+\) and A2\(^+\), and the constant-time dynamic split graph
algorithm. Along the way it implemented two substantial published data structures from scratch
--- the worst-case dynamic trees of Sleator and Tarjan and the globally biased binary trees of
Bent, Sleator and Tarjan --- and it verified the whole system with exhaustive and randomised
tests against independent oracles. \Cref{tab:final-status} gives the final status of every
claim.

\begin{table}[ht]
\centering\small
\caption[Final status of every claim of the paper]{Final status of every claim of Ibarra's paper
in this implementation. All bounds are worst-case.}\label{tab:final-status}
\begin{tabularx}{\textwidth}{@{}>{\raggedright\arraybackslash}p{3.3cm}lX>{\raggedright\arraybackslash}p{2.3cm}@{}}
\toprule
\textbf{Claim (paper)} & \textbf{Bound} & \textbf{Where in this report} & \textbf{Status}\\
\midrule
Static clique tree (\S2, \S4) & \(O(m+n)\) & \cref{ch:static}, \cref{tab:static-cost} & \ok\\
A1 preprocessing & \(O(m+n)\) & \cref{prop:a1-pre} (sparse representation) & \ok\\
A1 four operations & \(O(n)\) & \cref{thm:a1-bounds} & \ok\\
A2 \op{Insert-Query} & \(O(\log^2 n)\) & \cref{prop:a2-iq} & \ok\\
A2 \op{Insert} & \(O(n)\) & \cref{thm:a2-insert} (linear rebuild in case 3) & \ok\textsuperscript{*}\\
A2 \op{Delete} & \(O(n\log n)\) & \cref{sec:a2-insert-on} and before & \ok\\
A2 \op{Delete-Query} & \(O(n)\) & \cref{tab:a2-summary} & \ok\\
A2 preprocessing & \(O(m+n\log n)\) & \cref{ch:preprocessing}: implemented \(\Theta(n^2+mn)\) & \notok\\
A1\(^+\)/A2\(^+\) \op{Delete-Query} & \(O(\sqrt m)\) & \cref{prop:yellin} & \ok\\
A1\(^+\)/A2\(^+\) updates & \(+O(m+n)\) & \cref{prop:yellin} & \ok\\
A1\(^+\)/A2\(^+\) preprocessing & \(O(mn+n^2)\) & \cref{prop:yellin} & \ok\\
Split graphs (\S8) & \(O(1)\) / \(O(m+n)\) & \cref{ch:split} & \ok\\
Connected components (\S4.1 remark) & \(O(1)\) query & not implemented (outside the four operations) & ---\\
\bottomrule
\end{tabularx}
\par\smallskip\raggedright\footnotesize\textsuperscript{*}\,Case 3 (double subsumption) relies on the
\(O(k)\) bound of the bottom-up construction of globally biased trees, which the code attributes
to BST85 Theorem~12; it could not be checked against BST85 itself (\cref{sec:bst85}). The
dynamic-tree bounds themselves were checked against ST83 (\cref{thm:st83}).
\end{table}

\section{Lessons learned}
\begin{itemize}
\item \textbf{Initialisation costs are real.} Two of the paper's bounds (A1 and A2
  preprocessing) silently depend on how memory is initialised. For A1 a different
  representation solved the problem; for A2 it cannot be solved with an explicit matrix.
\item \textbf{Worst-case dynamic trees are much harder than amortised ones.} An amortised
  splay-based link--cut tree is short; the worst-case Sleator--Tarjan structure with globally
  biased trees required about 1{,}250 lines (\file{worst_case_dynamic_tree} and
  \file{globally_biased_binary_tree}) and the largest share of the testing effort.
\item \textbf{A bound for one operation can hide a rare expensive case.} The double-subsumption
  case of A2's \op{Insert} looks harmless in the paper but needs \(\Theta(n)\) dynamic-tree
  operations if implemented literally.
\item \textbf{Oracles beat examples.} Exhaustive and randomised oracle tests reach update
  sequences that hand-made examples rarely cover (long chains of splits, absorptions and
  exchanges), which is why they form the backbone of the verification.
\end{itemize}

\section{Future work}
\begin{enumerate}
\item \textbf{A2 preprocessing.} Either find a representation that answers
  \(|K_x\cap K_y|\) for closest pairs fast enough without an explicit \(W\) (which, by
  \cref{thm:bmm}, would have consequences for matrix multiplication), or establish an
  unconditional lower bound in a suitable model. Both are open research questions.
\item \textbf{Connected components.} Implement the ET-tree extension sketched in Ibarra \S4.1
  (Henzinger and King~\cite{henzinger1999}) to answer connectivity queries in \(O(1)\).
\item \textbf{Ibarra's open problems.} The paper (\S9) asks for dynamic maintenance of a
  maximum independent set, a minimum colouring and a minimum clique cover of a chordal graph;
  the classical algorithms of Gavril~\cite{gavril1972} need a perfect elimination ordering,
  which is hard to maintain dynamically.
\item \textbf{Formal verification.} A machine-checked proof that the dynamic-tree code
  maintains the ST83/BST85 invariants would close the remaining gap between the tested code and
  the published bounds.
\item \textbf{Larger experiments.} Measure on large real-world chordal graphs (e.g.\ elimination
  graphs of sparse matrices) to locate the practical break-even point between A1 and A2.
\end{enumerate}

%% =====================================================================
\begin{thebibliography}{99}
\addcontentsline{toc}{chapter}{References}

\bibitem{ibarra2008} L.~Ibarra. Fully dynamic algorithms for chordal graphs and split graphs.
\emph{ACM Transactions on Algorithms}, 4(4), Article~40, 20 pages, 2008.
\href{https://doi.org/10.1145/1383369.1383371}{doi:10.1145/1383369.1383371}.

\bibitem{aho1974} A.~V.~Aho, J.~E.~Hopcroft and J.~D.~Ullman. \emph{The Design and Analysis of
Computer Algorithms}. Addison-Wesley, 1974.

\bibitem{alman2025} J.~Alman, R.~Duan, V.~Vassilevska Williams, Y.~Xu, Z.~Xu and R.~Zhou. More
asymmetry yields faster matrix multiplication. In \emph{Proc.\ ACM--SIAM Symposium on Discrete
Algorithms (SODA)}, 2025. arXiv:2404.16349 (\(\omega<2.371339\)).

\bibitem{openai2026} OpenAI. \emph{An Upper Bound of $9/4$ for the Matrix Multiplication Exponent}. Preprint, 2 October 2026. \url{https://github.com/openai/math/blob/main/preprints/Matrix-Multiplication-Nine-Fourths-October-2-2026/paper.pdf}. The result concerns exact arithmetic over $\mathbb{C}$.

\bibitem{shpilka2026} A.~Shpilka. \emph{Exposition of OpenAI's Matrix Multiplication in Exponent $9/4$ (and an explicit algorithm)}. Electronic Colloquium on Computational Complexity, TR26-281, 2026. \url{https://eccc.weizmann.ac.il/report/2026/281/}.

\bibitem{bent1985} S.~W.~Bent, D.~D.~Sleator and R.~E.~Tarjan. Biased search trees.
\emph{SIAM Journal on Computing}, 14(3):545--568, 1985.

\bibitem{bernstein1981} P.~A.~Bernstein and N.~Goodman. Power of natural semijoins.
\emph{SIAM Journal on Computing}, 10:751--771, 1981.

\bibitem{blair1993} J.~R.~S.~Blair and B.~Peyton. An introduction to chordal graphs and clique
trees. In \emph{Graph Theory and Sparse Matrix Computation}, IMA Volumes in Mathematics and its
Applications 56, pages 1--29. Springer, 1993. (Also Oak Ridge National Laboratory report
ORNL/TM-12203, 1992.)

\bibitem{briggs1993} P.~Briggs and L.~Torczon. An efficient representation for sparse sets.
\emph{ACM Letters on Programming Languages and Systems}, 2(1--4):59--69, 1993.

\bibitem{buneman1974} P.~Buneman. A characterization of rigid circuit graphs.
\emph{Discrete Mathematics}, 9:205--212, 1974.

\bibitem{foldes1977} S.~F\"oldes and P.~L.~Hammer. Split graphs. \emph{Congressus
Numerantium}, 19:311--315, 1977.

\bibitem{fulkerson1965} D.~R.~Fulkerson and O.~A.~Gross. Incidence matrices and interval
graphs. \emph{Pacific Journal of Mathematics}, 15(3):835--855, 1965.

\bibitem{gavril1972} F.~Gavril. Algorithms for minimum coloring, maximum clique, minimum
covering by cliques, and maximum independent set of a chordal graph. \emph{SIAM Journal on
Computing}, 1(2):180--187, 1972.

\bibitem{gavril1974} F.~Gavril. The intersection graphs of subtrees in trees are exactly the
chordal graphs. \emph{Journal of Combinatorial Theory, Series B}, 16:47--56, 1974.

\bibitem{golumbic1980} M.~C.~Golumbic. \emph{Algorithmic Graph Theory and Perfect Graphs}.
Academic Press, 1980.

\bibitem{grone1984} R.~Grone, C.~R.~Johnson, E.~M.~S\'a and H.~Wolkowicz. Positive definite
completions of partial Hermitian matrices. \emph{Linear Algebra and its Applications},
58:109--124, 1984.

\bibitem{hammer1981} P.~L.~Hammer and B.~Simeone. The splittance of a graph.
\emph{Combinatorica}, 1:275--284, 1981.

\bibitem{henzinger1999} M.~R.~Henzinger and V.~King. Randomized dynamic graph algorithms with
polylogarithmic time per operation. \emph{Journal of the ACM}, 46(4):502--516, 1999.

\bibitem{holee1989} C.~Ho and R.~C.~T.~Lee. Counting clique trees and computing perfect
elimination schemes in parallel. \emph{Information Processing Letters}, 31:61--68, 1989.

\bibitem{lundquist1990} M.~Lundquist. \emph{Zero patterns, chordal graphs, and matrix
completions}. PhD thesis, Clemson University, 1990.

\bibitem{rose1976} D.~J.~Rose, R.~E.~Tarjan and G.~S.~Lueker. Algorithmic aspects of vertex
elimination on graphs. \emph{SIAM Journal on Computing}, 5(2):266--283, 1976.

\bibitem{sleator1983} D.~D.~Sleator and R.~E.~Tarjan. A data structure for dynamic trees.
\emph{Journal of Computer and System Sciences}, 26(3):362--391, 1983.

\bibitem{tarjan1984} R.~E.~Tarjan and M.~Yannakakis. Simple linear-time algorithms to test
chordality of graphs, test acyclicity of hypergraphs, and selectively reduce acyclic
hypergraphs. \emph{SIAM Journal on Computing}, 13(3):566--579, 1984.

\bibitem{walter1978} J.~R.~Walter. Representations of chordal graphs as subtrees of a tree.
\emph{Journal of Graph Theory}, 2:265--267, 1978.

\bibitem{yellin1994} D.~M.~Yellin. An algorithm for dynamic subset and intersection testing.
\emph{Theoretical Computer Science}, 129:397--406, 1994.

\end{thebibliography}
%% ---- end expanded/14_concl.tex ----
%% ---- begin expanded/15_appendix.tex ----
%% =====================================================================
\appendix
\chapter{Build and Run Guide}\label{app:build}
%% =====================================================================
The implementation needs a C++17 compiler (tested with g++~13.3) and CMake~3.16 or newer; it has
no other dependencies. The Python audit scripts additionally need Python~3 and NetworkX.

\noindent\textbf{Obtain and build.}\par\nobreak
\begin{lstlisting}[style=shell]
sudo apt update && sudo apt install -y git g++ cmake
git clone https://github.com/khemsinghbhati/MTP-3rd-sem.git
cd MTP-3rd-sem
cmake -S . -B build -DCMAKE_BUILD_TYPE=Release -DMTP_BUILD_BENCHMARKS=ON
cmake --build build -j
\end{lstlisting}

\noindent\textbf{Run the tests.}\par\nobreak
\begin{lstlisting}[style=shell]
ctest --test-dir build --output-on-failure        # all 23 tests, about 45 s
ctest --test-dir build -R ibarra_a1               # only the A1 tests
\end{lstlisting}

\noindent\textbf{Run the programs.}\par\nobreak
\begin{lstlisting}[style=shell]
./build/mtp_cli          data/example_7_vertices.txt
./build/mtp_dynamic_cli  data/example_7_vertices.txt data/dynamic_example_ops.txt
./build/mtp_dynamic_cli  data/example_7_vertices.txt data/dynamic_example_ops.txt --fast-delete
./build/mtp_a2_cli       data/example_7_vertices.txt data/dynamic_example_ops.txt
./build/mtp_a2_cli       data/example_7_vertices.txt data/dynamic_example_ops.txt --fast-delete
./build/mtp_indexed_cli  data/example_7_vertices.txt data/dynamic_example_ops.txt
\end{lstlisting}

\noindent\textbf{Run the benchmarks and audits.}\par\nobreak
\begin{lstlisting}[style=shell]
./build/mtp_benchmarks              > a1_star.csv
./build/mtp_benchmark_a2_st         > a1_vs_a2_path.csv
./build/mtp_a2_preprocessing_audit  > a2_empty.csv
./build/mtp_a2_two_hub_stress
./build/mtp_a2_operation_audit
./build/mtp_benchmark_indexed
python3 -m pip install -r audit/requirements.txt
python3 audit/verify_chordal_identities.py 17
python3 audit/verify_w_updates.py
\end{lstlisting}

\paragraph{Use the library.} A minimal program that maintains a chordal graph with A1:
\begin{lstlisting}[style=cpp,numbers=none]
#include "mtp/sparse_ibarra_a1.hpp"
#include <iostream>
int main() {
    // the running example: 7 vertices, 10 edges
    auto g = mtp::Graph::from_edges(7, {{0,1},{1,3},{0,2},{2,3},{2,4},
                                       {2,5},{4,5},{3,6},{0,3},{1,2}});
    mtp::SparseIbarraChordalGraphA1 dyn(std::move(g));   // O(n+m)
    std::cout << dyn.insert_query(0, 4) << '\n';          // 1 (yes)
    dyn.insert_edge(0, 4);                                // O(n)
    std::cout << dyn.insert_query(4, 6) << '\n';          // 0 (no)
    std::cout << dyn.maximal_clique_count() << '\n';      // 4
}
\end{lstlisting}
Link it against the library \code{mtp_graphs} (\code{target_link_libraries(app mtp_graphs)}).

%% =====================================================================
\chapter{File-by-File Reference}\label{app:files}
%% =====================================================================
\begin{longtable}{@{}>{\raggedright\arraybackslash}p{0.42\textwidth}r>{\raggedright\arraybackslash}p{0.43\textwidth}@{}}
\caption[Files of the repository]{Source, test and benchmark files (line counts at commit
\code{651b236}).}\label{tab:files}\\
\toprule
\textbf{File} & \textbf{Lines} & \textbf{Contents}\\
\midrule
\endfirsthead
\toprule
\textbf{File} & \textbf{Lines} & \textbf{Contents}\\
\midrule
\endhead
\bottomrule
\endfoot
\multicolumn{3}{@{}l}{\emph{Core library (\code{include/mtp/*.hpp}, \code{src/*.cpp})}}\\
\file{graph.hpp/.cpp} & 40 / 95 & \code{Graph}; linear bulk loader (\cref{sec:graph-class})\\
\file{mcs.hpp/.cpp} & 20 / 96 & MCS, PEO test (\cref{sec:mcs-code,sec:peo-code})\\
\file{maximal_cliques.hpp/.cpp} & 15 / 63 & cliques from a PEO (\cref{sec:cliques-code})\\
\file{clique_tree_blair_peyton.hpp/.cpp} & 10 / 68 & Blair--Peyton clique tree (\cref{sec:bp})\\
\file{clique_graph.hpp/.cpp} & 34 / 51 & explicit \(W_G\) (test baseline)\\
\file{clique_tree.hpp/.cpp} & 35 / 139 & Kruskal clique tree + checker (test baseline)\\
\file{sparse_ibarra_a1.hpp/.cpp} & 78 / 334 & \textbf{A1} (\cref{ch:a1})\\
\file{dynamic_ibarra.hpp/.cpp} & 66 / 348 & dense A1 used by A2 (\cref{sec:dense-a1})\\
\file{globally_biased_binary_tree.hpp/.cpp} & 94 / 593 & BST85 trees (\cref{sec:bst85})\\
\file{worst_case_dynamic_tree.hpp/.cpp} & 67 / 494 & ST83 forest (\cref{sec:st83})\\
\file{dynamic_ibarra_a2.hpp/.cpp} & 57 / 293 & \textbf{A2} (\cref{ch:a2})\\
\file{indexed_dynamic_ibarra.hpp/.cpp} & 57 / 386 & indexed comparison variant (\cref{sec:indexed})\\
\file{fast_delete_ibarra.hpp/.cpp} & 37 / 114 & modified Yellin index (\cref{ch:fastdelete})\\
\file{dynamic_split_ibarra.hpp/.cpp} & 55 / 165 & split graphs (\cref{ch:split})\\
\midrule
\multicolumn{3}{@{}l}{\emph{Programs (\code{src/})}}\\
\file{main.cpp} & 77 & \code{mtp_cli}\\
\file{dynamic_main.cpp} & 90 & \code{mtp_dynamic_cli} (A1, A1\(^+\))\\
\file{a2_main.cpp} & 92 & \code{mtp_a2_cli} (A2, A2\(^+\))\\
\file{indexed_main.cpp} & 92 & \code{mtp_indexed_cli}\\
\midrule
\multicolumn{3}{@{}l}{\emph{Tests (\code{tests/})}}\\
\file{test_algorithms.cpp} & 212 & MCS, PEO, cliques\\
\file{test_clique_tree.cpp} & 192 & baseline clique tree, Prüfer oracle\\
\file{test_linear_clique_tree.cpp} & 197 & Blair--Peyton tree\\
\file{test_ibarra_dynamic.cpp} & 147 & dense A1\\
\file{test_sparse_ibarra_a1.cpp} & 147 & sparse A1\\
\file{test_sparse_ibarra_a1_fast.cpp} & 148 & sparse A1\(^+\)\\
\file{test_sparse_a1_stress.cpp} & 118 & sparse A1 adversarial\\
\file{test_ibarra_fast.cpp} & 148 & dense A1\(^+\)\\
\file{test_ibarra_fast_stress.cpp} & 53 & Yellin index stress\\
\file{test_ibarra_a2.cpp} & 247 & A2 versus A1\\
\file{test_a2_linear_insertion.cpp} & 65 & A2 double subsumption\\
\file{test_worst_case_dynamic_tree.cpp} & 155 & ST forest versus naive forest\\
\file{test_linear_st_forest.cpp} & 102 & linear ST construction\\
\file{test_bst85_linear_builder.cpp} & 67 & linear biased-tree build\\
\file{test_globally_biased_binary_tree.cpp} & 112 & join/split\\
\file{test_ibarra_indexed.cpp} & 158 & indexed variant\\
\file{test_ibarra_split.cpp} & 62 & split recognition\\
\file{test_ibarra_split_stats.cpp} & 73 & four split numbers\\
\file{test_deterministic_graph_lookup.cpp} & 43 & \code{Graph} membership\\
\midrule
\multicolumn{3}{@{}l}{\emph{Benchmarks (\code{benchmarks/})}}\\
\file{benchmark_dynamic.cpp} & 61 & \code{mtp_benchmarks}\\
\file{benchmark_a2_st.cpp} & 51 & \code{mtp_benchmark_a2_st}\\
\file{a2_preprocessing_audit.cpp} & 32 & \code{mtp_a2_preprocessing_audit}\\
\file{a2_two_hub_stress.cpp} & 84 & \code{mtp_a2_two_hub_stress}\\
\file{a2_operation_audit.cpp} & 81 & \code{mtp_a2_operation_audit}\\
\file{benchmark_indexed.cpp} & 47 & \code{mtp_benchmark_indexed}\\
\end{longtable}

%% =====================================================================
\chapter{Complete Program Traces}\label{app:traces}
%% =====================================================================
This appendix reproduces, without editing, the input files and the output of the programs for
every example used in the report. The commands were run (on 23 November 2026, commit
\code{651b236}) in a directory containing the listed files; the files
\file{example_7_vertices.txt}, \file{nonchordal_cycle4.txt} and \file{dynamic_example_ops.txt}
are in \file{data/} in the repository, the others are listed here in full. Note that A2 numbers
the cliques differently from A1, because it builds its clique tree by replaying insertions.

%% ---- begin parts/16_traces_generated.tex ----
\subsection*{Static program on the running example}
\begin{lstlisting}[style=shell,title={\small\texttt{example\_7\_vertices.txt}}]
7 10
0 1
1 3
0 2
2 3
2 4
2 5
4 5
3 6
0 3
1 2
\end{lstlisting}
\begin{lstlisting}[style=shell]
$ ./build/mtp_cli example_7_vertices.txt
MCS-derived ordering: 4 5 1 2 0 3 6
Maximal cliques (3):
{ 6 3 }
{ 3 0 2 1 }
{ 2 5 4 }
Clique tree T edges:
K1 -- K0 [weight=1]
K2 -- K1 [weight=1]
\end{lstlisting}
\subsection*{Static program on the 4-cycle (rejected, exit code 2)}
\begin{lstlisting}[style=shell,title={\small\texttt{nonchordal\_cycle4.txt}}]
4 4
0 1
1 2
2 3
3 0
\end{lstlisting}
\begin{lstlisting}[style=shell]
$ ./build/mtp_cli nonchordal_cycle4.txt
MCS-derived ordering: 2 1 0 3
The graph is not chordal: MCS ordering is not a PEO.
(exit code 2)
\end{lstlisting}
\subsection*{A1 on the running example with the repository's sample script}
\begin{lstlisting}[style=shell,title={\small\texttt{dynamic\_example\_ops.txt}}]
# Ibarra (2008), Sections 3-4.
# The queries leave the graph unchanged; accepted updates modify the clique tree.
Show
Insert-Query 0 4
Insert 0 4
Delete-Query 0 4
Delete 0 4
Insert-Query 4 6
Show
\end{lstlisting}
\begin{lstlisting}[style=shell]
$ ./build/mtp_dynamic_cli example_7_vertices.txt dynamic_example_ops.txt
Maximal cliques (3):
 K0 = { 6 3 }
 K1 = { 3 2 1 0 }
 K2 = { 2 5 4 }
 K1 -- K0 weight=1
 K2 -- K1 weight=1
Insert-Query 0 4 => yes
Insert 0 4 => yes
Delete-Query 0 4 => yes
Delete 0 4 => yes
Insert-Query 4 6 => no
Maximal cliques (3):
 K0 = { 6 3 }
 K1 = { 3 2 1 0 }
 K2 = { 2 5 4 }
 K1 -- K0 weight=1
 K2 -- K1 weight=1
\end{lstlisting}
\subsection*{A1 on the running example (\cref{fig:trace-ex})}
\begin{lstlisting}[style=shell,title={\small\texttt{ex\_ops.txt}}]
Show
Insert 0 4
Show
Insert-Query 4 6
Insert-Query 1 6
Delete-Query 0 1
Delete 0 1
Show
\end{lstlisting}
\begin{lstlisting}[style=shell]
$ ./build/mtp_dynamic_cli example_7_vertices.txt ex_ops.txt
Maximal cliques (3):
 K0 = { 6 3 }
 K1 = { 3 2 1 0 }
 K2 = { 2 5 4 }
 K1 -- K0 weight=1
 K2 -- K1 weight=1
Insert 0 4 => yes
Maximal cliques (4):
 K0 = { 6 3 }
 K1 = { 3 2 1 0 }
 K2 = { 2 5 4 }
 K3 = { 2 0 4 }
 K1 -- K0 weight=1
 K1 -- K3 weight=2
 K3 -- K2 weight=2
Insert-Query 4 6 => no
Insert-Query 1 6 => yes
Delete-Query 0 1 => yes
Delete 0 1 => yes
Maximal cliques (5):
 K0 = { 6 3 }
 K1 = { 3 2 0 }
 K2 = { 2 5 4 }
 K3 = { 2 0 4 }
 K4 = { 3 2 1 }
 K1 -- K3 weight=2
 K1 -- K0 weight=1
 K3 -- K2 weight=2
 K1 -- K4 weight=2
\end{lstlisting}
\subsection*{A1: deletion queries after an insertion (\cref{sec:delete-query-math})}
\begin{lstlisting}[style=shell,title={\small\texttt{ex\_ops2.txt}}]
Insert 0 4
Delete-Query 0 2
Delete-Query 2 4
Delete-Query 3 6
Delete 3 6
Show
\end{lstlisting}
\begin{lstlisting}[style=shell]
$ ./build/mtp_dynamic_cli example_7_vertices.txt ex_ops2.txt
Insert 0 4 => yes
Delete-Query 0 2 => no
Delete-Query 2 4 => no
Delete-Query 3 6 => yes
Delete 3 6 => yes
Maximal cliques (4):
 K0 = { 3 2 1 0 }
 K1 = { 2 5 4 }
 K2 = { 2 0 4 }
 K3 = { 6 }
 K0 -- K2 weight=2
 K2 -- K1 weight=2
 K0 -- K3 weight=0 (dummy)
\end{lstlisting}
\subsection*{A1 on the star \(K_{1,3}\): double subsumption (\cref{fig:trace-star})}
\begin{lstlisting}[style=shell,title={\small\texttt{star.txt}}]
4 3
0 1
0 2
0 3
\end{lstlisting}
\begin{lstlisting}[style=shell,title={\small\texttt{star\_ops.txt}}]
Show
Insert-Query 1 3
Insert 1 3
Show
Delete-Query 0 2
Delete-Query 0 1
Delete 1 3
Show
\end{lstlisting}
\begin{lstlisting}[style=shell]
$ ./build/mtp_dynamic_cli star.txt star_ops.txt
Maximal cliques (3):
 K0 = { 3 0 }
 K1 = { 0 2 }
 K2 = { 0 1 }
 K1 -- K0 weight=1
 K2 -- K0 weight=1
Insert-Query 1 3 => yes
Insert 1 3 => yes
Maximal cliques (2):
 K0 = { 0 2 }
 K1 = { 0 1 3 }
 K1 -- K0 weight=1
Delete-Query 0 2 => yes
Delete-Query 0 1 => yes
Delete 1 3 => yes
Maximal cliques (3):
 K0 = { 0 3 }
 K1 = { 0 2 }
 K2 = { 0 1 }
 K2 -- K1 weight=1
 K2 -- K0 weight=1
\end{lstlisting}
\subsection*{A1 on the star \(K_{1,3}\): insertion with an edge exchange}
\begin{lstlisting}[style=shell,title={\small\texttt{star\_ops2.txt}}]
Show
Insert-Query 1 2
Insert 1 2
Show
\end{lstlisting}
\begin{lstlisting}[style=shell]
$ ./build/mtp_dynamic_cli star.txt star_ops2.txt
Maximal cliques (3):
 K0 = { 3 0 }
 K1 = { 0 2 }
 K2 = { 0 1 }
 K1 -- K0 weight=1
 K2 -- K0 weight=1
Insert-Query 1 2 => yes
Insert 1 2 => yes
Maximal cliques (2):
 K0 = { 3 0 }
 K1 = { 0 1 2 }
 K1 -- K0 weight=1
\end{lstlisting}
\subsection*{A1 on the path \(P_5\)}
\begin{lstlisting}[style=shell,title={\small\texttt{path.txt}}]
5 4
0 1
1 2
2 3
3 4
\end{lstlisting}
\begin{lstlisting}[style=shell,title={\small\texttt{path\_ops.txt}}]
Show
Insert-Query 0 2
Insert-Query 0 3
Insert 0 2
Show
Insert-Query 0 3
\end{lstlisting}
\begin{lstlisting}[style=shell]
$ ./build/mtp_dynamic_cli path.txt path_ops.txt
Maximal cliques (4):
 K0 = { 4 3 }
 K1 = { 3 2 }
 K2 = { 2 1 }
 K3 = { 1 0 }
 K1 -- K0 weight=1
 K2 -- K1 weight=1
 K3 -- K2 weight=1
Insert-Query 0 2 => yes
Insert-Query 0 3 => no
Insert 0 2 => yes
Maximal cliques (3):
 K0 = { 4 3 }
 K1 = { 3 2 }
 K2 = { 1 0 2 }
 K1 -- K0 weight=1
 K2 -- K1 weight=1
Insert-Query 0 3 => yes
\end{lstlisting}
\subsection*{A1 on the fan with an isolated vertex (\cref{fig:dummy,fig:insert-query-no})}
\begin{lstlisting}[style=shell,title={\small\texttt{fan.txt}}]
6 7
0 1
0 2
1 2
1 3
2 3
2 4
3 4
\end{lstlisting}
\begin{lstlisting}[style=shell,title={\small\texttt{fan\_ops.txt}}]
Show
Insert-Query 0 4
Insert-Query 0 3
Insert 0 3
Show
\end{lstlisting}
\begin{lstlisting}[style=shell]
$ ./build/mtp_dynamic_cli fan.txt fan_ops.txt
Maximal cliques (4):
 K0 = { 5 }
 K1 = { 4 3 2 }
 K2 = { 2 3 1 }
 K3 = { 1 2 0 }
 K1 -- K0 weight=0 (dummy)
 K2 -- K1 weight=2
 K3 -- K2 weight=2
Insert-Query 0 4 => no
Insert-Query 0 3 => yes
Insert 0 3 => yes
Maximal cliques (3):
 K0 = { 5 }
 K1 = { 4 3 2 }
 K2 = { 2 1 0 3 }
 K1 -- K0 weight=0 (dummy)
 K2 -- K1 weight=2
\end{lstlisting}
\subsection*{A2 on the running example (same script as for A1)}
\begin{lstlisting}[style=shell]
$ ./build/mtp_a2_cli example_7_vertices.txt ex_ops.txt
Maximal cliques (3):
 K0 = { 2 4 5 }
 K1 = { 0 1 2 3 }
 K2 = { 3 6 }
 K0 -- K1 weight=1
 K1 -- K2 weight=1
Insert 0 4 => yes
Maximal cliques (4):
 K0 = { 2 4 5 }
 K1 = { 0 2 4 }
 K2 = { 0 1 2 3 }
 K3 = { 3 6 }
 K2 -- K1 weight=2
 K2 -- K3 weight=1
 K1 -- K0 weight=2
Insert-Query 4 6 => no
Insert-Query 1 6 => yes
Delete-Query 0 1 => yes
Delete 0 1 => yes
Maximal cliques (5):
 K0 = { 2 4 5 }
 K1 = { 0 2 4 }
 K2 = { 1 2 3 }
 K3 = { 0 2 3 }
 K4 = { 3 6 }
 K3 -- K2 weight=2
 K3 -- K4 weight=1
 K3 -- K1 weight=2
 K1 -- K0 weight=2
\end{lstlisting}
\subsection*{A2\(^+\) on the running example}
\begin{lstlisting}[style=shell]
$ ./build/mtp_a2_cli example_7_vertices.txt ex_ops.txt --fast-delete
Maximal cliques (3):
 K0 = { 2 4 5 }
 K1 = { 0 1 2 3 }
 K2 = { 3 6 }
 K0 -- K1 weight=1
 K1 -- K2 weight=1
Insert 0 4 => yes
Maximal cliques (4):
 K0 = { 2 4 5 }
 K1 = { 0 2 4 }
 K2 = { 0 1 2 3 }
 K3 = { 3 6 }
 K2 -- K1 weight=2
 K2 -- K3 weight=1
 K1 -- K0 weight=2
Insert-Query 4 6 => no
Insert-Query 1 6 => yes
Delete-Query 0 1 => yes
Delete 0 1 => yes
Maximal cliques (5):
 K0 = { 2 4 5 }
 K1 = { 0 2 4 }
 K2 = { 1 2 3 }
 K3 = { 0 2 3 }
 K4 = { 3 6 }
 K3 -- K2 weight=2
 K3 -- K4 weight=1
 K3 -- K1 weight=2
 K1 -- K0 weight=2
\end{lstlisting}
%% ---- end parts/16_traces_generated.tex ----

%% =====================================================================
\chapter{Glossary}\label{app:glossary}
%% =====================================================================
\begin{description}[style=nextline,leftmargin=2.2em]
\item[Absorption (subsumption)] An old maximal clique becomes non-maximal because it is
  contained in a new clique; its tree node is contracted into the new node.
\item[Chordal graph] A graph in which every cycle of length \(\ge4\) has a chord.
\item[Clique tree] A tree on the maximal cliques in which the cliques containing any vertex form
  a subtree; edge weights are intersection sizes.
\item[Closest pair] For a non-edge \(uv\), the nodes \(x\in T_u\), \(y\in T_v\) at minimum
  distance in the clique tree.
\item[Dashed / solid edge] In a Sleator--Tarjan tree, solid edges are stored inside path
  representations; dashed edges join a path to its parent node.
\item[Double subsumption] An insertion in which both \(K_x\) and \(K_y\) are absorbed.
\item[Dummy edge] A weight-0 clique-tree edge joining cliques of different components.
\item[Expose] The Sleator--Tarjan procedure that makes a node-to-root path solid.
\item[Globally biased tree] A binary search tree with weighted leaves whose depth is
  logarithmic in the inverse relative weight, with worst-case join and split (BST85).
\item[Heavy edge] An edge to a child \(v\) with \(2\size(v)>\size(\text{parent})\).
\item[Maximal clique] A clique not contained in a larger clique.
\item[MCS] Maximum cardinality search: repeatedly pick a vertex with the most numbered
  neighbours.
\item[Minimal separator] A minimal vertex set separating two vertices; in a chordal graph these
  are exactly the intersections along clique-tree edges.
\item[Oracle] An independent, simpler computation used to check answers in tests.
\item[Path set] The biased tree of the dashed child paths of a node, ordered by weight.
\item[PEO] Perfect elimination ordering: every vertex's later neighbours form a clique.
\item[Split graph] A graph whose vertices split into a clique and an independent set.
\item[Witness] For a vertex \(v\), a stored clique node containing \(v\) (used by A2).
\item[Worst-case bound] A bound that holds for every single execution of an operation.
\end{description}
%% ---- end expanded/15_appendix.tex ----
\end{document}
